% DER_XVIII-a: Grothendieck's text with the commentary set in, then the overview, inside a group that applies the chunk's own definitions
\dgrchapter{XVIII}{Chapter XVIII}{Catégories et ensembles accessibles}
\dgrpart{0}{XVIII-a}{Part 1 of 5, manuscript pages 50--138}{XVIII-a\enspace\textperiodcentered\enspace Chapter XVIII, part 1 of 5\enspace\textperiodcentered\enspace manuscript pages 50--138}{Chapter XVIII \textperiodcentered\ \textup{XVIII-a}}
\dgrchunkbegin
% DER_XVIII-a: the definitions of the chunk's preamble, applied inside its group
\let\printedpage\relax\newcommand{\printedpage}[1]{\if\relax\detokenize{#1}\relax\else\marginpar{\footnotesize p.~#1}\fi}
\let\msp\relax\newcommand{\msp}[1]{\marginpar{\footnotesize [#1]}}
\let\uHom\relax\newcommand{\uHom}{\underline{\mathrm{Hom}}}
\let\Ens\relax\newcommand{\Ens}{\mathrm{Ens}}
\let\Fais\relax\newcommand{\Fais}{\mathrm{Fais}}
\let\pstop\relax\newcommand{\pstop}{\mathrm{pstop}}
\let\cM\relax\newcommand{\cM}{\mathcal{M}}
\let\cN\relax\newcommand{\cN}{\mathcal{N}}
\let\cX\relax\newcommand{\cX}{\mathcal{X}}
\let\cY\relax\newcommand{\cY}{\mathcal{Y}}
\let\cF\relax\newcommand{\cF}{\mathcal{F}}
\let\Hom\relax\newcommand{\Hom}{\mathrm{Hom}}
\let\Top\relax\newcommand{\Top}{\mathrm{Top}}
\let\Ob\relax\newcommand{\Ob}{\mathrm{Ob}\,}
\let\defeq\relax\newcommand{\defeq}{\overset{\text{déf}}{=}}
\let\bt\relax\newcommand{\bt}{\boxtimes}
\let\tbt\relax\newcommand{\tbt}{\mathbin{\widetilde{\boxtimes}}}
\let\linj\relax\newcommand{\linj}{\lhook\joinrel\longrightarrow}
\let\pr\relax\newcommand{\pr}{\operatorname{pr}}
\let\card\relax\newcommand{\card}{\operatorname{card}}
\let\Fl\relax\newcommand{\Fl}{\operatorname{Fl}}
\let\Filt\relax\newcommand{\Filt}{\operatorname{Filt}}
\let\Ind\relax\newcommand{\Ind}{\operatorname{Ind}}
\let\hookxrightarrow\relax\newcommand{\hookxrightarrow}[2][]{\lhook\joinrel\xrightarrow[#1]{#2}}
\let\PF\relax\newcommand{\PF}{\mathrm{PF}}
\let\Aut\relax\newcommand{\Aut}{\operatorname{Aut}}
\let\uFilt\relax\newcommand{\uFilt}{\underline{\mathrm{Filt}}}
\let\circled\relax\newcommand{\circled}[1]{\tikz[baseline=(c.base)]\node[draw,circle,inner sep=1pt](c){#1};}

% ---- from der-ch18_p001-005.tex
\printedpage{} % no page number printed
% head: (none)

\begin{center}
{\Large\textsf{LES DÉRIVATEURS}}

\bigskip

\textbf{ALEXANDRE GROTHENDIECK}

\bigskip\bigskip

\textbf{Chapitre XVIII}

\medskip

\textbf{Catégories et ensembles accessibles}
\end{center}

\newpage

\printedpage{} % no page number printed
% head: (none)

Ce texte a été déchiffré et transcrit en \LaTeXe{} par M. Künzer. Il a été édité par M. Künzer, J. Malgoire, et G. Maltsiniotis. La transcription est aussi fidèle que possible au manuscrit. Seules quelques corrections mineures évidentes ont été effectuées. Pour les rares ajouts ou commentaires des éditeurs, ainsi que pour la numérotation originale des pages du manuscrit, les caractères de machine à écrire \texttt{[typewriter]} entre crochets sont utilisés. Quelques notes de bas de page ajoutées par les éditeurs sont signalées par «~N.~Éd.~». La numérotation des pages de ce chapitre commence à la page 50, car à l'origine ce chapitre était la suite du chapitre XIX qui se termine au milieu de la page 50.

Cette édition est provisoire. Les remarques, commentaires et corrections sont bienvenus. Envoyer un message à :

\textbf{georges.maltsiniotis@imj-prg.fr}

\hfill G. Maltsiniotis\hspace*{2cm}

\newpage

\printedpage{1}
% head: A. Grothendieck : Dérivateurs, Ch. XVIII, Édition des Universités de Montpellier II et Paris VII
% foot: Version 17 décembre 2025

\msp{page 50}

\section*{4\quad Paratopos et l'opération $\mathcal{M}\boxtimes\mathcal{N}$}

\subsection*{4.1\quad Pseudo-topos. Problème du produit.}

J'appelle \textit{catégorie pseudo-$\mathfrak{U}$-topos}, ou simplement \textit{pseudo-topos}, une $\mathfrak{U}$-catégorie $\mathcal{M}$ stable par $\mathfrak{U}$-limites inductives, et qui admet une \textit{petite} sous-catégorie $\mathcal{M}_0$ \textit{strictement génératrice} (i.e. génératrice par épimorphismes stricts). Alors $\mathcal{M}$ admet aussi des $\mathfrak{U}$-limites projectives (SGA 4 I 8.12.7). (Mais l'exemple de (Gr), cf. loc. cit. 8.12.9, montre que $\mathcal{M}^\circ$ n'est par forcément, pour autant, un pseudo-$\mathfrak{U}$-topos, faute à $\mathcal{M}$ d'admettre une petite sous-catégorie strictement \textit{co}génératrice. Mais le théorème 2.6 \texttt{[du chapitre XIX]} donne des conditions sur $\mathcal{M}$ (surtout du type exactitude) % [sic] par forcément -> pas forcément, misprint; also the parenthesis opened at "(Mais" is never closed
\msp{page 51}
pour que $\mathcal{M}$ admette assez d'injectifs -- ce qui impliquera alors l'existence d'une petite sous-catégorie strictement cogénératrice, donc que $\mathcal{M}^\circ$ est également un pseudo-topos (\footnote{C'est le cas notamment si $\mathcal{M}$ est un topos.}).

\dgrblock{dgrb957}{Why the theory of coefficients leads to cardinal arithmetic}
\begin{dcomment}{context}{p50}{Why the theory of coefficients leads to cardinal arithmetic}
Grothendieck is looking for categories of coefficients in which constructions on diagrams can be made uniformly. The first sentence already contains the two ingredients that will govern the chapter: arbitrary small colimits, which make extensions possible, and a small family of objects, which makes those extensions controllable. His geometrical language is deliberate. Calling a category a category of sheaves encourages us to regard a colimit-preserving functor as an inverse image, even when no topology and no left exactness have been supplied. The subsequent tensor problem will test how much of this geometrical picture survives.

There is unusually direct historical evidence for the place of this digression. In the opening paragraph of the letter of 2 April 1991 to Thomason, Grothendieck describes work on accessible categories as a return to SGA~4, I,~9, needed for his broader homotopical project; the editor identifies Chapter~XVIII explicitly \citep[opening paragraph, printed p.~2]{Grothendieck1991Letter}. Thus the relevance of the present pages does not have to be inferred from a modern theory of derivators. The connection to factorization methods is also internal to the manuscript: Chapter~XIX, Theorem~2.6, uses accessibility, cardinal filtrations, and exactness assumptions to construct the required lifting factorizations \citep[ms. pp.~25--26]{Grothendieck1990XIX}.

We keep the universe fixed. ``Small'' means small relative to $\mathfrak U$, and a $\mathfrak U$-category has $\mathfrak U$-small hom-sets. The whole category need not be small. A small full subcategory $S\subset\mathcal M$ is strictly generating in the sense used here when its canonical families of maps reconstruct objects and their morphisms; equivalently in this cocomplete setting, the restricted Yoneda functor
\[
 \rho_S:\mathcal M\longrightarrow [S^{\mathrm{op}},\Ens],
 \qquad x\longmapsto\mathcal M(-,x)|_S
\]
is fully faithful. Such a subcategory is called dense in current terminology. In particular,
\[
 x\simeq\mathop{\mathrm{colim}}_{(s\to x)\in S/x}s.
\]
The category $S/x$ is small: its objects form a small union of hom-sets and its arrows a small union of subsets of hom-sets of $S$. This observation is where the size of the generating family first does mathematical work.

A pseudo-topos category here is consequently a cocomplete category with a small dense subcategory. It is not assumed accessible in the modern sense, and the adjective does not assert the exactness axioms of a topos. Later, the paratopos condition adds a uniform presentability bound on a small generating family; the resulting categories are the locally presentable categories. We establish that identification in the commentary to Section~4.3, with the cardinal convention made explicit. It should not be read back into the first definition.

For the diagram-theoretic motivation, the commentary on Chapter~I, manuscript p.~4 (``From derived coefficients to diagram categories''), and on Chapter~XII, manuscript p.~225 (``Smoothness and properness depend on a coefficient theory''), provides the appropriate continuity. Here we develop the additional size arguments, rather than repeating the construction of a derivator from a relative category.
\end{dcomment}


Je visualise les catégories pseudo-topos comme étant la catégorie des “faisceaux” sur un objet $\mathcal{X}$, visualisé comme une sorte d'“espace” que je distingue de la catégorie $\mathcal{M}=\Fais(\mathcal{X})$ des faisceaux sur $\mathcal{X}$. C'est $\mathcal{X}$ qui sera pour ceci le “pseudo-topos”. Les pseudo-topos forment une 2-catégorie, la variance dans $\mathcal{X}$ étant opposée de celle dans $\Fais(\mathcal{X})=\mathcal{M}$. De façon précise, si $\mathcal{X}$ correspond à $\mathcal{M}$, $\mathcal{Y}$ à $\mathcal{N}$, un morphisme $f$ de $\mathcal{X}$ dans $\mathcal{Y}$ sera donné par un foncteur
\[
f^* \;:\; \mathcal{N} \longrightarrow \mathcal{M}
\tag*{(4.1.1)}
\]
en sens inverse, commutant aux $\varinjlim$ inductives quelconques (\footnote{Dire ici que les objets de $\mathcal{M}$ sont visualisés comme des “ouverts” généralisés de $\mathcal{X}$. On a $\mathcal{M}\xrightarrow{\sim}\uHom^!(\mathcal{M}^\circ,\Ens)$, i.e. $\mathcal{M}$ s'interprète bel et bien comme les “faisceaux” sur cette catégorie $\mathcal{M}$ d'ouverts \ldots}). (\textbf{N.B.} si $\mathcal{X}$, $\mathcal{Y}$ sont même des topos, i.e. $\mathcal{M}$, $\mathcal{N}$ des catégories-topos, et si on veut que $f$ soit un morphisme de topos, on exigera de plus que $f^*$ soit de plus \textit{exact à gauche}. On dira qu'un \textit{morphisme de pseudo-topos est exact} si de plus $f^*$ est exact à gauche.)

\msp{page 52}
En vertu du dual de loc. cit. 8.12.7, cela équivaut au fait que $f^*$ admette un adjoint à droite $f_*$,
\[
f_* \;:\; \mathcal{M} \longrightarrow \mathcal{N}\;,
\tag*{(4.1.2)}
\]
lequel commute donc aux petites $\varprojlim$\,. La donnée de $f_*$ satisfaisant à cette propriété détermine donc $f^*$ à isomorphisme unique près. Mais en général, il ne suffit pas de se donner un tel $f_*$, pour pouvoir affirmer qu'il corresponde à un $f^*$, i.e. qu'il admette un adjoint à droite\,; pour ceci, il faudrait qu'on sache que les foncteurs $\mathcal{M}\longrightarrow\Ens$ qui commutent aux petites $\varprojlim$ sont représentables, ce qui sera le cas si $\mathcal{M}$ admet une petite sous-catégorie \textit{co}génératrice, i.e. si $\mathcal{M}^\circ$ est également une catégorie pseudo-topos. % [sic] qu'il admette un adjoint à droite -> adjoint à gauche, f^* is the left adjoint of f_*

\newpage

\printedpage{2}
% head: A. Grothendieck : Dérivateurs, Ch. XVIII, Édition des Universités de Montpellier II et Paris VII
% foot: Version 17 décembre 2025

Il est sans doute judicieux de procéder comme dans le cas des topos, et d'appeler \textit{morphisme de pseudo-topos} $f:\mathcal{X}\longrightarrow\mathcal{Y}$ un quadruple
\[
(f^*,f_*,\alpha,\beta)
\tag*{(4.1.3)}
\]
\msp{page 53}
de deux foncteurs adjoints $(f^*,f_*)$, et des morphismes d'adjonction. Quand $f^*$ est une équivalence, ce qui équivaut à $f_*$ une équivalence, ou $\alpha$ et $\beta$ des isomorphismes, on dira qu'on a une \textit{équivalence de pseudo-topos}, et à toutes fins pratiques, on pourra alors regarder $\mathcal{X}$ et $\mathcal{Y}$ comme des objets essentiellement identiques.

\dgrblock{dgrb958}{Adjoints, representability, and the direction of a morphism}
\begin{dcomment}{exposition}{p52}{Adjoints, representability, and the direction of a morphism}
Grothendieck wants to recover the direct image from the inverse image, just as for a geometric morphism of topoi. The relevant existence theorem is the dual special adjoint functor theorem. In the present situation it says that a colimit-preserving functor out of a pseudo-topos category has a right adjoint. The representability statement used is SGA~4, I, Corollaries~8.12.7--8.12.8, applied to the opposite category \citep{SGA4I1972}.

Here is the application with its variance exposed. If $L:\mathcal N\to\mathcal M$ preserves small colimits and $m\in\mathcal M$, then
\[
 n\longmapsto\mathcal M(Ln,m):\mathcal N^{\mathrm{op}}\longrightarrow\Ens
\]
sends colimits in $\mathcal N$ to limits of sets. Representability provides an object $Rm$ and bijections
\[
 \mathcal M(Ln,m)\simeq\mathcal N(n,Rm).
\]
Yoneda makes these choices functorial in $m$ and gives $L\dashv R$. Applying the same argument to the diagonal $\mathcal M\to\mathcal M^I$, which preserves colimits, supplies all small limits in $\mathcal M$. No assertion that $\mathcal M^{\mathrm{op}}$ is another pseudo-topos is needed for either step.

The converse problem starts with $R$, not with $L$. Preservation of limits alone does not give a left adjoint on an arbitrary complete category. The small cogenerating hypothesis on $\mathcal M$ is a sufficient condition; it permits the undualized representability theorem to be applied to $m\mapsto\mathcal N(n,Rm)$. The printed phrase ``un adjoint \`a droite'' at the end of manuscript p.~52 has the wrong direction: the missing adjoint to $f_*$ is the left adjoint $f^*$. This is an edition slip, retained in the source and registered separately.

The asymmetry can already be seen in groups. There are arbitrarily large simple groups, for example the finitary alternating groups on sufficiently large infinite sets. Given a small family of groups $(G_i)$, choose a simple group $H$ whose cardinal exceeds every $|G_i|$. Every homomorphism $H\to G_i$ is trivial, since a nontrivial map from a simple group is injective. Such maps cannot separate the identity endomorphism of $H$ from the trivial endomorphism. Thus no small family cogenerates groups, although finite free groups give a small generating supply. This independently explains the obstruction invoked in SGA~4, I,~8.12.9; it is not an objection to the existence of right adjoints to colimit-preserving functors out of groups.

Finally, an adjunction includes its unit and counit. The equivalence between descriptions by left and right adjoints reverses natural transformations: a transformation $L\Rightarrow L'$ corresponds to a transformation $R'\Rightarrow R$. This is the reason for the opposite category in (4.1.5). Invertible unit and counit characterize an adjoint equivalence. Requiring $f^*$ also to preserve finite limits gives the usual extra condition for a geometric morphism when the coefficient categories are topoi; it is an additional requirement, not a consequence of having an adjoint.
\end{dcomment}


Pour deux pseudo-topos $\mathcal{X}$, $\mathcal{Y}$, j'ai envie d'introduire (comme dans le cas des topos) un pseudo-topos produit
\[
\mathcal{X}\times\mathcal{Y}\;,
\tag*{(4.1.4)}
\]
qui ait les propriétés d'un 2-produit dans la 2-catégorie des pseudo-topos (\footnote{Il apparaît finalement que ce qu'on construit n'est \textit{pas} un 2-produit dans la 2-catégorie envisagée, donc il vaut mieux le noter par un signe tel que $\mathcal{X}\boxtimes\mathcal{Y}$. Mais si $\mathcal{X}$, $\mathcal{Y}$ sont deux topos, $\mathcal{X}\boxtimes\mathcal{Y}$ est un 2-produit dans la 2-catégorie des topos.}). Je devrais préciser que par définition
\[
\begin{aligned}
\uHom_{\pstop}(\mathcal{X},\mathcal{Y}) \;&\simeq\; \uHom_!(\mathcal{F}(\mathcal{Y}),\mathcal{F}(\mathcal{X}))^\circ\\
&\simeq\; \uHom^{!'}(\mathcal{F}(\mathcal{X}),\mathcal{F}(\mathcal{Y})),
\end{aligned}
\tag*{(4.1.5)}
\]
où l'indice ! (resp. l'exposant !) indique la sous-catégorie pleine dans la catégorie $\uHom$ envisagée formée des foncteurs qui commutent aux petites $\varinjlim$ (resp. aux petites
\msp{page 54}
$\varprojlim$\,), et l'accent dans le $\uHom^{!'}$ dans le dernier membre indique qu'on exige de plus que le foncteur en question ait un adjoint à gauche.

Soient donc
\[
\mathcal{M} \;=\; \mathcal{F}(\mathcal{X})\;, \qquad \mathcal{N} \;=\; \mathcal{F}(\mathcal{Y})\;,
\tag*{(4.1.6)}
\]
et soit $\mathcal{Z}$ un troisième pseudo-topos, décrit par
\[
\mathcal{P} \;=\; \mathcal{F}(\mathcal{Z})\;.
\tag*{(4.1.7)}
\]
On veut expliciter, à équivalence près, la catégorie produit $\uHom_{\pstop}(\mathcal{Z},\mathcal{X})\times\uHom_{\pstop}(\mathcal{Z},\mathcal{Y})$, i.e.
\[
\uHom_!(\mathcal{M},\mathcal{P})^\circ \;\times\; \uHom_!(\mathcal{N},\mathcal{Q})^\circ.
\tag*{(4.1.8)}
\] % [sic] \mathcal{Q} -> \mathcal{P}, \mathcal{Q} is not defined; \mathcal{Z} is described by \mathcal{P} (4.1.7)
Je vais introduire pour cela une catégorie
\printedpage{3}
% head: A. Grothendieck : Dérivateurs, Ch. XVIII, Édition des Universités de Montpellier II et Paris VII
% foot: Version 17 décembre 2025
\msp{page 55}
\[
\mathcal{M}\boxtimes\mathcal{N}
\tag*{(4.1.9)}
\]
par la formule
\[
\mathcal{M}\boxtimes\mathcal{N} \;\simeq\; \uHom^{!!}(\mathcal{M}^\circ\times\mathcal{N}^\circ,\Ens)\;,
\tag*{(4.1.10)}
\]
où le signe $\uHom^{!!}$ désigne la sous-catégorie pleine de la catégorie $\uHom$, formée des bi-contrafoncteurs % doubt resolved: retain the diplomatic bi-contrafoncteurs reading; see corrections.csv, PDF leaf 5
\[
F \;:\; \mathcal{M}^\circ\times\mathcal{N}^\circ \longrightarrow \Ens\;,
\tag*{(4.1.10.1)}
\]
tels que pour tout $x$ dans $\mathcal{M}$, le foncteur
\[
y \longmapsto F(x,y) \;:\; \mathcal{N}^\circ \longrightarrow \Ens
\tag*{(4.1.10.2)}
\]
soit dans $\uHom^!(\mathcal{N}^\circ,\Ens)$, i.e. transforme petites $\varinjlim$ en $\varprojlim$\,, et que, symétriquement, pour tout $y$ dans $\mathcal{N}$, le foncteur
\[
x \longmapsto F(x,y) \;:\; \mathcal{M}^\circ \longrightarrow \Ens
\tag*{(4.1.10.3)}
\]
soit dans $\uHom^!(\mathcal{M}^\circ,\Ens)$. Mais cela signifie aussi que ces foncteurs sont \textit{représentables}\,; donc on trouve aussi
\[
\mathcal{M}\boxtimes\mathcal{N} \xleftarrow{\sim} \uHom^!(\mathcal{N}^\circ,\mathcal{M})
\tag*{(4.1.11)}
\]
et
\[
\mathcal{M}\boxtimes\mathcal{N} \xleftarrow{\sim} \uHom^!(\mathcal{M}^\circ,\mathcal{N})
\tag*{(4.1.12)}
\]
\msp{page 56}
en utilisant (dans chacune de ces formules) dans deux côtés les équivalences de catégories
\[
\left\{
\begin{aligned}
&\mathcal{M} &&\xrightarrow{\sim} &&\uHom^!(\mathcal{M}^\circ,\Ens)\\
&\mathcal{N} &&\xrightarrow{\sim} &&\uHom^!(\mathcal{N}^\circ,\Ens)\;.
\end{aligned}
\right.
\tag*{(4.1.13)}
\]
On peut dire aussi que la catégorie $\mathcal{M}\boxtimes\mathcal{N}$ est équivalente à celle des \textit{couples} de foncteurs
\[
\left\{
\begin{aligned}
&\varphi &&: &&\mathcal{M}^\circ &&\longrightarrow &&\mathcal{N}\\
&\psi &&: &&\mathcal{N}^\circ &&\longrightarrow &&\mathcal{M}\;,
\end{aligned}
\right.
\tag*{(4.1.14)}
\]
munis de deux morphismes d'adjonction
\[
\left\{
\begin{aligned}
&\psi\varphi^\circ &&\xleftarrow{\alpha} &&\mathrm{id}_{\mathcal{M}}\\
&\varphi\psi^\circ &&\xleftarrow{\beta} &&\mathrm{id}_{\mathcal{N}}
\end{aligned}
\right.
\tag*{(4.1.15)}
\]
satisfaisant les deux compatibilités habituelles exprimant l'adjonction de $\varphi^\circ$ et de $\psi$ (ou, de façon équivalente, de $\psi^\circ$ et de $\varphi$). Le foncteur
\[
(\varphi,\psi) \longmapsto \varphi \qquad (\text{resp.}\;\; (\varphi,\psi) \longmapsto \psi)
\]
donne une équivalence de la catégorie de ces couples avec $\uHom^!(\mathcal{M}^\circ,\mathcal{N})$ (resp. avec $\uHom^!(\mathcal{N}^\circ,\mathcal{M})$), de sorte qu'on a un carré commutatif d'équivalences canoniques de foncteurs :

\dgrblock{dgrb959}{The actual product and the proposed tensor}
\begin{dcomment}{exposition}{p54}{The actual product and the proposed tensor}
The attempted product is useful because it makes Grothendieck identify the universal property he really needs. Two inverse-image functors into a coefficient category and one functor preserving colimits separately in two variables carry different information. The former already have a simple representing category.

Let $\mathcal M,\mathcal N,\mathcal P$ be pseudo-topos categories, and denote by $\operatorname{Cocont}(\mathcal A,\mathcal B)$ the category of colimit-preserving functors and natural transformations. There is an equivalence
\[
 \operatorname{Cocont}(\mathcal M\times\mathcal N,\mathcal P)
 \simeq
 \operatorname{Cocont}(\mathcal M,\mathcal P)
 \times\operatorname{Cocont}(\mathcal N,\mathcal P).
\]
To prove it, write $0$ for initial objects and use the two inclusions $m\mapsto(m,0)$ and $n\mapsto(0,n)$. Every pair $(m,n)$ is canonically the coproduct of $(m,0)$ and $(0,n)$. A cocontinuous $F$ therefore satisfies
\[
 F(m,n)\simeq F(m,0)\amalg F(0,n).
\]
Conversely, two cocontinuous functors $G,H$ determine $F(m,n)=G(m)\amalg H(n)$; interchange of colimits shows that $F$ is cocontinuous. A natural transformation between two such functors is uniquely the coproduct of its restrictions. These constructions are mutually inverse up to the indicated canonical isomorphisms. Moreover, $\mathcal M\times\mathcal N$ is a pseudo-topos category: colimits are coordinatewise, and the union of the two small strictly generating families $(s,0)$ and $(0,t)$ is strictly generating.

After reversing the variance and the two-cells as in (4.1.5), this is precisely the bicategorical product of pseudo-topoi. The tensor category which the text is about to define solves a different universal problem, once its existence has been established:
\[
 \operatorname{Cocont}(\mathcal M\boxtimes\mathcal N,\mathcal P)
 \simeq \operatorname{Cocont}_{2}(\mathcal M\times\mathcal N,\mathcal P),
\]
where the subscript $2$ means preservation of colimits in each variable separately. We will use this formula only in the accessible setting or in the explicit examples where it is proved.

The difference is visible without large categories of coefficients. For sets, separate continuity of the contravariant bifunctor gives
\[
 F(X,Y)\simeq F(1,1)^{X\times Y},
\]
since $X$ and $Y$ are coproducts of singleton sets. Hence $\Ens\boxtimes\Ens\simeq\Ens$, whereas the category representing the first displayed universal property is $\Ens\times\Ens$. These categories are not equivalent: the terminal object of $\Ens$ has two subobjects, while the terminal object of $\Ens\times\Ens$ has four. The footnote on manuscript p.~53 correctly retracts the product interpretation. Within the smaller bicategory of topoi and geometric morphisms, imposing left exactness changes the problem and recovers the geometrical product; the two bicategories must not be interchanged.
\end{dcomment}


% ---- from der-ch18_p006-010.tex
\printedpage{4}
% head: A. Grothendieck : Dérivateurs, Ch. XVIII, Édition des Universités de Montpellier II et Paris VII

\msp{page 57}

\dgrblock{dgrb960}{Birepresentability and contravariant adjunctions}
\begin{dcomment}{exposition}{p56}{Birepresentability and contravariant adjunctions}
Grothendieck's definition packages a candidate coefficient object by all of its pairings with objects in the two factors. This is why a bifunctor, rather than a pair of functors, is the natural starting point. Fix $F:\mathcal M^{\mathrm{op}}\times\mathcal N^{\mathrm{op}}\to\Ens$ preserving limits separately. Representability in the first variable gives objects $\psi(y)\in\mathcal M$ such that
\[
 F(x,y)\simeq\mathcal M(x,\psi(y)).
\]
An arrow $y'\to y$ in $\mathcal N$ induces a transformation $F(-,y)\to F(-,y')$ and therefore an arrow $\psi(y)\to\psi(y')$. This determines a functor $\psi:\mathcal N^{\mathrm{op}}\to\mathcal M$. It preserves limits because each representable $\mathcal M(x,-)$ does, and the family of all representables detects a limiting cone. Representing in the other variable gives $\varphi:\mathcal M^{\mathrm{op}}\to\mathcal N$ and the natural bijection
\[
 \mathcal N(y,\varphi(x))\simeq F(x,y)
 \simeq\mathcal M(x,\psi(y)).
\]
Thus $\varphi^{\mathrm{op}}:\mathcal M\to\mathcal N^{\mathrm{op}}$ is left adjoint to $\psi$. The apparent two units in (4.1.15) are the unit and the opposite of the counit of this ordinary adjunction. Their triangular identities, not just the existence of two comparison arrows, are required.

The argument also works on morphisms. A transformation $F\to F'$ gives compatible transformations $\psi\to\psi'$ and $\varphi\to\varphi'$ by Yoneda. Conversely, either transformation determines the bifunctor transformation. This proves the equivalences (4.1.11)--(4.1.16), including their two-cell content. An object of the tensor is therefore a contravariant adjunction, not a cocontinuous functor from $\mathcal M$ to $\mathcal N$.

Small limits of such bifunctors are calculated pointwise. Indeed, if $(F_i)$ is a diagram and $(x_j)$ a colimit diagram in $\mathcal M$, then
\[
 \left(\lim_iF_i\right)(\mathrm{colim}_jx_j,y)
 \simeq\lim_i\lim_jF_i(x_j,y)
 \simeq\lim_j\lim_iF_i(x_j,y).
\]
The interchange is the universal property of limits in sets, with no filteredness assumption. The same argument treats the other variable. This establishes completeness. It does not construct coproducts or a left adjoint to the forthcoming inclusion in a presheaf category. The sentence on p.~57 announcing that the tensor is a pseudo-topos is a proposed conclusion; the questions immediately following it and the later footnote to Lemma~4.2.3 show that it is not yet proved.

There is also a precise size benefit to restricting both variables to small dense subcategories. Once Corollary~4.2.2 is established, transformations between two candidate tensors form a subset of a product indexed by $\operatorname{Ob}(S\times T)$ of small sets. This proves local smallness of the candidate tensor before cocompleteness is known. The restriction theorem addresses the size of hom-sets; accessibility will address the size of constructions.
\end{dcomment}


\[
\begin{tikzcd}[column sep=small]
 & \cM \tbt \cN \arrow[dl, "\sim" sloped] \arrow[dr, "\sim" sloped] & \\
\uHom^{!}(\cM^{\circ},\cN) \arrow[dr, "\sim" sloped] & & \uHom^{!}(\cN^{\circ},\cM) \arrow[dl, "\sim" sloped] \\
 & \underbrace{\cM \bt \cN}_{\defeq \uHom^{!!}(\cM^{\circ}\times\cN^{\circ},\Ens)} &
\end{tikzcd}
\tag*{(4.1.16)}
\]

\texttt{[où $\cM \tbt \cN$ est la catégorie des couples définie dans la page précédente.]}

Je dis que $\cM \bt \cN$ \textit{est une catégorie pseudo-topos}. Il est clair, \texttt{[en]} outre, qu'elle admet des petites $\varprojlim$ , qui se calculent argument par argument dans $\uHom(\cM^{\circ}\times\cN^{\circ},\Ens)$, ou au choix dans $\uHom^{!}(\cM^{\circ},\cN)$, ou $\uHom^{!}(\cN^{\circ},\cM)$. Comment construire des petites $\varinjlim$ ? Il n'est pas question d'invoquer un argument de cogénérateurs. Il faut d'ailleurs trouver aussi une petite sous-catégorie strictement génératrice. Visiblement, il y a quelque chose de non trivial à démontrer !

\section*{4.2\quad Cafouillages sur $\cM \bt \cN$}

Revenons à la situation générale d'un pseudo-topos $\cX$, décrit par la catégorie pseudo-topos $\cM$. Si $\cN$ est une catégorie quelconque, on définit la catégorie

\msp{page 58}
\[
\mathrm{Fais}(\cX,\cN) \qquad \text{ou} \qquad \cF(\cX,\cN)
\]
des faisceaux sur $\cX$, à valeurs dans $\cN$, par
\[
\cF(\cX,\cN) \ \defeq\ \uHom^{!}(\cM^{\circ},\cN) \ .
\tag*{(4.2.1)}
\]
On a donc un foncteur pleinement fidèle
\[
\cF(\cX,\cN) \ \defeq\ \uHom^{!}(\cM^{\circ},\cN) \ \linj\ \Hom^{!}(\cN^{\circ},\overline{\cM}) \ ,
\tag*{(4.2.2)}
\]
où $\overline{\cM}$ désigne la sous-catégorie de $\cM^{\wedge}$ (\textbf{N.B.} $\cM^{\wedge}$ n'est pas en général une $\mathcal{U}$-catégorie qu'à cela ne tienne) formée des foncteurs \textit{représentables} $\cM^{\circ}\longrightarrow\Ens$ (laquelle est une $\mathfrak{U}$-catégorie, équivalente à $\cM$ par l'inclusion canonique $\cM \linj \overline{\cM}$). Donc on a à isomorphisme unique près une factorisation de (4.2.2) en
\[
\cF(\cX,\cN) \ \overset{\text{pl. fid.}}{\lhook\joinrel\relbar\joinrel\relbar\joinrel\longrightarrow}\ \uHom^{!}(\cN^{\circ},\cM) \ ,
\tag*{(4.2.3)}
\]
l'image essentielle étant formée des foncteurs
\[
\psi \ :\ \cN^{\circ} \longrightarrow \cM
\]
tels que pour tout $x$ dans $\cM$,
\[
y \longmapsto \Hom(y,\psi(x)) \ :\ \cN^{\circ} \longrightarrow \Ens
\]
soit représentable, i.e. des foncteurs qui
% foot: Version 17 décembre 2025
\printedpage{5}
% head: A. Grothendieck : Dérivateurs, Ch. XVIII, Édition des Universités de Montpellier II et Paris VII
\msp{page 59}
admettent un \textit{adjoint à gauche}. Quand $\cN$ lui-même est une catégorie pseudo-topos, cette condition est automatique (elle résulte de l'hypothèse que $\psi$ commute aux petites $\varprojlim$ ). Donc dans ce cas, (4.2.3) est une équivalence de catégories.

Soit maintenant $S$ ($S$ comme “site”) une \textit{petite} sous-catégorie de $\cM$, génératrice par épimorphismes stricts. On sait que cela signifie que le foncteur composé
\[
\cM \xrightarrow{\ \sim\ } \uHom^{!}(\cM^{\circ},\Ens) \xrightarrow{\text{restriction à } S} S^{\wedge} \ (\ \defeq \uHom(S^{\circ},\Ens)\ )
\tag*{(4.2.4)}
\]
est pleinement fidèle, et ceci implique plus généralement, que pour toute catégorie $\cN$, le foncteur restriction
\[
\cF(\cX,\cN) \ \defeq\ \uHom^{!}(\cM^{\circ},\cN) \ \overset{\text{pl. fid.}}{\lhook\joinrel\relbar\joinrel\relbar\joinrel\longrightarrow}\ \uHom(S^{\circ},\cN)
\tag*{(4.2.5)}
\]
est pleinement fidèle. En effet, le deuxième membre, via le plongement pleinement fidèle
\[
\cN \ \overset{\text{pl. fid.}}{\lhook\joinrel\relbar\joinrel\relbar\joinrel\longrightarrow}\ \uHom^{!}(\cN^{\circ},\Ens) \ \linj\ \cN^{\wedge} \ ,
\]
se plonge dans
\[
\uHom(S^{\circ},\uHom^{!}(\cN^{\circ},\Ens)) \ \simeq\ \Hom^{!}(\cN^{\circ},S^{\wedge}) \ ,
\]
et la pleine fidélité de (4.2.5) équivaut donc à celle du foncteur composé
\[
\cF(\cX,\cN) \ \defeq\ \uHom^{!}(\cM^{\circ},\cN) \ \longrightarrow\ \uHom^{!}(\cN^{\circ},S^{\wedge}) \ \simeq\ \uHom(S^{\circ},\uHom^{!}(\cN^{\circ},\Ens)) \ .
\tag*{(4.2.6)}
\]

\msp{page 60}
Or on a commutativité dans le carré
\[
\begin{tikzcd}[column sep=large, row sep=large]
\cF(\cX,\cN) \defeq \uHom^{!}(\cM^{\circ},\cN) \arrow[r, "\text{\scriptsize(4.2.5)}"] \arrow[d, hook, "\text{\scriptsize(4.2.3)}"', "\text{\scriptsize pleinement fidèle}"] & \uHom(S^{\circ},\cN) \arrow[d, hook, "\text{\scriptsize pleinement fidèle}"] \\
\uHom^{!}(\cN^{\circ},\cM) \arrow[r, hook, "\text{\scriptsize(4.2.4)}", "\text{\scriptsize pleinement fidèle}"'] & \underbrace{\uHom^{!}(\cN^{\circ},S^{\wedge})}_{\simeq\, \uHom(S^{\circ},\uHom^{!}(\cN^{\circ},\Ens))}
\end{tikzcd}
\tag*{(4.2.7)}
\]
(où les flèches verticales sont des équivalences, si $\cN$ est lui-même une catégorie pseudo-topos). Notons d'autre part que l'on a une équivalence
\[
\uHom(S^{\circ},\cN) \ \xleftarrow{\ \sim\ }\ \uHom^{!}(S^{\wedge\circ},\cN) \ \defeq\ \cF(\Top(S),\cN) \ ,
\tag*{(4.2.8)}
\]
où $\Top(S)$ est le pseudo-topos (en fait, un topos) défini par la catégorie pseudo-topos $S^{\wedge}$. Pour ceci, notons qu'on a
\[
\uHom^{!}(S^{\wedge\circ},\cN) \ \xrightarrow{\ \sim\ }\ \uHom_{!}(S^{\wedge},\cN^{\circ})^{\circ} \ ,
\]
et d'autre part le foncteur restriction à $S \linj S^{\wedge}$ donne
\[
\uHom_{!}(S^{\wedge},\cN^{\circ}) \ \xrightarrow{\ \sim\ }\ \uHom(S,\cN^{\circ}) \ ,
\]
d'où
% foot: Version 17 décembre 2025
\printedpage{6}
% head: A. Grothendieck : Dérivateurs, Ch. XVIII, Édition des Universités de Montpellier II et Paris VII
\msp{page 61}
\[
\uHom(S^{\wedge\circ},\cN) \ \xrightarrow{\ \sim\ }\ (\uHom(S,\cN^{\circ}))^{\circ} \ \xrightarrow{\ \sim\ }\ \uHom(S^{\circ},\cN) \ ,
\]
ce qui n'est autre que (4.2.8). Ainsi :

\dgrblock{dgrb961}{Density reconstructs a continuous functor}
\begin{dcomment}{exposition}{p59}{Density reconstructs a continuous functor}
The move to a small site is the first effective reduction in the tensor problem. Grothendieck needs to remember a limit-preserving functor on a small collection of test objects, and density supplies exactly that reconstruction. It is helpful to prove full faithfulness directly, so that no completeness of an arbitrary target is silently inserted.

Let $S\subset\mathcal M$ be small and dense, and let $F,G:\mathcal M^{\mathrm{op}}\to\mathcal N$ preserve all small limits of $\mathcal M^{\mathrm{op}}$. This formulation includes the requirement that the relevant image cones have limits in $\mathcal N$. For each $x$, density gives $x=\mathrm{colim}_{S/x}s$, hence
\[
 F(x)\simeq\lim_{(S/x)^{\mathrm{op}}}F(s),\qquad
 G(x)\simeq\lim_{(S/x)^{\mathrm{op}}}G(s).
\]
Given a transformation $\eta:F|_{S^{\mathrm{op}}}\to G|_{S^{\mathrm{op}}}$, compose the projections from $F(x)$ with the maps $\eta_s$. They form a cone to the diagram defining $G(x)$, and therefore give a unique map $\eta_x:F(x)\to G(x)$. If $u:x\to x'$ is an arrow in $\mathcal M$, both composites from $F(x')$ to $G(x)$ have the same projections to every $G(s)$ indexed by $s\to x$; they are equal. Thus $\eta_x$ is natural in $x$. Uniqueness follows from the same limiting projections. This proves (4.2.5) for any $\mathcal N$ in which the functors under consideration exist.

The symmetry construction has an equally transparent formula. Given $F$, represent in $\mathcal M$ the functor
\[
 x\longmapsto\mathcal N(y,Fx),
\]
for each $y\in\mathcal N$. The resulting $\psi:\mathcal N^{\mathrm{op}}\to\mathcal M$ lies in the image exactly when, for every $x$, the presheaf
\[
 y\longmapsto\mathcal M(x,\psi(y))
\]
is representable in $\mathcal N$. This is the well-typed formula intended on manuscript p.~58. The printed $\Hom(y,\psi(x))$ interchanges the variables and cannot be evaluated with the stated types. In terms of adjunctions, $\psi$ must have a left adjoint $\mathcal M\to\mathcal N^{\mathrm{op}}$.

Now take both coefficient categories to be pseudo-topos categories. Restrict first along $S\subset\mathcal M$ and then along $T\subset\mathcal N$. The argument just given proves
\[
 \mathcal M\boxtimes\mathcal N\lhook\joinrel\longrightarrow
 [(S\times T)^{\mathrm{op}},\Ens]
\]
fully faithful. It also explains why the construction respects natural transformations, an issue that cannot be settled merely by asserting that every object is a colimit of generators. The functors being reconstructed must carry those presentations to the appropriate limits. This qualification will reappear in the proof of Theorem~4.4.16.
\end{dcomment}


\textbf{Proposition 4.2.1.} \textit{Soient $\cM=\cF(\cX)$ une catégorie pseudo-topos, munie d'une sous-catégorie pleine strictement génératrice $S$, $\cN$ une $\mathfrak{U}$-catégorie quelconque, alors le foncteur “restriction à $S$” suivant est pleinement fidèle.}
\[
\cF(\cX,\cN) \ \defeq\ \uHom^{!}(\cM^{\circ},\cN) \ \longrightarrow\ \uHom(S^{\circ},\cN) \ \xleftarrow{\ \sim\ }\ \underbrace{\uHom(S^{\wedge\circ},\cN)}_{\defeq\, \cF(\Top(S),\cN)} \ ,
\tag*{(4.2.1.1)}
\]
\textit{et il s'insère dans le carré commutatif (4.2.7) de foncteurs pleinement fidèles.}

\textbf{Corollaire 4.2.2.} \textit{Supposons que $\cN=\cF(\cY)$ soit également une catégorie pseudo-topos, munie d'une sous-catégorie pleine $T$ strictement génératrice. Alors}
\[
\cF(\cX,\cN) \ \simeq\ \cF(\cY,\cM) \ \simeq\ \underbrace{\uHom^{!!}(\cM^{\circ}\times\cN^{\circ},\Ens)}_{=\,\cM\bt\cN} \ ,
\]
\textit{et le foncteur restriction}
\[
\begin{aligned}
\cM\bt\cN \ \defeq\ \uHom^{!!}(\cM^{\circ}\times\cN^{\circ},\Ens) \quad &\longrightarrow \quad \uHom(S^{\circ}\times T^{\circ},\Ens)\\
&\simeq \quad (S\times T)^{\wedge}\\
&\defeq \quad \cF(\Top(S\times T))
\end{aligned}
\tag*{(4.2.2.1)}
\]
\textit{est pleinement fidèle.}

La première assertion a été vue dans 4.1. Pour la pleine fidélité de (4.2.2.1), on note que ce foncteur s'interprète, à équivalence de catégories près, comme le composé du foncteur pleinement fidèle (4.2.1.1), et du foncteur
\msp{page 62}
similaire relatif à $\cN$ (donc lui aussi pleinement fidèle)
\[
\cF(\cY,S^{\wedge}) \ \defeq\ \uHom^{!}(\cN^{\wedge},S^{\wedge}) \ \xrightarrow{\ \text{restr. à } T\ }\ \underbrace{\uHom(T^{\circ},S^{\wedge})}_{\uHom(S^{\circ}\times T^{\circ},\Ens)} \ ,
\tag*{(4.2.2.2)} % [sic] \uHom^{!}(\cN^{\wedge},S^{\wedge}) -> \uHom^{!}(\cN^{\circ},S^{\wedge}), cf. (4.2.2.3) where the same F(Y,S^) is Hom!(N°,S^)
\]
compte tenu de l'équivalence de symétrie
\[
\cF(\cY,S^{\wedge}) \ \defeq\ \uHom^{!}(\cN^{\circ},S^{\wedge}) \ \simeq\ \uHom^{!}(S^{\wedge\circ},\cN) \ .
\tag*{(4.2.2.3)}
\]

\dgrblock{dgrb962}{The missing completeness hypothesis in the presheaf equivalence}
\begin{dcomment}{examples}{p61}{The missing completeness hypothesis in the presheaf equivalence}
The restriction theorem is sound, but the displayed equivalence (4.2.8) asks for more: every functor $S^{\mathrm{op}}\to\mathcal N$ is supposed to extend continuously to $\widehat S^{\mathrm{op}}$. For an arbitrary $\mathcal N$ this is false. The free-cocompletion argument in the text is being applied with target $\mathcal N^{\mathrm{op}}$, which needs the colimits used by that argument.

Take $S$ to be the category with one object and only its identity, so that $\widehat S=\Ens$, and take $\mathcal N=\operatorname{FinSet}$. A functor $S^{\mathrm{op}}\to\mathcal N$ can select the two-element set $2$. If it extended to a continuous $F:\Ens^{\mathrm{op}}\to\operatorname{FinSet}$, the countable coproduct decomposition $\mathbb N=\coprod_{n\in\mathbb N}1$ would force
\[
 F(\mathbb N)\simeq\prod_{n\in\mathbb N}2.
\]
No finite set represents this product in $\operatorname{FinSet}$: testing the product universal property on $1$ would give a bijection from a finite set to the infinite set $2^{\mathbb N}$. Thus restriction is not essentially surjective. This counterexample concerns the claimed equivalence, not the already proved full faithfulness of restriction.

Here is the sufficient hypothesis and its proof. Suppose $\mathcal N$ has all small limits and let $A:S^{\mathrm{op}}\to\mathcal N$. For a presheaf $P$, define
\[
 \overline A(P)=\lim_{(S/P)^{\mathrm{op}}}A(s).
\]
The category of elements $S/P$ is small. For $n\in\mathcal N$, the universal property gives
\[
 \mathcal N(n,\overline A(P))
 \simeq\operatorname{Nat}\bigl(P,\mathcal N(n,A(-))\bigr).
\]
The right-hand side carries colimits in $P$ to limits of sets. Since representables detect limits, $\overline A$ carries colimits of presheaves to limits in $\mathcal N$. For $P=h_s$, Yoneda identifies $\overline A(h_s)$ with $A(s)$. The density argument of the preceding block gives uniqueness on objects and on transformations. Therefore
\[
 [S^{\mathrm{op}},\mathcal N]\simeq
 \operatorname{Cont}(\widehat S^{\mathrm{op}},\mathcal N)
\]
under the completeness hypothesis.

There are also two missing superscripts in the printed formulas on p.~61: the initial equivalence and the final term of (4.2.1.1) use unrestricted $\uHom(\widehat S^{\mathrm{op}},\mathcal N)$ where the continuous functors $\uHom^{!}$ are intended. Adding these superscripts repairs the notation but does not repair the missing hypothesis on $\mathcal N$. Corollary~4.2.2 takes $\mathcal N$ to be a pseudo-topos category, which is complete, so the corollary survives intact.
\end{dcomment}


Au total, quand $\cM$, $\cN$ sont tous deux des catégories pseudo-topos, munies de sous-catégories strictement génératrices $S$ resp. $T$, on a un plongement pleinement fidèle canonique
\[
\cM\bt\cN \ \overset{\rho}{\linj}\ (S\times T)^{\wedge} \ ,
\tag*{(4.2.9)}
\]
% foot: Version 17 décembre 2025
\printedpage{7}
% head: A. Grothendieck : Dérivateurs, Ch. XVIII, Édition des Universités de Montpellier II et Paris VII
qu'on interprète par la suite comme $(f_{S}^{\cM}\times f_{T}^{\cN})_{*}$, où
\[
\left\{
\begin{aligned}
f_{S}^{\cM} \ &:\ \cX \ \longrightarrow\ \Top S\\
f_{T}^{\cN} \ &:\ \cY \ \longrightarrow\ \Top T
\end{aligned}
\right.
\tag*{(4.2.10)}
\]
sont les morphismes de pseudo-topos, donnés par les plongements pleinement fidèles (foncteurs image directe)
\[
\left\{
\begin{aligned}
(f_{S}^{\cM})_{*} \ &:\ \cM \ \longrightarrow\ S^{\wedge}\\
(f_{T}^{\cN})_{*} \ &:\ \cN \ \longrightarrow\ T^{\wedge} \ ,
\end{aligned}
\right.
\tag*{(4.2.11)}
\]
lesquels admettent bel et bien des adjoints à gauche
\msp{page 63}
\[
\left\{
\begin{aligned}
(f_{S}^{\cM})^{*} \ &:\ S^{\wedge} \ \longrightarrow\ \cM\\
(f_{T}^{\cN})^{*} \ &:\ T^{\wedge} \ \longrightarrow\ \cN \ ,
\end{aligned}
\right.
\tag*{(4.2.12)}
\]
à savoir les prolongements canoniques (par commutativité aux petites $\varinjlim$ ) des foncteurs d'inclusion $S \linj \cM$, $T \linj \cN$.

La question essentielle, à présent, c'est visiblement de construire un adjoint à gauche pour le foncteur (4.2.9). Il revient au même, pour $(s,t)$ dans $S\times T$, de définir un objet
\[
s\bt t \ \in\ \Ob\cM\bt\cN \ ,
\tag*{(4.2.13)}
\]
tel que l'on ait, pour $F$ variable dans $\cM\bt\cN$, un isomorphisme fonctoriel en $F$,
\[
\underbrace{\rho(F)(s,t)}_{\defeq\, F(s,t)} \ \simeq\ \Hom(s\bt t,F) \ .
\]
Donc il faut prouver ceci :

\textbf{Lemme 4.2.3} (\footnote[4]{Pas prouvé ; peut-être faux !}). \textit{Soient $\cM$, $\cN$ deux catégories pseudo-topos, et $s\in\Ob\cM$, $t\in\Ob\cN$, et considérons sur $\cM\bt\cN\defeq\uHom^{!!}(\cM^{\circ}\times\cN^{\circ},\Ens)$ le foncteur $F\longmapsto F(s,t)$. Ce foncteur est} représentable \textit{(par un objet que je note $s\bt t$).}

L'idée naturelle serait de définir $s\bt t$ par
\[
(s\bt t)(x,y) \ =\ s(x)\times t(y) \ =\ \Hom_{\cM}(x,s)\times\Hom_{\cN}(y,t) \ ,
\]

\msp{page 64}
ce qui définit bien un foncteur
\[
\cM^{\circ}\times\cN^{\circ} \longrightarrow \Ens \ .
\]
Malheureusement, pour un des arguments fixé, disons pour $y$ fixé, ce foncteur n'est \textit{pas} compatible aux limites projectives, p.ex. il ne transforme pas l'objet final de $\cM^{\circ}$ correspondant à l'objet initial de $\cM$) en objet final de Ens, mais en $t(y)=\Hom_{\cN}(y,t)$. % [sic] "de $\cM$)" -> closing parenthesis without opening one, unbalanced as printed
% foot: Version 17 décembre 2025
\printedpage{8}
% head: A. Grothendieck : Dérivateurs, Ch. XVIII, Édition des Universités de Montpellier II et Paris VII
Néanmoins, je pense que $s\bt t$ doit bel et bien se définir comme produit de deux objets de $\cM\bt\cN$, dépendant fonctoriellement de $s$ resp. de $t$ (\footnote[5]{également faux, cf. plus bas.}). On est donc amené dès à présent, qu'on le veuille ou non, à essayer de définir des foncteurs canoniques, commutant aux petites $\varinjlim$ ,
\[
\cM \longrightarrow \cM\bt\cN \longleftarrow \cN
\tag*{(4.2.13)} % [sic] (4.2.13) -> (4.2.14), the label (4.2.13) is already used for s⊠t ∈ Ob M⊠N
\]
(lesquels seront les foncteurs image inverse pour les projections canoniques
\begin{tikzcd}[row sep=small, column sep=small, baseline=(\tikzcdmatrixname-2-1.base)]
 & \cX \\
\cX\times\cY \arrow[ur] \arrow[dr] & \\
 & \cY
\end{tikzcd}
) (\footnote[6]{Le cas où $\cM$, $\cN$ sont des catégories-topos, nous montre que ces foncteurs devraient être décrits comme $x\longmapsto x\bt e_{\cN}$, $y\longmapsto e_{\cM}\bt y$. Mais il n'est pas clair qu'on ait en dehors du cas des topos, et même en admettant l'existence du $x\bt y$, la formule (valable pour un topos) $x\bt y\simeq(x\bt e_{\cN})\times(e_{\cM}\bt y)$.}).

\dgrblock{dgrb963}{The reflector is the actual existence problem}
\begin{dcomment}{exposition}{p63}{The reflector is the actual existence problem}
At this point the tensor has a faithful realization as a category of presheaves. Grothendieck wants to build colimits in it by first building ordinary presheaf colimits and then reflecting them. The individual coefficient categories already illustrate this procedure without any accessibility assumption.

For the restricted nerve $\rho_S:\mathcal M\to\widehat S$, define
\[
 L_S(P)=\mathop{\mathrm{colim}}_{(s,a)\in S/P}s.
\]
For each $x\in\mathcal M$, a map from this colimit to $x$ is a compatible family of maps $s\to x$, one for each element $a\in P(s)$. Equivalently, it is a natural transformation $P\to\rho_Sx$. Thus $L_S\dashv\rho_S$. Density says that the counit $L_S\rho_Sx\to x$ is invertible. This constructs the morphisms (4.2.10)--(4.2.12) and shows exactly where cocompleteness of $\mathcal M$ enters.

For $\rho:\mathcal E=\mathcal M\boxtimes\mathcal N\hookrightarrow\widehat{S\times T}$ the same formula would require the colimits one is trying to prove exist. If a reflector $L$ exists, then
\[
 s\boxtimes t=L(h_{(s,t)}),\qquad
 \mathcal E(s\boxtimes t,F)\simeq F(s,t),
\]
and arbitrary colimits in $\mathcal E$ are $L$ applied to the corresponding colimits of nerves. This is a consequence of the adjunction, not a prior recipe for constructing $L$.

The restriction to generator pairs in Lemma~4.2.3 is therefore natural. Their evaluation functors are the coordinate functions of the embedding. If $\mathcal E$ were already cocomplete and these evaluations representable, their representing objects could be extended from $S\times T$ to all presheaves by the displayed colimit construction. At the present stage cocompleteness is not available. The wording ``Pas prouv\'e; peut-\^etre faux!'' should be retained as a genuine record of the state of the argument. The accessible case will have a positive answer; the general case needs a separate existence theorem and is not settled by the formal Yoneda calculations.
\end{dcomment}


\section*{4.3\quad L'opération $\cM\bt\cN$ et l'axiome d'accessibilité}

J'ai passé la journée hier, essayant d'établir l'existence des $x\bt y$ dans $\cM\bt\cN$, sans y parvenir – du moins pas sans faire une hypothèse supplémentaire d'accessibilité,
\msp{page 65}
dont il devient évident à présent qu'il faut la faire de toutes façons. Je vais quand même énoncer le positif de mes perplexités.

\dgrblock{dgrb964}{Why the naive formula fails, and the abelian calculation}
\begin{dcomment}{development}{p64}{Why the naive formula fails, and the abelian calculation}
Grothendieck has reached a useful obstruction: the external product of two representable presheaves need not be a separately continuous bifunctor. We can sharpen the obstruction to an exact criterion and then calculate an example in which the proposed product decomposition is impossible. These are elementary consequences of the definitions, developed here to locate the failed geometrical intuition. We checked the tensor constructions in \citet[Section~6.5]{Kelly1982} and the algebraic examples in \citet[p.~95]{Freyd1966}; the latter already contains the tensor calculus for module theories. We make no priority claim for the following formulations or proofs.

\textbf{The external-product criterion.}\ours{}\quad
Let $\mathcal M,\mathcal N$ be pseudo-topos categories. The bifunctor
\[
 H_{s,t}(x,y)=\mathcal M(x,s)\times\mathcal N(y,t)
\]
belongs to $\mathcal M\boxtimes\mathcal N$ if and only if $s$ and $t$ are terminal in their respective categories.

\emph{Proof.} If $H_{s,t}$ preserves the empty limit in the first variable, evaluate at the initial object $0_{\mathcal M}$. For every $y$ we obtain
\[
 1\simeq H_{s,t}(0_{\mathcal M},y)
 \simeq\mathcal N(y,t).
\]
This says exactly that $t$ is terminal. Interchanging the variables proves that $s$ is terminal. Conversely, if both are terminal, $H_{s,t}$ is the constant singleton bifunctor and preserves every limit separately. This proves both directions, including the degenerate case of a one-object coefficient category.\hfill$\square$

For the tensor itself, there is a particularly instructive classical calculation:
\[
 \operatorname{Ab}\boxtimes\operatorname{Ab}\simeq\operatorname{Ab},
 \qquad A\boxtimes B\longleftrightarrow A\otimes_{\mathbb Z}B.
\]
We give the full argument because it identifies the universal elements, rather than only the ambient category. Under (4.1.11), an object is a continuous $\psi:\operatorname{Ab}^{\mathrm{op}}\to\operatorname{Ab}$. Composing with the forgetful functor to sets gives a continuous contravariant set-valued functor, hence $U\psi(B)\simeq\operatorname{Ab}(B,C)$ for some $C$. The group structure on these sets is natural in $B$. By Yoneda, its multiplication is a morphism $\mu:C\times C\to C$ in $\operatorname{Ab}$, with an internal unit. That unit is necessarily zero; since $\mu$ is a group homomorphism and has zero as its unit,
\[
 \mu(c,d)=\mu(c,0)+\mu(0,d)=c+d.
\]
The inverse is consequently the usual negative. Thus the internal abelian group structure is forced, and
\[
 \psi(B)=\operatorname{Ab}(B,C)
\]
with its ordinary pointwise group structure. Natural transformations between two such functors are exactly homomorphisms between their representing groups; they automatically preserve these structures. Under the tensor description, the associated bifunctor is
\[
 F_C(A,B)=\operatorname{Ab}\bigl(A,\operatorname{Ab}(B,C)\bigr)
 \simeq\operatorname{Ab}(A\otimes_{\mathbb Z}B,C).
\]
The last bijection follows directly by sending a homomorphism $A\to\operatorname{Ab}(B,C)$ to its bilinear evaluation map $(a,b)\mapsto f(a)(b)$, and then using the defining generators and bilinear relations of $A\otimes_{\mathbb Z}B$. It proves that evaluation at $(A,B)$ is represented by $A\otimes_{\mathbb Z}B$.

Now the terminal object of $\operatorname{Ab}$ is $0$. The proposed formula involving terminal coefficients would give
\[
 (A\boxtimes0)\times(0\boxtimes B)=0,
\]
whereas $\mathbb Z\boxtimes\mathbb Z$ corresponds to $\mathbb Z$. This is a complete counterexample to the proposed decomposition outside topoi. The failure has no size-theoretic cause: abelian groups are locally finitely presentable and their tensor exists. What fails is the identification of a bilinear operation with a cartesian product operation.
\end{dcomment}


\dgrblock{dgrb965}{The modern tensor comparison and its precise scope}
\begin{dcomment}{later}{p64}{The modern tensor comparison and its precise scope}
The construction becomes a well-behaved tensor on locally presentable categories. Its universal property involves functors preserving colimits separately; its objects can be realized by the separately continuous bifunctors used in (4.1.10). Kelly's construction of tensors for small categories with prescribed colimits supplies an earlier categorical framework for this operation \citep[Section~6.5, equation~(6.24)]{Kelly1982}. Thus this part of the comparison predates the manuscript. The accessible presentation developed below connects the small presentation with the large category of coefficients.

The ordinary topos assertion also has a precise predecessor: \citet[Theorem~2.3, pp.~50--52]{Pitts1985} proves that the product of Grothendieck topoi is their tensor as cocomplete categories, in a more general indexed setting. The construction uses the topology on a product of sites generated by covers in each variable. Specializing the base to sets gives the assertion in Grothendieck's footnote on p.~53 and the external-product formula in the footnote on p.~64. Pitts's Definition~2.2 explicitly separates this theorem from an unrestricted existence claim for arbitrary locally small cocomplete categories.

The higher analogue is especially close in its formulas. For presentable $\infty$-categories, \citet[Proposition~4.8.1.15, Lemma~4.8.1.16, and Proposition~4.8.1.17]{Lurie2017HA} constructs the tensor and identifies it with the category of limit-preserving functors $\mathcal C^{\mathrm{op}}\to\mathcal D$. Examples~4.8.1.19--4.8.1.20 identify its restriction to $\infty$-topoi with their product and identify spaces as the tensor unit. These are structural analogues of the present formulas, with mapping spaces in place of hom-sets.

There is no implication here that replacing sets by spaces converts every assertion in these pages into a correct higher assertion. Ordinary filtered colimits, strict categories of elements, and cardinal estimates on actual hom-sets all have to be replaced by their appropriate higher versions. Conversely, the derivator formalism by itself does not establish an accessibility bound on an arbitrary coefficient category. The present chapter is supplying such bounds before the homotopical constructions are made. The unconditional pseudo-topos tensor problem and the positive accessible tensor theorem must therefore remain separate entries in the audit.
\end{dcomment}


\dgrblock{dgrb966}{The representability problem behind the tensor}
\begin{dcomment}{exposition}{p65:Je vais quand même énoncer le positif de mes perplexités.}{The representability problem behind the tensor}
Grothendieck has found the right category of bilinear coefficients; what is missing is the ability to make objects in it.  A separately continuous bifunctor is easy to recognize, and limits of such bifunctors remain separately continuous.  Colimits are different: their pointwise values generally destroy the continuity conditions.  The proposed object $x\boxtimes y$ would solve this problem at its smallest visible scale, by representing the operation which reads a coefficient at $(x,y)$.

Write $\mathcal E=\mathcal M\boxtimes\mathcal N$ and $P=(S\times T)^\wedge$, and let $\alpha:\mathcal E\hookrightarrow P$ be restriction.  If a reflection $L\dashv\alpha$ exists, then
\[
 \mathcal E\bigl(Lh_{(s,t)},F\bigr)
 \cong P(h_{(s,t)},\alpha F)
 \cong F(s,t).
\]
Thus $s\boxtimes t=Lh_{(s,t)}$.  The same adjunction supplies every small colimit:
\[
 \operatorname*{colim}_{i\in I}^{\mathcal E}F_i
 \cong L\!\left(\operatorname*{colim}_{i\in I}^{P}\alpha F_i\right).
\]
The images of representables generate $\mathcal E$ densely, because every presheaf is the colimit of its representables and $L$ preserves that colimit.  These are the sound implications $c\Rightarrow b'$ and $c\Rightarrow a$ in Proposition 4.3.1.  No left exactness of $L$ is needed.

There is a genuine difficulty in the reverse implication as written.  Lemma 4.3.2 assumes that $\mathcal E$ has small \emph{limits}, although its proof constructs a \emph{colimit} of representing objects.  The printed variance $\mathcal E^\circ\to\mathrm{Ens}$ is also reversed: $\xi\mapsto\alpha(\xi)(x)$ is covariant.  These readings occur in the edition, not only in the transcription.  The valid elementary lemma says that if $\mathcal E$ is cocomplete and each evaluation is corepresented by $r(x)$, then
\[
 LF=\operatorname*{colim}_{(x,u)\in X/F}r(x)
 \quad\hbox{satisfies}\quad
 \mathcal E(LF,e)\cong P(F,\alpha e).
\]
Indeed, applying $\mathcal E(-,e)$ turns this colimit into the limit of the evaluations $\alpha(e)(x)$ over the category of elements of $F$, which is precisely the set of natural transformations $F\to\alpha(e)$.  The representing objects form a functor $r:X\to\mathcal E$ by Yoneda, so the displayed colimit has the required coherent transition maps.

Completeness alone cannot replace cocompleteness in that elementary lemma.  Let $\mathcal E$ be the opposite of the ordered category of ordinals belonging to the fixed universe.  It has every small limit, since a small family of ordinals has a supremum, but it has no initial object.  Take $X=\{*\}$ and $\alpha(e)=\mathcal E(0,e)$.  Its only evaluation is corepresentable.  A left adjoint would nevertheless have to send the empty set to an initial object of $\mathcal E$, which does not exist.  This example concerns the lemma's generality; it is not a counterexample to the tensor construction, where $\alpha$ has additional properties.  The written argument for $b'\Rightarrow c$ therefore needs an existence theorem beyond this lemma.  The accessibility hypothesis introduced next provides one, as we prove below.
\end{dcomment}


\textbf{Proposition-perplexité 4.3.1.} \textit{Soient $\cM$, $\cN$ deux catégories pseudo-topos, d'où $\cM\bt\cN$ comme dans 4.1. Soient $S\subset\cM$, $T\subset\cN$ deux petites sous-catégories strictement génératrices, et considérons l'inclusion pleinement fidèle correspondante :}
\[
\begin{tikzcd}[column sep=large]
\underbrace{\cM\bt\cN}_{\defeq\, \uHom^{!!}(\cM^{\circ}\times\cN^{\circ},\Ens)} \arrow[r, hook, "\alpha"'] & S^{\wedge}\bt T^{\wedge} \ \simeq\ (S\times T)^{\wedge}. \arrow[l, dotted, bend right, "\varphi\,?"']
\end{tikzcd}
\tag*{(4.3.1.1)}
\]
\textit{Considérons les conditions suivantes :}
\begin{enumerate}[label={}]
\item[\textit{a)}] \textit{$\cM\bt\cN$ est un pseudo-topos.}
\item[\textit{b)}] \textit{Pour $x\in\Ob\cM$, $y\in\Ob\cN$, $x\bt y\in\Ob\cM\bt\cN$ existe, i.e. le foncteur (commutant aux petites $\varprojlim$ )}
\[
F \ \longmapsto\ F(x,y) \ , \qquad \cM\bt\cN \longrightarrow \Ens
\tag*{(4.3.1.2)}
\]
\textit{est représentable.}
\item[\textit{b$'$)}] \textit{Comme b), avec $x\in\Ob S$, $y\in\Ob T$.}
\item[\textit{c)}] \textit{Le foncteur $\alpha$ (lequel commute aux petites $\varprojlim$ ) admet un adjoint à gauche}
\[
\varphi \ :\ (S\times T)^{\wedge} \longrightarrow \cM\bt\cN \ .
\]
\end{enumerate}
\textit{On a alors les implications}
\[
b \ \Longrightarrow\ b' \ \Longleftrightarrow\ c \ \Longrightarrow\ a \ .
\tag*{(4.3.1.3)}
\]
% foot: Version 17 décembre 2025

% ---- from der-ch18_p011-015.tex
\setcounter{footnote}{6}

\printedpage{9}
% head: A. Grothendieck : Dérivateurs, Ch. XVIII, Édition des Universités de Montpellier II et Paris VII

\noindent\textit{De plus, si c est satisfait, on a ceci : pour $s\in S$, $t\in T$, considérons}
\[
(s,t)\ \in\ (S\times T)\ \subset\ (S\times T)^\wedge\ \simeq\ S^\wedge\boxtimes T^\wedge\ ,
\]
\textit{qui s'explicite aussi comme}
\[
\left\{
\begin{aligned}
&\rlap{$\displaystyle (s,t)\ =\ \pr_1^*(s)\times\pr_2^*(t)\ \simeq\ s\overset{S^\wedge,T^\wedge}{\boxtimes}t\in\Ob\,(S\times T)^\wedge\ ,$}\\
&\textit{i.e.}\ (s,t)(s',t')\ =\ s(s')\times t(t') &&=\ \Hom_S(s',s)\times\Hom_T(t',t)\\
& &&=\ \Hom_{S\times T}((s',t'),(s,t))\ ,
\end{aligned}
\right.
\tag*{(4.3.1.4)}
\]
\[
\begin{tikzcd}
S\times T \arrow[r, "\mathrm{pr}_2"] \arrow[d, "\mathrm{pr}_1"] & T\\
S\ ,
\end{tikzcd}
\]
\msp{page 66}\textit{on a}
% note: the label (4.3.1.4) is printed for this display as well as for the preceding one
\[
s\overset{\mathcal M,\mathcal N}{\boxtimes}t\ \simeq\ \varphi(s,t)\ ,
\tag*{(4.3.1.4)}
\]
\textit{et $\varphi(S\times T)=\{s\boxtimes t\mid s\in S,\ t\in T\}$ est une partie strictement génératrice de $\mathcal M\boxtimes\mathcal N$ (}\footnote{\textbf{N.B.} On a bien sûr $(s,t)=s\overset{S^\wedge,T^\wedge}{\boxtimes}t$.}\textit{).}

\noindent\textbf{N.B.} Je ne connais pas d'exemple où les conditions envisagées ne soient satisfaites, mais doute pourtant à présent que même la plus faible le soit toujours, i.e. que $\mathcal M\boxtimes\mathcal N$ soit un pseudo-topos dès que $\mathcal M$ et $\mathcal N$ le sont.

\noindent\textsc{Démonstration de} 4.3.1. L'implication b $\Longrightarrow$ b$'$ est tautologique. L'implication b$'$ $\Longleftrightarrow$ c résulte du

% marginal vertical bar printed beside Lemme 4.3.2 (from "Lemme 4.3.2." to "soit représentable.")
\noindent\textbf{Lemme 4.3.2.} \textit{Soient $\mathcal E$, $X$ des $\mathfrak U$-catégories avec $X$ petite, $\mathcal E$ stable par petites $\varprojlim$ ,}
\[
\alpha\ :\ \mathcal E\longrightarrow X^\wedge
\]
\textit{un foncteur. Pour que $\alpha$ admette un adjoint à gauche $\varphi$, il faut et il suffit que pour tout $x\in\Ob X$, le foncteur}
\[
\left\{
\begin{array}{ccl}
\xi & \longmapsto & \Hom_{X^\wedge}(x,\alpha(\xi))\ =\ \alpha(\xi)(x)\\
\mathcal E^\circ & \longrightarrow & \mathrm{Ens}
\end{array}
\right.
\]
\textit{soit représentable.}

\noindent C'est évidemment nécessaire. La suffisance résulte du fait que $X\subset X^\wedge$ est stricte%
% foot: Version 17 décembre 2025
\printedpage{10}%
% head: A. Grothendieck : Dérivateurs, Ch. XVIII, Édition des Universités de Montpellier II et Paris VII
\msp{page 67}%
ment génératrice dans $X^\wedge$, i.e. que tout objet $F$ de $X^\wedge$ est limite inductive $F=\varinjlim\limits_I F_i$, les $F_i$ dans $X$ (on prend $I=X/F$), d'où
\[
\begin{aligned}
\Hom(F,\alpha(\xi))\ &\simeq\ \varprojlim_I\ \underbrace{\Hom(F_i,\alpha(\xi))}_{\Hom(\rho(F_i),\xi)}\\
&\simeq\ \Hom(\varinjlim_I\,(\rho(F_i),\xi))\ ,
\end{aligned}
\]
donc le foncteur $\xi\longmapsto\Hom(F,\alpha(\xi))$ est représenté par $\rho(F)\overset{\text{déf}}{=}\varinjlim\limits_I\rho(F_i)$.

\noindent Prouvons c $\Longrightarrow$ a. Pour l'existence des petites $\varinjlim$ dans $\mathcal M\boxtimes\mathcal N$, et l'existence d'une partie strictement génératrice, on invoque le lemme général :

% marginal vertical bar printed beside Lemme 4.3.3 (from "Lemme 4.3.3." to "comme")
\noindent\textbf{Lemme 4.3.3.} \textit{Soit $\alpha:\mathcal E\hookrightarrow\mathcal F$ un foncteur pleinement fidèle admettant un adjoint à gauche $\rho$.}
\begin{enumerate}[label={}]
\item[\textcircled{a}] \textit{Si les $\varinjlim$ d'un type donné $I$ existent dans $\mathcal F$, elles existent dans $\mathcal E$, et se calculent comme}
\[
\varinjlim_I{}_{\mathcal E}\,\xi_i\ =\ \rho(\varinjlim_I{}_{\mathcal F}\,\alpha(\xi_i))\ .
\tag*{(4.3.3.1)}
\]
\item[\textcircled{b}] \textit{Si $U$ est une petite partie strictement génératrice de $\mathcal F$, $\rho(U)$ est une petite partie strictement génératrice dans $\mathcal F$.}
\end{enumerate}

\noindent En effet, si $\xi$ désigne le deuxième membre \texttt{[de l'égalité 4.3.3.1]}, et $\eta$ un objet variable de $\mathcal E$, on aura
\[
\begin{aligned}
\Hom_{\mathcal E}(\xi,\eta)\ &\overset{\text{adj.}}{\simeq}\ \Hom_{\mathcal F}(\varinjlim_I\alpha(\xi_i),\alpha(\eta))\\
&\simeq\ \varprojlim_I\quad\underbrace{\Hom_{\mathcal F}(\alpha(\xi_i),\alpha(\eta))}_{\simeq\,\Hom_{\mathcal E}(\xi_i,\eta)\ \ \text{car $\alpha$ pleinement fidèle}}\quad .
\end{aligned}
\]
\msp{page 68}
\[
\simeq\ \ \varprojlim_I\Hom_{\mathcal E}(\xi_i,\eta)\ ,
\]
ce qui prouve que $\xi$ est une limite inductive des $\xi_i$ dans $\mathcal E$ (\footnote{Pour prouver b), on note que la pleine fidélité de $\alpha$ implique que $\rho$ est essentiellement surjectif ; comme $\rho$ commute aux petites $\varinjlim$ et que tout objet de $\mathcal F$ est de la forme $\varprojlim\limits_I x_i$, les $x_i$ dans $U$, on aura $\rho(x)=\varinjlim\limits_I\rho(x_i)$, les $\rho(x_i)$ dans $\rho(U)$, on gagne.}).

\noindent Ceci prouve que (moyennant c)) $\mathcal M\boxtimes\mathcal N$ est stable par petites $\varinjlim$ , donc il reste à prouver que $\mathcal M\boxtimes\mathcal N$ admet une petite famille strictement génératrice. Cela résultera du b) dans le lemme, puisque $S\times T$ engendre strictement $(S\times T)^\wedge$. Cela prouve en même temps la dernière assertion de la proposition, sous réserve de prouver la formule (4.1.3.4). Il faut donc prouver un isomorphisme fonctoriel en $F\in\Ob\mathcal M\boxtimes\mathcal N$, % [sic] (4.1.3.4) -> (4.3.1.4), the formula s ⊠ t ≃ φ(s,t) to be proved is labelled (4.3.1.4) on p. 9
\[
\Hom(\rho(s,t),F)\ \simeq\ F(s,t)\ ,
\]
mais par adjonction on a
\[
\Hom(\rho(s,t),F)\ \simeq\ \Hom((s,t),\alpha(F))\ =\ \alpha(F)(s,t)\ =\ F(s,t)\ ,
\]
OK.

% foot: Version 17 décembre 2025
\printedpage{11}
% head: A. Grothendieck : Dérivateurs, Ch. XVIII, Édition des Universités de Montpellier II et Paris VII

\noindent\msp{page 69}\textbf{Proposition 4.3.4.} \textit{Soit $\mathcal M$ une $\mathfrak U$-catégorie. Les conditions suivantes sont équivalentes.}
\begin{enumerate}[label={}]
\item[\textit{a)}] \textit{$\mathcal M$ est une catégorie $\mathfrak U$-pseudo-topos, i.e. elle est stable par $\mathfrak U$-limites inductives, et admet une $\mathfrak U$-petite sous-catégorie $S$ strictement génératrice.}
\item[\textit{b)}] \textit{Il existe une $\mathfrak U$-petite catégorie $S$, et un foncteur $i:\mathcal M\hookrightarrow S^\wedge$} pleinement fidèle, \textit{ayant un adjoint à gauche.}
\end{enumerate}

\noindent\textsc{Démonstration}. On a a) $\Longrightarrow$ b), car il suffit de prendre pour $S$ dans b) celui spécifié dans a), l'hypothèse sur $S$ implique que le foncteur canonique
\[
\mathcal M\ \longrightarrow\ S^\wedge\ ,\qquad\quad \text{composé de } \mathcal M\hookrightarrow\ \mathcal M^\wedge\ \xrightarrow{\ \text{restriction}\ }S^\wedge
\]
est pleinement fidèle. Qu'il admette un adjoint à gauche résulte du lemme 4.3.2. L'implication b) $\Longrightarrow$ a) résulte du lemme 4.3.3.

\dgrblock{dgrb967}{A reflective presentation and its generators}
\begin{dcomment}{exposition}{p69:Qu'il admette un adjoint}{A reflective presentation and its generators}
Grothendieck now moves from the particular tensor problem to the presentation of its possible factors.  The useful feature of a pseudo-topos is that its objects can be reconstructed from a fixed small category of probes.  This is the precise content of the restricted Yoneda embedding; it is stronger than asking the probes merely to distinguish parallel arrows.

For a small full strictly generating subcategory $S\subset\mathcal M$, put
\[
 N(x)(s)=\mathcal M(s,x),\qquad
 L(F)=\operatorname*{colim}_{(s,u)\in S/F}s.
\]
The colimit exists because $\mathcal M$ is cocomplete.  The universal property of a colimit and Yoneda give $\mathcal M(LF,x)\cong\operatorname{Nat}(F,Nx)$.  Strict generation says that the canonical diagram $S/x\to\mathcal M$ has colimit $x$; equivalently, the counit $LNx\to x$ is an isomorphism.  Hence $N$ is fully faithful.  This proves Proposition 4.3.4 directly, without the incorrect completeness hypothesis of Lemma 4.3.2.  The density statement is the one formulated in \citet[Expos\'e I, Proposition 7.2(i)]{SGA4I1972}.

Conversely, suppose $L:P(S)\rightleftarrows\mathcal M:N$ is a reflection.  If $D:I\to\mathcal M$ is a small diagram, then $L\operatorname{colim}_I ND$ represents its cocones, by adjunction and full faithfulness.  Thus $\mathcal M$ is cocomplete.  For $x\in\mathcal M$, apply $L$ to the category-of-elements presentation of $Nx$:
\[
 x\cong LNx
 \cong\operatorname*{colim}_{(s,u)\in S/Nx}Lh_s.
\]
This presentation is natural in $x$ and exhibits the family $Lh_s$ as dense.  The target of this family is $\mathcal M$; correspondingly, the last $\mathcal F$ in the printed Lemma 4.3.3(b) must be read mathematically as $\mathcal E$.  The source wording is retained in the text and recorded in the corrections register.

The adjunction also explains the accessibility assertion in Corollary 4.3.5.  Suppose all $Lh_s$ preserve sufficiently filtered colimits under their Hom functors.  Since $S$ is small, choose one regular cardinal $\kappa$ which works for every $s$.  Colimits in a presheaf category are pointwise, so
\[
 N\operatorname{colim}_I x_i\cong\operatorname{colim}_I Nx_i
 \quad\text{for every $\kappa$-filtered $I$}
\]
is equivalent to the same assertion after evaluating at each $s$, hence to $\kappa$-presentability of all $Lh_s$.  Every $x$ is a small colimit of these objects.  Choose a larger regular cardinal $\lambda$ than $\kappa$ and than the number of arrows in one such presentation of $x$.  In sets, $\lambda$-filtered colimits commute with limits of fewer than $\lambda$ arrows.  Applying this to the Hom formula for that presentation proves that $x$ is $\lambda$-presentable.  Thus each object has a rank, although the ranks of all objects need not share one bound.
\end{dcomment}


\noindent\textbf{Corollaire 4.3.5.} \textit{Soit $\mathcal M$ une $\mathfrak U$-catégorie. Pour que $\mathcal M$ soit un $\mathfrak U$-pseudo-topos} accessible \textit{(i.e. dont tous les objets sont accessibles), il faut et il suffit qu'il existe $S$ dans} Cat \textit{et un foncteur pleinement fidèle et} accessible $\mathcal M\overset{\alpha}{\hookrightarrow}S^\wedge$\textit{, admettant un adjoint à gauche $\rho$.}

\noindent\msp{page 70}\textsc{Démonstration}. Il suffit de noter que si $S$, $\alpha$ sont comme dans 4.3.4 b), alors $\alpha$ est accessible si et seulement si les $\rho(S)$ sont formées d'objets accessibles dans $\mathcal M$ ; comme cette partie est strictement génératrice, on sait que cela implique que \textit{tout} objet de $\mathcal M$ est accessible. (Cf. SGA 4 I 9.9.)

\dgrblock{dgrb968}{From accessible probes to an ind-presentation}
\begin{dcomment}{exposition}{p70:Il suffit de noter que si}{From accessible probes to an ind-presentation}
Theorem 4.3.6 seeks an intrinsic description of the categories for which the tensor construction should work.  In contemporary terminology its main conditions describe a locally presentable category.  The indexing convention needs care: throughout these pages, ``grand devant $\pi$'' permits simultaneous solutions to diagrams with at most $\pi$ arrows.  With $\pi$ infinite, this is modern $\kappa$-filteredness for $\kappa=\pi^+$.  An object called $\pi$-accessible here is therefore $\pi^+$-presentable, not generally $\pi$-presentable.  The implication from a cocomplete category with a small strong generator of presentable objects to local presentability is also stated in \citet[Theorem 1.20]{AdamekRosicky1994}.  For the dense generator used here, a direct proof clarifies the role of the ranks.

Assume $\mathcal M$ is cocomplete and has a small full dense subcategory $S$ of accessible objects.  Choose an infinite regular $\kappa$ such that every object of $S$ is $\kappa$-presentable.  Enlarge $S$ to a small full subcategory $C$ closed, up to isomorphism, under colimits of diagrams with fewer than $\kappa$ arrows.  Here is a construction which also proves the required smallness.  Start with $S$; at each successor stage choose colimits of every such diagram in the preceding stage, and take the full subcategory on the resulting set of objects.  At a limit stage take the union of the preceding object sets.  Carry out $\kappa$ stages.  All sets used remain in the universe.  By regularity, a diagram with fewer than $\kappa$ arrows in the union appears at some earlier stage, so its chosen colimit also belongs to the union.

All objects of $C$ remain $\kappa$-presentable.  For a diagram $D:K\to\mathcal M$ of that size, the formula
\[
 \mathcal M(\operatorname{colim}_K D,-)
 \cong\operatorname{lim}_{K^{\mathrm{op}}}\mathcal M(D(-),-)
\]
and the corresponding interchange theorem in sets prove closure of presentable objects under these colimits.  At every stage this applies to the objects just constructed.

For any $x\in\mathcal M$, the category $C/x$ is $\kappa$-filtered.  Given a diagram of fewer than $\kappa$ arrows in $C/x$, take the colimit of its underlying diagram in $C$.  Its structure maps to $x$ induce a map from this colimit to $x$, giving a cocone in $C/x$.  This includes the empty diagram, because $C$ contains an initial object.  Moreover
\[
 x\cong\operatorname*{colim}_{(c,u)\in C/x}c.
\]
To see density after enlarging $S$, a cocone from $C/x$ to $y$ restricts to one from $S/x$, hence to a unique map $x\to y$.  For each $c\to x$, its agreement with the given map $c\to y$ can be tested after every map $s\to c$, and density of $S$ gives that agreement.  Thus the restriction has neither lost nor added cocones.

The functor which takes a $\kappa$-filtered diagram of objects of $C$ to its colimit in $\mathcal M$ is now essentially surjective.  It is fully faithful because, for diagrams $(c_i)_{i\in I}$ and $(d_j)_{j\in J}$, presentability gives
\[
 \mathcal M(\operatorname{colim}_I c_i,\operatorname{colim}_J d_j)
 \cong\operatorname*{lim}_{i\in I^{\mathrm{op}}}
          \operatorname*{colim}_{j\in J}\mathcal M(c_i,d_j),
\]
the defining Hom formula for ind-objects.  Consequently $\mathcal M\simeq\operatorname{Ind}_\kappa C$, with the subscript here in the modern convention.  When $\kappa=\pi^+$ this is Grothendieck's $\operatorname{Ind}_\pi C$.

This argument proves the structural part of the theorem without asserting that the particular slice-cardinality filtration chosen on page 72 already satisfies all the axioms of a cardinal filtration.  That stronger numerical statement is exactly where the subsequent attempted proof encounters difficulty.  The converse structural implication uses the cocompleteness of $\operatorname{Ind}_\kappa C$ when $C$ has the indicated small colimits; Section 4.4 supplies it.  No inaccessible cardinal beyond the standing universe is required for this passage.
\end{dcomment}


\dgrblock{dgrb969}{Why accessibility supplies the tensor reflector}
\begin{dcomment}{exposition}{p70:Il suffit de noter que si}{Why accessibility supplies the tensor reflector}
The accessible presentation gives a concrete answer to the existence problem which opened Section 4.3.  It replaces a class of continuity requirements by a set of requirements on a small category.  We give the construction, because that change of size is the mathematical reason the added axiom works.

Take a common regular cardinal $\kappa$ and small full subcategories $C\subset\mathcal M$, $D\subset\mathcal N$ of $\kappa$-presentable objects, closed under colimits of fewer than $\kappa$ arrows, with $\mathcal M\simeq\operatorname{Ind}_\kappa C$ and $\mathcal N\simeq\operatorname{Ind}_\kappa D$.  Write $P=(C\times D)^\wedge$.  In $P$ consider the set $\Sigma$ of maps
\[
 \operatorname*{colim}_{i\in K}^{P}h_{(c_i,d)}
       \longrightarrow h_{(\operatorname{colim}_K^C c_i,d)}
\]
for all diagrams in $C$ of fewer than $\kappa$ arrows and all $d\in D$, together with the analogous maps in the second variable.  There is a set of such maps, since the categories are small and one may choose a set of representatives of all indexing shapes of the stated size.  A presheaf $F$ is orthogonal to $\Sigma$ exactly when it takes those colimits to limits separately: orthogonality means that precomposition with each displayed map induces a bijection into $F$.

The orthogonal subcategory $\Sigma^\perp\subset P$ is reflective.  To see the size control explicitly, choose a regular $\mu$ such that all domains and codomains in $\Sigma$ are $\mu$-presentable presheaves.  This is possible by writing each as a small colimit of representables and taking one bound for the set $\Sigma$.  Starting with $X_0\in P$, perform the following two operations at every successor stage.  First push out the coproduct of all maps $s:A\to B$ in $\Sigma$ along all maps $A\to X_\alpha$, thus adjoining every required extension.  Then coequalize every pair of maps $B\rightrightarrows Y$ into that new object $Y$ which agree on $A$, for every $s\in\Sigma$.  At limit stages take the colimit.  Continue for $\mu$ stages.

For an orthogonal object $Z$, each operation induces a bijection from maps out of the new stage into $Z$ to maps out of the preceding stage into $Z$.  The pushout operation does so because extensions into $Z$ exist uniquely; the coequalizer operation does so because those extensions agree whenever they agree on $A$.  The same holds at limits.  Thus $X_0\to X_\mu$ is already universal for maps into orthogonal objects.  The endpoint is itself orthogonal: a map $A\to X_\mu$ factors through one stage by presentability and is extended at the next; two maps $B\rightrightarrows X_\mu$ agreeing on $A$ factor through a common stage, their equality on $A$ holds at a later stage, and the coequalizer step then makes them equal.  The ordinal $\mu$ is $\mu$-filtered, which justifies each factorization.  Hence this construction is the reflector $L:P\to\Sigma^\perp$.

To identify this category with Grothendieck's tensor, first recall a useful description of the restricted ind-completion.  A presheaf $H$ on a small $\kappa$-cocomplete category $C$ belongs to $\operatorname{Ind}_\kappa C$ if and only if it carries the indicated colimits in $C$ to limits of sets.  In one direction, representables have that property, and $\kappa$-filtered colimits commute with those limits.  In the other, a diagram of fewer than $\kappa$ arrows in $C/H$ has a cocone: take its underlying colimit in $C$ and use the compatible elements to obtain the element of $H$ at the colimit.  The canonical category-of-elements presentation then proves the assertion.

Each factor is complete, and its reflective presheaf embedding computes limits.  Indeed, objects in a full reflective subcategory are precisely those orthogonal to every reflection unit.  A limit of such objects remains orthogonal to every unit, since Hom into a limit is a limit; orthogonality to its own unit then makes that unit invertible.  Thus the presheaf limit belongs to the reflective subcategory.

It follows that a separately continuous $F\in\Sigma^\perp$ gives a functor $D^{\mathrm{op}}\to\mathcal M$ preserving limits of fewer than $\kappa$ arrows.  The dual functor $D\to\mathcal M^{\mathrm{op}}$ therefore preserves the corresponding colimits.  By the ind-completion universal property proved in Section 4.4, it extends uniquely to a colimit-preserving functor $\mathcal N\to\mathcal M^{\mathrm{op}}$.  Dualizing and applying the Hom pairing recovers a separately continuous bifunctor on $\mathcal M^{\mathrm{op}}\times\mathcal N^{\mathrm{op}}$.  Conversely, restriction of such a bifunctor lies in $\Sigma^\perp$, and its values on general objects are determined by their ind-presentations.  These constructions are inverse, including natural transformations.

Thus $\mathcal M\boxtimes\mathcal N\simeq\Sigma^\perp$ is a reflective subcategory of $P$.  It is cocomplete, with colimits obtained by reflection, and has the small dense family $Lh_{(c,d)}$.  Its inclusion preserves $\mu$-filtered colimits, since orthogonality to maps with $\mu$-presentable endpoints is preserved by those colimits.  Accordingly the reflected representables are $\mu$-presentable, and the tensor is again locally presentable.  Finally every evaluation $F\mapsto F(x,y)$ has a representing object: it preserves all limits and is accessible, or, equivalently, it is obtained from the small evaluation objects using the cocontinuity in $x$ and $y$ and the now-established colimits.  For the latter explicit construction, choose presentations $x=\operatorname{colim}_{C/x}c$, $y=\operatorname{colim}_{D/y}d$ and take
\[
 x\boxtimes y
   =\operatorname*{colim}_{(c\to x,d\to y)\in(C/x)\times(D/y)}Lh_{(c,d)}.
\]
Applying $\operatorname{Hom}(-,F)$ verifies the representing property directly.  This proves all the existence assertions of Proposition 4.3.1 under the accessibility hypothesis, without using the disputed general implication $b'\Rightarrow c$.
\end{dcomment}


\noindent Je vais essayer de donner un énoncé plus complet :

\noindent\textbf{Théorème 4.3.6.} \textit{Soit $\mathcal M$ une $\mathfrak U$-catégorie. Conditions équivalentes :}
\begin{enumerate}[label={}]
\item[\textit{a)}] \textit{$\mathcal M$ stable par petites $\varinjlim$ , et elle admet une petite sous-catégorie pleine $S$, strictement génératrice, et formée d'objets accessibles de $\mathcal M$.}
\item[\textit{a$'$)}] \textit{$\mathcal M$ est une catégorie $\mathfrak U$-pseudo-topos faiblement accessible, i.e. satisfait a), avec en plus} tout \textit{élément de $\mathcal M$ (pas seulement ceux de $S$) accessible dans $\mathcal M$.}
\item[\textit{b)}] \textit{Il existe une petite catégorie $S$, et un foncteur $\alpha:\mathcal M\hookrightarrow S^\wedge$} pleinement fidèle et accessible, \textit{tel que $\alpha$ admette un adjoint à gauche.}
\item[\textit{c)}] \msp{page 71}\textit{$\mathcal M$ est stable par petites $\varinjlim$ , est accessible, et admet une filtration cardinale} $(\Filt^\pi\mathcal M)_{\pi\geqslant\pi_0}$.
\item[\textit{d)}] \textit{Il existe un cardinal $\pi_1$, une petite catégorie $\mathcal M_{\pi_1}$, stable par limites inductives de cardinal $\leqslant\pi_1$, tels qu'on ait une équivalence de catégories}
\[
\mathcal M\ \xleftarrow{\ \sim\ }\ \Ind_{\pi_1}\mathcal M_{\pi_1}\ ,
\tag*{(4.3.6.1)}
\]
\textit{où le deuxième membre désigne la sous-catégorie pleine de} $\Ind(\mathcal M_{\pi_1})$ \textit{formée des ind-objets $(X_i)_{i\in I}$ de $\mathcal M_{\pi_1}$ tels que $I$ soit grand devant $\pi_1$.}
\end{enumerate}

\noindent\textsc{Démonstration de} 4.3.6 (\footnote{Pour une démonstration complète (celle donnée ici canule à partir d'un certain moment), cf. §4.6 (p. 139-148).}). On a déjà vu a$'$ $\Longrightarrow$ a $\Longrightarrow$ b $\Longrightarrow$ a$'$, i.e. l'équivalence de a, a$'$, b, soit PTA cette condition. Nous allons prouver les implications
\[
\mathrm{PTA}\ \Longrightarrow\ \mathrm{c}\ \Longrightarrow\ \mathrm{d}\ \Longrightarrow\ \mathrm{PTA}\ ,
\]
% foot: Version 17 décembre 2025
\printedpage{12}%
% head: A. Grothendieck : Dérivateurs, Ch. XVIII, Édition des Universités de Montpellier II et Paris VII
ce qui établira le théorème.

\noindent\textit{Prouvons PTA $\Longrightarrow$ c.} Soit $S$ une petite sous-catégorie pleine strictement génératrice. Ainsi, \msp{page 72}pour tout $x$ dans $\mathcal M$ on a
\[
x\ =\ \varinjlim_{s\in S/x}s
\tag*{(4.3.6.2)}
\]
(cf. SGA 4 I 7.2 (i)). Soit $\pi_0=\card\Fl(S)$, et pour tout cardinal $\pi\geqslant\pi_0$, soit
\[
\begin{array}{lcl}
\Filt^\pi\mathcal M & = & \text{sous-catégorie pleine de $\mathcal M$ formée des objets $x$ tels que}\\
 & & \card\,(\Fl(S/x))\leqslant\pi\ .
\end{array}
\tag*{(4.3.6.3)}
\]
Je veux prouver que l'on trouve ainsi une filtration cardinale sur $\mathcal M$, i.e. que les conditions a) b) c) de SGA 4 I 9.12 sont satisfaites. C'est prouvé dans loc. cit. 9.13 sous des hypothèses légèrement différentes, la condition d'accessibilité de $\mathcal M$ étant remplacée par celle (beaucoup moins sympathique!) que les familles de flèches épimorphiques strictes dans $\mathcal M$ sont épimorphiques strictes universelles. Cette condition n'est utilisée que dans la démonstration de la condition b), à savoir la stabilité de $\Filt^\pi\mathcal M$ par $\varinjlim\limits_I$ de cardinal $\leqslant\pi$ (sans avoir d'ailleurs à supposer $I$ filtrante, ni a fortiori grand devant $\pi_0$). Il faut d'ailleurs dire que la définition de $\Filt^\pi\mathcal M$ donnée dans loc. cit. n'est pas tout à fait la même que (4.3.6.3), \msp{page 73}mais c'est
% the printed line break "car-/dinal" inside the display is kept as printed
\[
\begin{array}{lcl}
\Filt'^{\pi}\mathcal M & = & \text{sous-catégorie pleine de $\mathcal M$ formée des objets $x$ tels qu'il}\\
 & & \text{existe une famille épimorphique stricte $x_i\xrightarrow{\alpha_i}x$ de car-}\\
 & & \text{dinal $\leqslant\pi$, avec $x_i\in S\ \forall\ i\in I$.}
\end{array}
\tag*{(4.3.6.3$'$)}
\]
Ainsi on a
\[
\Filt^\pi\mathcal M\ \subset\ \Filt'^{\pi}\mathcal M\ ,
\]
mais dans loc. cit. on prouve que pour $x$ dans $\Filt'^{\pi}\mathcal M$, on a
\[
\card\Fl(S/x)\ \leqslant\ \pi^{\pi_0}\ ,
\]
d'où
\[
\Filt^\pi\mathcal M\ \subset\ \Filt'^{\pi}\mathcal M\ \subset\ \Filt^{\pi^{\pi_0}}\mathcal M\ ,
\]
donc si $\pi=\pi^{\pi_0}$, p.ex. $\pi=2^c$ avec $c\geqslant\pi_0$, alors on trouve
\[
\Filt^\pi\mathcal M\ =\ \Filt'^{\pi}\mathcal M\ .
\]

\noindent Donc les deux filtrations on l'air $\pm$ équivalentes. Mais la démonstration donnée de la stabilité de $\Filt'^{\pi}$ semble vraiment faire appel de façon essentielle à la condition d'universalité des familles épimorphiques strictes (donc épimorphiques effectives, avec les hypothèses faites impliquant l'existence de petites $\varprojlim$ ). Donc il me faut bon gré, mal gré, refaire le travail! % [sic] on l'air -> ont l'air, verb avoir (3rd person plural) intended

% foot: Version 17 décembre 2025
\printedpage{13}
% head: A. Grothendieck : Dérivateurs, Ch. XVIII, Édition des Universités de Montpellier II et Paris VII

\noindent\msp{page 74}Condition a) : $\Filt^\pi$ équivalente à une petite catégorie. C'est immédiat, puisque par (4.3.6.2), $\Filt^\pi$ est formée de $\varprojlim\limits_I$ dans $\mathcal M$ d'objets de $S$ (petite catégorie), avec $I$ telle que $\card\Fl I\leqslant\pi$.

\noindent Condition b) : $\Filt^\pi$ stable par $\varprojlim\limits_I$ , quand $\card\Fl I\leqslant\pi$ (\footnote{quand $I$ filtrante et grande devant $\pi_0$ -- je ne le prouve que dans ce cas, hélas!}). Soit donc $(x_i)_{i\in I}$, où
\[
\begin{array}{ccc}
I & \longrightarrow & \mathcal M\\[1ex]
i & \longmapsto & x_i
\end{array}
\]
un système inductif dans $\mathcal M$, avec des $x_i$ dans $\Filt^\pi$. Considérons le produit fibré
\[
\begin{tikzcd}
I \arrow[d, "\varphi"'] & S/I = J \arrow[l, "\text{cofibré}"'] \arrow[d] \arrow[dr, dotted, "\psi"] & & \\
\mathcal M & S/\mathcal M \arrow[l, "\text{but}"', "\text{Cat-cofibré}"] \arrow[r, "\text{source}"] & S \arrow[r, hook] & \mathcal M\ ,
\end{tikzcd}
\tag*{(4.3.6.4)}
\]
où $J$ est Cat-cofibré sur $I$, sa fibre en $i\in\Ob I$ étant isomorphe à $S/x_i$, donc $\card\Fl J_i\leqslant\pi$. Comme par hypothèse $\card\Fl I\leqslant\pi$, on en conclut que $\card\Fl J\leqslant\pi^3=\pi$. (\textbf{N.B.} L'ensemble des flèches $v:j_0\longrightarrow j_1$ de $J$ au dessus d'une flèche $u:i_0\longrightarrow i_1$ de $I$ est “contenu” dans l'ensemble produit $\Ob J_{i_0}\times\Fl J_{i_1}$, par l'application
\[
v\longmapsto(\operatorname{source}v,v':u_*(j_0)\longrightarrow j_1)\ ,
\]
donc de cardinal majoré par $\pi^2$.) On \msp{page 75}a alors
\[
x\ \overset{\text{déf}}{=}\ \varinjlim_I x_i\ =\ \varinjlim_J\ \underbrace{\xi_j}_{=\,\psi(j)}\qquad\quad\mbox{(\footnotemark)}\ ,
\]
\footnotetext{On utilise la propriété de distributivité des $\varinjlim$ , relative à un foncteur Cat-cofibrant
\[
\begin{tikzcd}[ampersand replacement=\&]
J \arrow[r] \arrow[dr] \& I\\
\& \mathcal M\ ,
\end{tikzcd}
\]
explicitée dessus.}%
où cette fois la limite est une limite d'objets de $S$, prise dans $\mathcal M$. Il s'ensuit que $x\in\Filt'^{\pi}$, ce n'est pas tout à fait $x\in\Filt^\pi$. (Ça implique $x\in\Filt^{\pi^{\pi_0}}$, donc $x\in\Filt^\pi$ si $\pi=\pi^{\pi_0}$, p.ex. $\pi=2^c$ avec $c\geqslant\pi_0$.) Mais c'est le moment d'utiliser le fait que les objets de $S$ sont accessibles, et on prendra $\pi_0$ tel que les objets de $S$ soient $\pi_0$-accessibles, et cette fois on prendra quand même $I$ filtrante et grande devant $\pi_0$ (pour nous borner à ce qui est exigé dans la condition b) de loc. cit.). On doit prouver que $\card\,(\Fl(S/x))\leqslant\pi$ (ce qui signifie $x\in\Filt^\pi$). Mais comme les objets de $S$ sont $\pi_0$-accessibles et $I$ grande devant $\pi_0$, on voit aussitôt que
\[
S/x\ \simeq\ \varinjlim_i S/x_i\ ,
\]
% foot: Version 17 décembre 2025

% ---- from der-ch18_p016-020.tex
\setcounter{footnote}{11}

\printedpage{14}
% head: A. Grothendieck : Dérivateurs, Ch. XVIII, Édition des Universités de Montpellier II et Paris VII

d'où
\[
\Fl\,S/x \;\simeq\; \varinjlim_{i} \Fl\,S/x_i \;,
\]
donc
\[
\card\Fl(S/x) \;\leqslant\; (\underbrace{\card I}_{\leqslant\pi})\underbrace{(\sup_{i}\card\Fl(S/x_i))}_{\leqslant\pi} \;\leqslant \pi^2 = \pi \;,
\]

\hfill q.e.d.

\msp{page 76}
% margin: hand-drawn bracket mark in the left margin beside the next two lines
Prouvons c) (\footnote{\textit{Démonstration canulée}, ne marche que si $\pi^{\pi_0} = \pi$\,!!}), i.e. que pour tout $\pi' > \pi \geqslant \pi_0$ et tout $x \in \Ob\Filt^{\pi'}$, on peut écrire
\[
x \;=\; \varinjlim_{I} x_i \qquad I \text{ grande devant } \pi\text{, les } x_i \text{ dans } \Filt^{\pi}\text{, } \card I \leqslant \pi'^{\pi}. \tag*{(4.3.6.5)}
\]
% margin: hand-drawn bracket mark in the left margin beside the next three lines
La démonstration se fait comme dans loc. cit. (p. 151-152), en changeant légèrement la définition de $I$, comme désignant l'ensemble des sous-catégories pleines $i$ de $S/x$ qui sont filtrantes, de cardinal $\leqslant \pi$ \textit{et de plus grandes devant} $\pi_0$ (\footnote{Il n'est pas sûr que $I$ soit grande devant $\pi$, sauf si on suppose $\pi^{\pi_0} = \pi$\,!}). (On pourra supposer $S/x$ filtrante en augmentant $S$ de façon à ce que $S$ soit stable dans $\mathcal{M}$ par petites $\varinjlim$ finies, donc les catégories $S/x$ aussi, ce qui implique qu'elles sont filtrantes. Il est \texttt{[phrase incomplète]})

On a donc prouvé PTA $\Longrightarrow$ c).

\dgrblock{dgrb970}{Why the first filtration argument stops}
\begin{dcomment}{exposition}{p76:On a donc prouvé PTA}{Why the first filtration argument stops}
Grothendieck wants a presentation which controls not only the existence of small pieces but also their number.  The two proposed measurements record different information.  The condition $|\operatorname{Fl}(S/x)|\leq\pi$ counts \emph{all} maps from probes and their relations.  The primed condition asks for one strictly epimorphic family from at most $\pi$ probes.  The first implies the second, but a family which covers $x$ need not furnish lifts of every map from a probe into $x$.

This distinction matters already on page 73.  The estimate $|\operatorname{Fl}(S/x)|\leq\pi^{\pi_0}$ cited from SGA 4 is proved there for a universally epimorphic, pullbackable covering family; see \citet[Expos\'e I, Proposition 7.6 and Proposition 9.13]{SGA4I1972}.  That hypothesis allows a map into $x$ to be pulled back to a covering family and then encoded by maps among the probes.  Accessibility alone is not the hypothesis of that estimate.  Its use in the present argument is therefore not justified by the cited result.  This identifies a missing inference, rather than a counterexample to every possible version of the estimate.

Several parts of the argument do work.  There are only a small set of presheaves on $S$ whose categories of elements have at most $\pi$ arrows, up to isomorphism, so the unprimed layer is essentially small.  If $I$ is $\kappa$-filtered, the objects of $S$ are $\kappa$-presentable, $|\operatorname{Fl}I|\leq\pi$, and $|\operatorname{Fl}(S/x_i)|\leq\pi$, then
\[
 |\operatorname{Ob}(S/\operatorname{colim}_I x_i)|
 \leq\sum_{i\in I}|\operatorname{Ob}(S/x_i)|\leq\pi.
\]
For a fixed arrow $s\to t$ in $S$, the arrows over it are determined by a map $t\to\operatorname{colim}_I x_i$.  Thus the number of all arrows is at most $|\operatorname{Fl}S|\cdot\pi\leq\pi$.  This is a precise version of the calculation on pages 74--75.  It uses filtered colimits of the relevant size, not arbitrary colimits.  The two projective-limit signs on page 74 are source slips; the argument concerns colimits.

The proposed proof of the decomposition axiom on page 76 requires more.  Starting with at most $\pi$ arrows of $S/x$, one must close the resulting subcategory under simultaneous cocones for at most $\pi_0$ constraints.  The number of constraints can be $\pi^{\pi_0}$, so closure need not stay within $\pi$.  A fixed-point hypothesis such as $\pi^{\pi_0}=\pi$ permits the familiar closure construction: adjoin chosen cocones for all these constraints, iterate through an ordinal of cofinality greater than $\pi_0$, and keep cardinality $\pi$.  Without a bound of this kind, the ordered collection of proposed small subcategories has not been shown to be large enough in the required sense.  The manuscript explicitly retracts this step in its footnotes; the sentence claiming that PTA implies condition (c) is not supported by the preceding proof.

The failure can be seen without uncertain cardinal arithmetic.  Take $\mathcal M=\mathbf{Set}$, let $S$ be a skeleton of finite sets, and put $\pi_0=\pi=\aleph_0$.  Let $x$ be an infinite set.  The proposed poset $I$ consists of countable full subcategories of $S/x$ which are $\aleph_1$-filtered.  For each $t\in x$, the singleton subcategory on the map $1\to x$ with value $t$ is such a subcategory.  If countably many distinct such singleton subcategories had an upper bound $K\in I$, $\aleph_1$-filteredness would give a common cocone for their countable family of objects.  Its vertex would be a finite set mapping to $x$ whose image contains all those distinct values, which is impossible.  Thus this particular $I$ is not grand devant $\aleph_0$.  Sets still have the desired decompositions, for instance by their countable subsets; the failed part is the proposed indexing construction.

The repair should preserve that distinction.  The direct construction in the preceding block proves an ind-presentation from accessible probes.  It does not identify the provisional filtration with the later intrinsic filtration by presentability ranks.  The latter comparison is the subject of Section 4.5, and the manuscript itself points forward to Section 4.6 for a corrected proof of the theorem.  An audit can therefore retain the theorem's structural content while recording that this particular numerical proof is abandoned.
\end{dcomment}


\textit{Prouvons c) $\Longrightarrow$ d).} Soit
\[
\pi_1 \;=\; 2^c \qquad \text{avec} \qquad c \;\geqslant\; \pi_0 \;.
\]
On sait alors que
\[
\Filt^{\pi_1}\mathcal{M} \;=\; \mathcal{M}_{\pi_1} \;\overset{\text{déf}}{=}\; \text{sous-catégorie pleine de $\mathcal{M}$ des objets $\pi_1$-accessibles de $\mathcal{M}$}
\]
(loc. cit. 9.16), ce qui implique notamment que $\mathcal{M}_{\pi_1}$ est stable par $\varinjlim$ de cardinal $\leqslant \pi_1$, en vertu de loc. cit. 9.9 (page 145). De plus,
\msp{page 77}
par loc. cit. 9.18 b) (p. 156), on a
\[
\Ind_{\pi_1}(\mathcal{M}_{\pi_1}) \xrightarrow{\;\sim\;} \mathcal{M} \;,
\]
ce qui prouve d).

\textit{Prouvons d) $\Longrightarrow$ PTA.} Par loc. cit. 9.18 a) (p. 156), on voit que les objets de $\mathcal{M}_{\pi_1}$ sont $\pi_1$-accessibles dans $\Ind_{\pi_1}(\mathcal{M}_{\pi_1})$, d'autre part il est clair que $\mathcal{M}_{\pi_1}$ est une sous-catégorie pleine strictement génératrice de $\Ind_{\pi_1}(\mathcal{M}_{\pi_1})$, donc il reste à voir que $\Ind_{\pi_1}(\mathcal{M}_{\pi_1})$ est stable par petites $\varinjlim$ , i.e. le

\textbf{Lemme 4.3.7.} \textit{Soit $\mathcal{M}_\pi$ une $\mathfrak{U}$-catégorie stable par limites inductives de cardinal $\leqslant \pi$, où $\pi$ est un cardinal infini donné. Alors} $\Ind_\pi\mathcal{M}_\pi$ \textit{est stable par petites limites inductives. De plus, $\mathcal{M} = \Ind_\pi\mathcal{M}_\pi$, munie de $\mathcal{M}_\pi \longrightarrow \mathcal{M}$, est solution d'un problème $2$-universel, à savoir
\printedpage{15}
% head: A. Grothendieck : Dérivateurs, Ch. XVIII, Édition des Universités de Montpellier II et Paris VII
que pour toute catégorie $\mathcal{N}$ stable par petites limites inductives, le foncteur restriction induit une équivalence}
\[
\uHom_{!}(\mathcal{M},\mathcal{N}) \xrightarrow{\;\sim\;} \uHom_{!\,\pi}(\mathcal{M}_\pi,\mathcal{N}) \;, \tag*{(4.3.7.1)}
\]
\textit{où le signe ! (resp. $!\,\pi$) désigne la sous-catégorie pleine de la catégorie} $\uHom$\textit{, formée des foncteurs qui commutent aux petites} $\varinjlim$ \textit{(resp. aux limites inductives de cardinal $\leqslant \pi$).}

\msp{page 78}
Prouvons d'abord que les $\varinjlim$ de cardinal $\leqslant \pi$ \textit{dans} $\mathcal{M}_\pi$ sont aussi des $\varinjlim$ \textit{dans} $\mathcal{M} = \Ind_\pi\mathcal{M}_\pi$ (\footnote{Mais \textit{attention}, $\mathcal{M}_\pi \longrightarrow \Ind(\mathcal{M}_\pi)$ ne commute \textit{pas} à ces $\varinjlim$ !}). On a donc
\[
x \xleftarrow{\;\sim\;} \varinjlim_{I} x_i \qquad \textit{dans } \mathcal{M}_\pi \;,
\]
i.e.
\[
\Hom(x,y) \xrightarrow{\;\sim\;} \varprojlim_{I} \Hom(x_i,y) \qquad \text{si } y \in \Ob\mathcal{M}_\pi \;, \tag*{$(*)$}
\]
prouvons que cela reste vrai pour $y$ dans $\mathcal{M}$, i.e. $y$ de la forme
\[
y \;=\; \varinjlim_{J} y_j \;, \qquad J \text{ grande devant } \pi\text{, les } y_j \text{ dans } \mathcal{M}_\pi.
\]
On a en effet
\[
\begin{aligned}
\Hom(x,y) \quad &\simeq \quad \varinjlim_{J} \Hom(x,y_j)\\
&\overset{(*)}{\simeq} \quad \varinjlim_{J}\, \varprojlim_{I} \Hom(x_i,y_j)\\
&\overset{\text{SGA 4 I 9.8}}{\simeq} \quad \varprojlim_{I}\, \varinjlim_{J} \Hom(x_i,y_j)\\
&\simeq \quad \varprojlim_{I} \Hom(x_i,y) \;,
\end{aligned}
\]
au total
\[
\Hom(x,y) \;\simeq\; \varprojlim_{I} \Hom(x_i,y),
\]
q.e.d.

Pour prouver que les petites $\varinjlim$ sont représentables dans $\mathcal{M} = \Ind_\pi(\mathcal{M}_\pi)$, comme elles le sont dans $\Ind(\mathcal{M}_\pi)$ (SGA 4 I 8.9.5 b, on utilise l'hypothèse que les $\varinjlim$ \textit{finies} sont représentables dans $\mathcal{M}_\pi$), il reste à voir que $\Ind_\pi(\mathcal{M}_\pi)$ est stable dans $\mathcal{M}_\pi$ par ces limites.

\dgrblock{dgrb971}{Which colimits the ind-completion remembers}
\begin{dcomment}{examples}{p78:Pour prouver que les petites}{Which colimits the ind-completion remembers}
Grothendieck's warning on page 78 is essential to the next section.  An ind-completion always remembers the original category fully faithfully, but it remembers a colimit \emph{as a colimit} only when the objects allowed as targets impose the appropriate continuity condition.  Here the large filtering condition supplies precisely that condition.

Let $\kappa=\pi^+$, suppose $C$ has colimits of diagrams with fewer than $\kappa$ arrows, and let $x=\operatorname{colim}_{i\in I}x_i$ be such a colimit in $C$.  If $Y=\operatorname{colim}_{j\in J}h_{y_j}$ is a $\kappa$-filtered colimit of representables, then
\[
 \begin{split}
 \operatorname{Hom}(h_x,Y)
 &\cong\operatorname{colim}_{j\in J}C(x,y_j)\\
 &\cong\operatorname{colim}_{j\in J}
          \operatorname{lim}_{i\in I^{\mathrm{op}}}C(x_i,y_j)\\
 &\cong\operatorname{lim}_{i\in I^{\mathrm{op}}}
          \operatorname{colim}_{j\in J}C(x_i,y_j)
 \cong\operatorname{lim}_{i\in I^{\mathrm{op}}}\operatorname{Hom}(h_{x_i},Y).
 \end{split}
\]
The middle interchange is justified because there are fewer than $\kappa$ simultaneous representatives and compatibility equations to choose.  First choose a common upper stage for the representatives; then choose a common upper stage where the equations hold.  Injectivity follows by doing the same for the equations expressing equality.  Thus $C\to\operatorname{Ind}_\kappa C$ preserves the indicated colimit.  Notice that the colimit need not be the pointwise colimit in $C^\wedge$.

One explicit example separates the two completions.  Let $C$ be a skeleton of the category of countable sets, let $\pi=\aleph_0$, and put $\kappa=\aleph_1$.  Countable colimits of countable sets are countable.  Every set $X$ is the $\aleph_1$-filtered union of its countable subsets, since a countable collection of countable subsets has countable union.  A map from a countable set into such a union factors through one stage, and two maps which become equal agree at a later stage.  It follows that the restricted Yoneda functor identifies
\[
 \operatorname{Ind}_{\aleph_1}C\simeq\mathbf{Set}.
\]
In particular, the countable coproduct of the singleton objects is represented by the countable set $\mathbb N$.

In the ordinary filtered completion $\operatorname{Ind}_{\aleph_0}C$, however, consider the filtered diagram of finite subsets $A\subset\mathbb N$.  Its colimit is the presheaf
\[
 F(c)=\operatorname{colim}_{A\subset\mathbb N\text{ finite}}C(c,A)
     =\{f:c\to\mathbb N\mid f(c)\text{ is finite}\}.
\]
This is not $h_{\mathbb N}$: evaluating at $c=\mathbb N$ misses the identity map.  Since finite coproducts of representables agree with those in $C$, this $F$ is the countable coproduct of the singleton in the ordinary ind-completion.  The inclusion of $C$ therefore does not preserve this countable coproduct there.

The proposed final sentence of page 78 cannot establish closure of $\operatorname{Ind}_\pi C$ under colimits \emph{inside} ordinary $\operatorname{Ind}C$ either.  In this example its countable coproduct would be $F$, which is not $\aleph_1$-continuous on $C^\circ$: $F(\mathbb N)$ consists only of finite-image maps, whereas the product of countably many copies of $F(1)=\mathbb N$ contains every sequence.  Section 4.4 correctly turns instead to constructing colimits intrinsically in the restricted completion.
\end{dcomment}


\msp{page 79}

\section*{4.4\quad Étude de $\Ind_\pi(C)$ (pour $C$ petite catégorie stable par $\varinjlim$ de cardinal $\leqslant \pi$)}

\textbf{4.4.1.} Finalement, la démonstration du lemme 4.3.7 n'est pas tellement évidente que je l'anticipais, et il me faut lui consacrer des préliminaires techniques, que je vais réunir ici. Dans toute la suite, $C$ représente une $\mathfrak{U}$-catégorie équivalente à une petite catégorie.

\printedpage{16}
% head: A. Grothendieck : Dérivateurs, Ch. XVIII, Édition des Universités de Montpellier II et Paris VII

\textbf{Proposition 4.4.2.} \textit{Soit $F \in C^{\wedge}$ un foncteur ind-représentable, i.e. tel que $C/F$ soit filtrante. Soit d'autre part $i \longmapsto x_i$ un foncteur $I \longrightarrow C$, où $I$ est une petite catégorie filtrante, et soit}
\[
\alpha \;:\; x \;=\; \varinjlim_{I}{}_{C^{\wedge}} x_i \longrightarrow F \qquad \text{(\footnotemark)} \tag*{(4.4.2.1)}
\]
\footnotetext{La donnée de $\alpha$ équivaut aussi à celle d'un homomorphisme d'ind-objets $(x_i)_{i\in I} \longrightarrow F$, $F$ étant considéré comme ind-objet (généralisé) indexé par $C/F$.}
\textit{un système projectif de flèches $x_i \longrightarrow F$ (}\textbf{N.B.} \textit{la limite inductive $x$ est prise dans $C^{\wedge}$, et est calculée argument par argument), ce qui équivaut à une factorisation $\varphi_\alpha : I \longrightarrow C/F$ de $I \longrightarrow C$ en}
\[
I \xrightarrow{\;\varphi_\alpha\;} C/F \hookxrightarrow[\psi]{\textit{can.}} C \;.
\]
\textit{Ceci dit, pour que $\alpha$ (4.4.2.1) soit un isomorphisme dans} $\Ind'(C) \subset C^{\wedge}$ \textit{(sous-catégorie pleine des foncteurs ind-représentables sur $C$), il faut et il suffit que $\varphi_\alpha$ soit} cofinal\textit{, ou ce qui revient au même, $I$ étant filtrante, qu'il satisfasse les deux conditions habituelles :}
\begin{itemize}[leftmargin=2.5em]
\item[F\,1)] $\forall\, \xi \in \Ob C/F$, $\exists\, i \in \Ob I$ \textit{et une flèche} $\xi \longrightarrow \varphi_\alpha(i)$.
\item[F\,2)] $\forall\, \xi \in \Ob C/F$, \textit{et} $i \in \Ob I$, \textit{et toute double-flèche}

\msp{page 80}
\[
\xi \rightrightarrows \varphi_\alpha(i) \;,
\]
\textit{il existe une flèche $i \xrightarrow{\;u\;} j$ dans $I$ telle que $\varphi_\alpha(u)$ égalise la double-flèche.}
\end{itemize}
\textsc{Démonstration}. Si $\varphi_\alpha$ est cofinal, on en conclut par définition que $\varinjlim \varphi_\alpha \longrightarrow \varinjlim \psi\varphi_\alpha$ est un isomorphisme, or cette flèche n'est autre que $\alpha$. Inversement, supposons que cette flèche soit un isomorphisme, donc $C/F \xleftarrow{\;\sim\;} C/\varinjlim x_i$, et il reste à voir que le foncteur canonique
\[
I \;\longrightarrow\; C/\varinjlim{}_{C^{\wedge}} x_i
\]
est cofinal, i.e. satisfait F\,1) et F\,2). La vérification est $\pm$ tautologique, en utilisant le fait que pour tout $y$ dans $C$, on a
\[
\Hom(y,\, \varinjlim{}_{C^{\wedge}} x_i) \;\xleftarrow{\;\sim\;}\; \varinjlim{}_{\mathrm{Ens}}\Hom(y,x_i),
\]
et du “calcul” des limites inductives filtrantes dans Ens.

\dgrblock{dgrb972}{Cofinality as an elementwise presentation criterion}
\begin{dcomment}{exposition}{p80:et du “calcul” des limites inductives filtrantes dans Ens.}{Cofinality as an elementwise presentation criterion}
The point of Proposition 4.4.2 is to replace a statement about equality of ind-objects by a statement about the category indexing their elements.  That change makes later simultaneous constructions possible: one can arrange factorization and equality at stages, rather than manipulate abstract maps out of a completed category.

Let $F:C^{\mathrm{op}}\to\mathbf{Set}$, and write $C/F$ for its category of elements.  An object is a pair $(c,u)$ with $u\in F(c)$; a morphism $(c,u)\to(d,v)$ is a map $a:c\to d$ satisfying $F(a)(v)=u$.  There is always a canonical presentation
\[
 F\cong\operatorname*{colim}_{(c,u)\in C/F}h_c
 \qquad\text{in }C^\wedge.
\]
A map $\alpha:\operatorname{colim}_{i\in I}h_{x_i}\to F$ is consequently the same as compatible choices $u_i\in F(x_i)$, hence a functor $q:I\to C/F$ over $C$.

If $q$ is final, replacing $C/F$ by $I$ in the canonical presentation gives $\alpha$ as an isomorphism.  Conversely, suppose $I$ is filtered and $\alpha$ is an isomorphism.  For any $(c,u)\in C/F$, the identity
\[
 F(c)\cong\operatorname{colim}_{i\in I}C(c,x_i)
\]
represents $u$ by some arrow $c\to x_i$.  This is F1.  If two arrows $(c,u)\rightrightarrows q(i)$ are given, their underlying maps $c\rightrightarrows x_i$ have the same image in that colimit of sets.  Filteredness supplies a later arrow $i\to j$ at which the two composites agree.  This is F2.  These conditions make every comma category $(c,u)\downarrow q$ nonempty and connected: two objects can first be moved to a common index of $I$, and their remaining discrepancy can then be equalized.  Thus $q$ is final.

For $\kappa=\pi^+$, the same observation yields the clean formulation of Corollary 4.4.5:
\[
 F\in\operatorname{Ind}_\kappa C
 \quad\Longleftrightarrow\quad C/F\text{ is $\kappa$-filtered}.
\]
One implication follows by the presentation criterion and the stability of $\kappa$-filteredness under a final functor from a $\kappa$-filtered category.  For the other, the canonical presentation itself is a $\kappa$-filtered presentation.  If ind-objects are required to use ordered indexing sets, the replacement by a final $\kappa$-directed poset is justified in the next block.  What is not required is that each bounded collection of arrows of $C/F$ lie inside a subcategory with an actual terminal object.  The distinction is exactly the one exposed by the editors' idempotent example.
\end{dcomment}


% margin: vertical bar in the left margin beside Corollaire 4.4.3
\textbf{Corollaire 4.4.3.} \textit{Soit $\pi$ un ordinal infini. Pour que $F$ soit dans} $\Ind'_\pi(C)$ \textit{(image essentielle de} $\Ind_\pi(C)$ \textit{dans $C^{\wedge}$), il faut et il suffit qu'il existe une catégorie ordonnée grande devant $\pi$ $I$, et un foncteur cofinal $I \longrightarrow C/F$.}

Il nous faut donc donner des conditions, sur une catégorie $J$ $(= C/F)$, pour qu'il existe une catégorie ordonnée $I$ grande devant $\pi$, et un foncteur cofinal $I \longrightarrow J$.

\msp{page 81}
Une telle catégorie s'appellera \textit{grande devant} $\pi$. (Il est immédiat que dans le cas où $J$ est elle-même ordonnée, cette terminologie est cohérente avec celle déjà adoptée pour les ensembles ordonnés, i.e. que s'il existe une application croissante cofinale $\varphi : I \longrightarrow J$, avec $I$ grand devant $\pi$, alors $J$ est lui-même grand devant $\pi$.) Une telle catégorie est
\printedpage{17}
% head: A. Grothendieck : Dérivateurs, Ch. XVIII, Édition des Universités de Montpellier II et Paris VII
automatiquement filtrante (SGA 4 I 8.1.3 b)). Notons aussi que si on a un foncteur cofinal $J \longrightarrow J'$, si $J$ est grande devant $\pi$, $J'$ l'est aussi (par transitivité de la cofinalité).

\textbf{Proposition 4.4.4.} \textit{Soient $J$ une $\mathfrak{U}$-catégorie, $\pi$ un cardinal infini. Conditions équivalentes :}
\begin{itemize}[leftmargin=4em]
\item[\textit{a)}] \textit{$J$ est grande devant $\pi$, i.e. il existe un ensemble ordonné grand devant $\pi$ et un foncteur cofinal $\varphi : I \longrightarrow J$.}
\item[\textit{b)}] \textit{$J$ satisfait les deux conditions suivantes :}
\begin{itemize}[leftmargin=4em]
\item[$\mathrm{PF}_{\pi}\,1$] $\mathrm{Ord}\,J$ \textit{est grande devant $\pi$, i.e. pour toute famille $(y_\alpha)_{\alpha\in A}$ d'objets de $J$, avec} $\card A \leqslant \pi$\textit{, il existe un majorant $y$ dans $J$ (i.e. un $y$ et des flèches $y_\alpha \xrightarrow{\;v_\alpha\;} y$).}
\item[$\mathrm{PF}_{\pi}\,2$] \textit{Pour $y_0$, $y_1$ dans $J$ et tout ensemble $(u_\alpha)_{\alpha\in A}$ de flèches $y_0 \xrightarrow{\;u_\alpha\;} y_1$ avec} $\card A \leqslant \pi$\textit{, il existe une flèche $y_1 \xrightarrow{\;u'\;} y_2$ qui les égalise, i.e. telle que $u'u_\alpha = u'u_\beta \;\forall\, \alpha,\beta \in A$. (Il suffit de l'exiger quand} $\card A \geqslant 2$.\textit{)}
\end{itemize}
\item[\textit{c)}] \textit{Pour toute partie $A$ de} $\Fl J$\textit{, avec} $\card A \leqslant \pi$\textit{, existe une sous-catégorie $J'$ de $J$, avec} $\card\Fl J' \leqslant \pi$\textit{, telle que $A \subset \Fl J'$ et que $J'$ admette un objet final.}
\item[\textit{d)}] \textit{Soit $I$ l'ensemble des sous-catégories $J'$ de $J$ telles que} $\card J' \leqslant \pi$ \textit{et ayant un objet final}

\msp{page 82}
\textit{et ordonnons $I$ par inclusion. Alors $I$ est grand devant $\pi$, de plus le foncteur canonique $I \longrightarrow J$ $(J' \longmapsto e_{J'} =$ objet final de $J')$ est cofinal} (\footnote{Introduire $I(J)$ avant l'énoncé de la proposition. \textbf{N.B.} La condition de cofinalité $I(J) \longrightarrow J$ n'est pas automatique, comme on voit en prenant $J = B_G$, $G = \{\pm 1\}$.},\footnote{\textbf{N.B.} L'hypothèse de cofinalité est superflue si $J$ est rigide, i.e. tout endomorphisme d'un objet de $J$ est identique.},\footnote{\ N. Éd.\ \ Dans cette condition (d), le “\textit{de plus le foncteur \dots\ est cofinal}” est ajouté par Grothendieck \textit{a posteriori}, ce qui explique sa vaine tentative de l'établir dans l'implication d $\Longrightarrow$ a ci-dessous, implication qui devient évidente après cet ajout. Par ailleurs, les conditions (c) et (d) ne sont \textit{pas} impliquées par les conditions équivalentes (a) et (b). Par exemple, si $J$ est la catégorie ayant un seul objet et un seul morphisme non identique $p$ tel que $p^2 = p$, alors $J$ satisfait aux conditions (a) et (b), mais \textit{pas} aux conditions (c) et (d).}).
\end{itemize}
(\footnote{Autres conditions :
\begin{itemize}[leftmargin=3em]
\item[e)] Pour tout ensemble $A$ avec $\card A \leqslant \pi$, $\mathrm{diag}_J : J \longrightarrow J^A$ est cofinal.
\item[e$'$)] (Si $A_0$ est un ensemble donné, de cardinal $= \pi$.) $J \longrightarrow J^{A_0}$ est cofinal.
\end{itemize}
Voir aussi 4.9.6.9.1 et 4.9.6.9.2, pages 237, 238.})

\textsc{Démonstration}.
\begin{itemize}[leftmargin=7em,labelwidth=6em,align=left]
\item[a $\;\Longrightarrow\;$ b] vérification immédiate, en utilisant le critère F 1, F 2 de la cofinalité pour $I$ filtrante.
\item[b $\;\Longrightarrow\;$ c] par AQT \texttt{[âne qui trotte]}.
\item[c $\;\Longrightarrow\;$ d] par AQT \texttt{[âne qui trotte]}.
\end{itemize}
% foot: Version 17 décembre 2025

\printedpage{18}
% head: A. Grothendieck : Dérivateurs, Ch. XVIII, Édition des Universités de Montpellier II et Paris VII
\begin{itemize}[leftmargin=7em,labelwidth=6em,align=left]
\item[d $\;\Longrightarrow\;$ a] Considérons l'ensemble ordonné grand devant $\pi$ construit dans l'énoncé de d), et le \textit{foncteur}
\[
I \;\longrightarrow\; J
\]
donné par $J' \longmapsto e_{J'}$ $=$ objet final de $J'$. (Vérification que c'est un foncteur immédiate.) Ce foncteur est surjectif, donc satisfait la condition F 1 de cofinalité. La condition F 2 résulte aussi de d) en regardant les trois objets $J'$, $j \xrightarrow{\;u_1\;} e_{J'}$, $j \xrightarrow{\;u_2\;} e_{J'}$ de $I$,
\[
\begin{tikzcd}[column sep=small, row sep=small]
J', & e_{J'} \\
j \arrow[ur, shift left, "u_1"] \arrow[ur, shift right, "u_2"'] &
\end{tikzcd}
\]
et en prenant un $J''$ qui les majore tous trois, alors $e_{J'} \longrightarrow e_{J''}$ égalise $u_1$, $u_2$ (\footnote{Un peu bref, le cas où $e_{J'} = j$ échappe à cet argument, et demande qu'on \texttt{[phrase incomplète]}}).
\end{itemize}

% margin: vertical bar in the left margin beside Corollaire 4.4.5
\textbf{Corollaire 4.4.5.} \textit{Soit $F \in \Ob C^{\wedge}$. Pour que $F$ appartienne à} $\Ind'_\pi(C)$\textit{, il faut et il suffit que $C/F$ satisfasse à l'ensemble des quatre conditions équivalentes de 4.4.4.}

\textbf{Corollaire 4.4.6.} \textit{Soit $J$ une catégorie stable par} $\varinjlim$ \textit{de cardinal $\leqslant \pi$. Alors $J$ est grande devant $\pi$.}

En effet, les conditions $\mathrm{PF}_{\pi}\,1$, $\mathrm{PF}_{\pi}\,2$ sont visiblement vérifiées.

\dgrblock{dgrb973}{A corrected terminal-subcategory criterion}
\begin{dcomment}{development}{p82:En effet, les conditions}{A corrected terminal-subcategory criterion}
Grothendieck is trying to turn simultaneous solvability of equations into a particularly concrete indexing poset: its elements would be small subcategories already possessing a final object.  The editors explain why this construction is too restrictive.  We first give the correct general replacement, then show precisely when his proposed one works.  The filtered-category argument is classical; the second statement isolates and proves the repair suggested by his rigidity footnote.  We checked the formulation in SGA 4 I, Section 8.1, the present edition's note 18, and the later product criterion cited below; we did not find this exact bounded rigid formulation there.  The label on the second result records our proof of that formulation, not a claim of priority.

For a small category $J$ and an infinite cardinal $\pi$, put $\kappa=\pi^+$.  Conditions $\mathrm{PF}_\pi1$ and $\mathrm{PF}_\pi2$ are equivalent to saying that every diagram with fewer than $\kappa$ arrows has a cocone.  To prove the nontrivial direction, first choose a common target for all objects in the diagram.  For each arrow $a:i\to j$ there are then two maps from its source to that target, namely the chosen map at $i$ and the composite through $j$.  Use $\mathrm{PF}_\pi2$ to choose a map from the common target to an object $z_a$ equalizing that individual pair, then choose a common target $z$ of all the $z_a$.  The resulting maps from the original common target to $z$ need not coincide; apply $\mathrm{PF}_\pi2$ to this family.  Their common composite equalizes every required pair.  If there are no compatibility pairs, the first common target suffices.  The empty diagram is included in $\mathrm{PF}_\pi1$.

Such a small $J$ admits a final functor from a $\kappa$-directed poset.  Here is a construction which does not presume the false terminal-subcategory condition.  Begin with a discrete poset $P_0$ having a vertex over each object of $J$.  Inductively maintain a functor $q_\alpha:P_\alpha\to J$ such that every principal downset in $P_\alpha$ has fewer than $\kappa$ elements.  At the next stage, for every subset $A\subset P_\alpha$ of size less than $\kappa$, adjoin a new vertex above its downward closure.  That closure has size less than $\kappa$ by regularity, so a chosen cocone of its image defines the new vertex and all arrows into it.  Also, for every $p\in P_\alpha$ and every parallel pair $u,v:j\rightrightarrows q_\alpha(p)$, adjoin a vertex above the downset of $p$ whose transition map equalizes $u,v$.  This is possible by applying the cocone condition to that downset with a new formal source and those two arrows.  Distinct new vertices are left incomparable.  At limit stages take unions.

After $\kappa$ stages the union $P$ is a set, all its principal downsets remain of size less than $\kappa$, and any subset of $P$ of size less than $\kappa$ occurs at one earlier stage and hence has an upper bound.  Thus $P$ is $\kappa$-directed.  The functor $q:P\to J$ satisfies F1 because of $P_0$, and F2 because every parallel-pair requirement is addressed at a successor stage.  The preceding presentation criterion proves that $q$ is final.  This supplies the implication $b\Rightarrow a$ of Proposition 4.4.4 in the small case needed for $J=C/F$.

\textbf{Proposition (the rigid repair).}\ours{}\quad
Suppose $J$ is small and every endomorphism of every object of $J$ is an identity.  Then $J$ is $\kappa$-filtered if and only if every set of at most $\pi$ arrows of $J$ is contained in a subcategory with at most $\pi$ arrows and a terminal object.  In that event, the poset of these terminal subcategories is $\kappa$-directed, and the functor choosing their terminal objects is final.

\emph{Proof.} Let $A$ be the prescribed set of arrows.  Close it under identities and finite composition to obtain a subcategory $K$ with at most $\pi$ arrows.  Choose a cocone $\eta:K\to j$ in $J$.  Adjoin $j$, the arrows $\eta_k:k\to j$, and the identity of $j$.  If $j$ was not in $K$, the generated subcategory has no nonidentity arrow leaving $j$, and every path ending at $j$ equals the appropriate $\eta_k$ by naturality.  Hence $j$ is terminal.  If $j$ already belongs to $K$, rigidity forces $\eta_j=1_j$.  For an old arrow $a:k\to j$, naturality gives $a=\eta_j a=\eta_k$.  All newly generated paths ending at $j$ reduce to the same arrow, again by naturality.  Thus $j$ is terminal in this case too.  The enlarged subcategory still has at most $\pi$ arrows.

Conversely, the arrow set of a diagram of at most $\pi$ arrows is contained in such a terminal subcategory, whose terminal object gives a cocone.  For directedness, take at most $\pi$ terminal subcategories, form the union of their arrow sets, and apply the proved condition.  Their union has at most $\pi$ arrows.  To prove finality of the terminal-object functor, every object of $J$ occurs in its singleton subcategory, giving F1.  Given a terminal subcategory $K$ and two maps $j\rightrightarrows e_K$, enlarge the arrow set of $K$ by this pair and choose a terminal subcategory containing it.  The unique map from $e_K$ to the new terminal object equalizes the pair, giving F2.  This completes the proof.

The rigidity hypothesis addresses the precise obstruction.  In the one-object category with arrows $1,p$ and $p^2=p\ne1$, every family of parallel arrows is equalized by postcomposition with $p$, so the category is $\kappa$-filtered for every infinite $\kappa$.  A subcategory containing $p$ cannot have a terminal object: its only object has two endomorphisms.  Its only terminal subcategory is therefore the identity-only subcategory.  The resulting functor from the singleton indexing poset is not final, since the relevant comma category has two disconnected objects.  This proves the editors' failure of both (c) and (d), and explains why Corollary 4.4.5 must refer to (a) or (b), rather than to all four printed conditions.
\end{dcomment}


\dgrblock{dgrb974}{The infinite-product criterion}
\begin{dcomment}{later}{p82:En effet, les conditions}{The infinite-product criterion}
The diagonal conditions in note 19 point toward a valid alternative.  For a small $J$, finality of $J\to J^A$ for every $|A|\leq\pi$, including $A=\varnothing$, means that $J$-colimits of sets commute with these products.  Since $\pi$ is infinite, \citet[Theorem 3.1]{AdamekKoubekVelebil2000} identifies this condition with $\pi^+$-filteredness.  The restriction to infinite $\pi$ matters: finite-product commutation alone is the weaker condition of siftedness.

The single-power condition (e$'$) needs a nonemptiness clause.  For $J=\varnothing$ and a nonempty set $A_0$, both $J$ and $J^{A_0}$ are empty, so the diagonal is final vacuously.  For nonempty $J$, finality of the indicated diagonal first forces connectedness.  Padding a smaller family with constant singleton diagrams then reduces its product comparison to the $A_0$ comparison.  Thus (e$'$) recovers (e) with that clause added.
\end{dcomment}


\msp{page 83}
\textbf{4.4.6.} Je me donne un foncteur % [sic] 4.4.6. -> 4.4.7. (?), numéro déjà employé par le Corollaire 4.4.6 qui précède
\[
J \xrightarrow{\;\varphi\;} \Ind'_\pi C \;, \tag*{(4.4.6.1)}
\]
$J$ une petite catégorie, et on cherche une $\varinjlim$ dans $\Ind'_\pi C$, en postulant au besoin l'existence des $\varinjlim$ de cardinal $\leqslant \pi$ dans $C$, et l'hypothèse $\card\Fl J \leqslant \pi$. Le procédé le plus naturel pour ramener la question de l'existence des $\varinjlim$ dans $\Ind'_\pi C$, à la question similaire dans $C$, est explicité dans SGA 4 I 8.8.3. On considère l'inclusion pleinement fidèle
\[
C \;\hookxrightarrow{\;\alpha\;}\; \Ind'_\pi C \;, \tag*{(4.4.6.2)}
\]
laquelle commute aux $\varinjlim$ de cardinal $\leqslant \pi$ (celles qui sont représentables dans $C$, cf. p. 78 et référence à SGA 4 I 9.8), d'où
\[
\uHom(J,C) \;\hookxrightarrow{\;\alpha\;}\; \uHom(J,\Ind'_\pi C) \tag*{(4.4.6.3)}
\]
également pleinement fidèle, et commutant aux $\varinjlim$ de cardinal $\leqslant \pi$ (représentables dans le premier membre). Cette flèche s'insère dans le carré commutatif relatif à la donnée d'un $j \in \Ob J$, et des foncteurs d'inclusion
\[
e \xrightarrow{\;\varepsilon_j\;} J \;, \qquad e_j(e) \;=\; j \;; \tag*{(4.4.6.4)} % [sic] e_j(e) -> \varepsilon_j(e), la flèche est notée \varepsilon_j
\]
% foot: Version 17 décembre 2025

% ---- from der-ch18_p021-025.tex
\printedpage{19}
% head: A. Grothendieck : Dérivateurs, Ch. XVIII, Édition des Universités de Montpellier II et Paris VII
% foot: Version 17 décembre 2025

\msp{page 84}
\[
\begin{tikzcd}
\uHom(J,C) \arrow[r, "\alpha^J"] \arrow[d, "\varepsilon_j^*"] & \uHom(J,\Ind'_\pi C) \arrow[d, "{\varepsilon_j^{*\prime}}"] & \ni \varphi \arrow[d, mapsto] \\
C \arrow[r, "\alpha"] & \Ind'_\pi(C) & \ni \varphi(j)\ ,
\end{tikzcd}
\qquad (\text{\footnotemark[21]})
\tag*{(4.4.6.2)}
\]
\footnotetext[21]{Les foncteurs $\varepsilon_j^*$, $\varepsilon_j^{*\prime}$ sont les foncteurs évaluation en $j$.}
lequel induit un foncteur
\[
\underbrace{\uHom(J/C)/\varphi}_{\overset{\text{déf}}{=}\, I} \xrightarrow{\varepsilon_j^\varphi} C/\varphi(j)\ . \tag*{(4.4.6.3)} % [sic] \uHom(J/C) -> \uHom(J,C), everywhere else written Hom(J,C)
\]
Si $C$ est stable par $\varinjlim$ relatives à un ensemble $K$ du type de conditions $\leqslant \pi$ (p.ex. $\varinjlim$ finies), il en est de même de $\uHom(J,C)$, et aussi de $\uHom(J,C)/\varphi$, et de $C/\varphi(j)$, en vertu du fait que $\alpha^J$ et $\alpha$ commutent aux $\varinjlim$ du type envisagé, de plus (4.4.6.3) commute aux limites en question. Supposons donc $C$ stable tout au moins par $\varinjlim$ finies, donc les catégories source et but dans (4.4.6.3) sont filtrantes. (Le fait que le foncteur commute aux $\varinjlim$ finies ne semble pas utile dans le contexte présent.) Cela nous servira pour avoir le critère commode de cofinalité SGA 4 I 8.1.3 b) par F\,1, F\,2.

\dgrblock{dgrb975}{One index category for an entire diagram}
\begin{dcomment}{exposition}{p84:Cela nous servira pour avoir}{One index category for an entire diagram}
The next difficulty is coherence.  An individual ind-object has many filtered presentations, but presentations chosen separately at the vertices of a diagram need not carry its arrows and relations.  Grothendieck uses the category of all coherent approximations at once:
\[
 I=\operatorname{Fun}(J,C)/\varphi.
\]
An object of $I$ consists of a $C$-valued diagram $D$ and a natural transformation $\alpha D\to\varphi$, where $\alpha:C\to\operatorname{Ind}_\pi C$.  A morphism is a natural transformation over $\varphi$.  Evaluating at $j$ gives $\varepsilon_j^\varphi:I\to C/\varphi(j)$.  In general $I$ is a category, not an ordered set; the phrase calling it an ordered set on page 85 is not needed by the argument.

The purpose of Lemma 4.4.9 is to show that each evaluation is final.  We can prove its stronger assertion directly and avoid the temporary dependency on coproducts of arbitrary ind-objects.  Assume that $C$ admits the copowers indexed by the sets $J(j,j')$, all of cardinality at most $\pi$.  Define
\[
 (L_jc)(j')=\coprod_{a\in J(j,j')}c.
\]
An arrow $b:j'\to j''$ carries the summand indexed by $a$ to the summand indexed by $ba$.  This makes $L_jc$ a diagram in $C$.  A natural transformation $L_jc\to D$ is uniquely determined by the map from the identity summand $c\to D(j)$; its component on the $a$-summand must be $D(a)$ composed with that map.  Hence $L_j$ is left adjoint to evaluation on $C$-valued diagrams.

Page 78 proves that $\alpha$ preserves these particular coproducts.  Therefore, for any $\varphi:J\to\operatorname{Ind}_\pi C$, the same calculation gives
\[
 \operatorname{Nat}(\alpha L_jc,\varphi)
 \cong\operatorname{Hom}(\alpha c,\varphi(j)).
\]
This calculation uses coproducts of copies of the representable object $c$ only.  It does not require the still-pending existence of coproducts of arbitrary objects of $\operatorname{Ind}_\pi C$.

For an object $(c,f:\alpha c\to\varphi(j))$ of $C/\varphi(j)$, let $\widetilde f:\alpha L_jc\to\varphi$ be the corresponding natural transformation.  Set
\[
 R_j(c,f)=(L_jc,\widetilde f)\in I.
\]
A map from $R_j(c,f)$ to $(D,\eta)$ in $I$ is a transformation $\theta:L_jc\to D$ with $\eta\alpha\theta=\widetilde f$.  By the preceding calculation, this is exactly a map $g:c\to D(j)$ with $\eta_j\alpha(g)=f$, which is a map $(c,f)\to\varepsilon_j^\varphi(D,\eta)$ in the slice.  These bijections are natural, proving $R_j\dashv\varepsilon_j^\varphi$.  Every comma category $(c,f)\downarrow\varepsilon_j^\varphi$ has the resulting universal object as an initial object, so evaluation is final.

If $C$ has all colimits of cardinality at most $\pi$, then so does $I$: take a colimit of approximating diagrams pointwise in $C$, use preservation by $\alpha$, and form the induced map to $\varphi$.  Thus $I$ is $\pi^+$-filtered.  Finality of each evaluation gives
\[
 \varphi(j)\cong\operatorname*{colim}_{i\in I}\alpha D_i(j).
\]
If $|\operatorname{Fl}J|\leq\pi$, form $c_i=\operatorname{colim}_{j\in J}D_i(j)$ in $C$.  Testing against any target $Z$ and interchanging two limits of Hom sets gives
\[
 \operatorname{Hom}(\operatorname{colim}_{i\in I}\alpha c_i,Z)
 \cong\operatorname{lim}_{j\in J^{\mathrm{op}}}\operatorname{Hom}(\varphi(j),Z).
\]
The left colimit exists by its $\pi^+$-filtered presentation.  It therefore supplies the desired $J$-colimit in $\operatorname{Ind}_\pi C$.  This proves the content of Lemma 4.4.7 and Corollary 4.4.8 after their explicitly stated cofinality reservation has been discharged.
\end{dcomment}


\msp{page 85}
Rappelons que la cofinalité des foncteurs (4.4.6.3), quand elle a lieu ( ?), implique le fait que $\varphi$ est une limite inductive du foncteur composé
\[
\begin{tikzcd}[row sep=small]
I \arrow[r, "\psi_J"] \arrow[d, equal] & \uHom(J,C) \arrow[r, "\alpha^J"] \arrow[dd, dotted, "\varepsilon_j^*"] & \uHom(J,\Ind'_\pi(C)) \arrow[dd, dotted, "{\varepsilon_j^{*\prime}}"] \\
\uHom(J,C)/\varphi \arrow[d, "\varepsilon_j^\varphi"'] & & \\
C/\varphi(j) \arrow[r, dotted, "\psi_j"] & C \arrow[r, dotted, "\alpha"] & \Ind'_\pi(C)\ ,
\end{tikzcd}
\tag*{(4.4.6.4)}
\]
suivant l'ensemble ordonné filtrant $I$, plus précisément que $\varphi$ est la $\varinjlim$ “argument par argument en $j$”, ce qui signifie que pour tout $j \in \Ob J$, $\varphi(j) = \varepsilon_j^*(\varphi)$ est limite inductive de $\varepsilon_j^*(\alpha^J\psi) = (\alpha\psi_j)\varepsilon_j^\varphi$. (Or on sait en effet que $\varphi(j)$ est la limite de $\alpha\psi$, puisque $\varphi(j) \in \Ob C^\wedge$ est ind-représentable, donc aussi celle de $(\alpha\psi_j)\varepsilon_j^\varphi$ si $\varepsilon_j^\varphi$ est cofinal.)

Ainsi, admettant la cofinalité de (4.4.6.3) et interprétant $\psi_J$ dans (4.4.6.4)

\msp{page 86}
\[
\psi_J \ :\ I \longrightarrow \uHom(J,C)
\]
comme un foncteur
\[
I \times J \xrightarrow{\Psi} C\ , % [sic] \Psi -> \psi, the functor I x J -> C is called \psi in the next formula and in (4.4.7.1)
\]
on trouve la “représentation indicielle” de $\varphi$
\[
\varphi(j) \ =\ \underbrace{\varinjlim_{i\in I} \psi(i,j)}_{\substack{\text{limite inductive}\\ \text{prise dans } \Ind'_\pi(C)\\ \text{d'objets de } C \hookrightarrow \Ind'_\pi(C)}} \qquad ,
\]

\printedpage{20}
% head: A. Grothendieck : Dérivateurs, Ch. XVIII, Édition des Universités de Montpellier II et Paris VII
% foot: Version 17 décembre 2025
d'où
\[
\begin{aligned}
\underbrace{\varinjlim_{j\in J}{}_{\,\Ind'_\pi C}\varphi(j)}_{\text{sous réserve d'existence}} \quad &= \quad \varinjlim_{j\in J}\, \varinjlim_{i\in I} \psi(i,j)\\
&= \quad \varinjlim_{i\in I}\, \varinjlim_{j\in J} \psi(i,j)\ ,
\end{aligned}
\]
toutes les limites inductives étant prises dans $\Ind'_\pi(C)$. Mais si les $\varinjlim$ de type $J$ sont représentables dans $C$, et $\card J \leqslant \pi$, on sait que la $\varinjlim\limits_{j\in J} \psi(i,j)$ \textit{dans} $C$ est aussi leur $\varinjlim$ dans $\Ind'_\pi C$, et on trouve donc que $\varinjlim\limits_{j\in J} \varphi(j)$ existe dans $\Ind'_\pi(C)$ si et seulement si $\varinjlim\limits_{i\in I} \Psi(i)$ existe dans $\Ind'_\pi(C)$, où pour $i \in I$, on pose
\[
\Psi(i) \ =\ \varinjlim_{j\in J} \psi(i,j) \qquad \textit{dans } C.
\]

\msp{page 87}
Résumons ce que nous avons obtenu :

\textbf{Lemme 4.4.7.} \textit{Soit $C$ une petite catégorie, $\pi$ un cardinal infini, $J$ une catégorie telle que $\card \Fl J \leqslant \pi$, $\varphi : J \longrightarrow \Ind'_\pi(C)$ un foncteur. Supposons les $\varinjlim$ de type $J$ existant dans $C$, et que les foncteurs $\varepsilon_j^\varphi$ ($j \in \Ob J$) de (4.4.6.3) soient \textup{cofinaux} (ce qui s'explicite de façon simple surtout si source et but sont filtrants, chose acquise si $C$ est stable par $\varinjlim$ finies). Alors $\varphi$ a une “représentation indicielle canonique”}
\[
\psi \ :\ I \times J \longrightarrow C \tag*{(4.4.7.1)}
\]
\textit{avec}
\[
I \ =\ \uHom(J,C)/\varphi\ , \tag*{(4.4.7.2)}
\]
\textit{i.e. on a l'isomorphisme fonctoriel en $j \in \Ob J$}
\[
\varphi(j) \ \simeq\ \varinjlim_{I} \psi(i,j) \qquad (\text{\footnotemark[22]})\ . \tag*{(4.4.7.3)}
\]
\footnotetext[22]{limite inductive dans $\Ind'_\pi$ d'objets de $C$.}
\textit{Posons pour tout $i \in \Ob I$}
\[
\Psi(i) \ =\ \varinjlim_{J} \psi(i,j) \qquad (\text{\footnotemark[23]})\ . \tag*{(4.4.7.4)}
\]
\footnotetext[23]{limite inductive prise dans $C$, mais valable aussi dans $\Ind'_\pi C$.}
\textit{Alors $\varinjlim\limits_{j\in J} \varphi(j)$ (dans $\Ind'_\pi C$) existe si et seulement si}
\msp{page 88}
\textit{la $\varinjlim\limits_{i\in I} \Psi(i)$ existe (dans $\Ind'_\pi C$ également), et les deux sont égales.}

\textbf{Corollaire 4.4.8.} \textit{Supposons que $C$ soit stable par $\varinjlim$ de cardinal $\leqslant \pi$ (pas seulement par celles de type $J$). Alors $\uHom(J,C)$ et $I = \uHom(J,C)/\varphi$ ont la même propriété, $I$ est grande devant $\pi$, et $\varinjlim\limits_{j\in J} \varphi(j) \ =\ \varinjlim\limits_{i\in I} \Psi(i)$ existe dans $\Ind'_\pi(C)$, sous réserve de cofinalité des foncteurs (4.4.6.3).}

\printedpage{21}
% head: A. Grothendieck : Dérivateurs, Ch. XVIII, Édition des Universités de Montpellier II et Paris VII
% foot: Version 17 décembre 2025
Il suffit en effet de vérifier ce

\textbf{Lemme 4.4.8.1.} \textit{Soit $I$ une catégorie grande devant $\pi$, et $\Psi : I \longrightarrow C$ un foncteur. Alors $\varinjlim\limits_{i} \Psi(i)$ existe dans $\Ind'_\pi(C)$, et c'est la limite inductive dans $C^\wedge$, ou encore “c'est” l'ind-objet grand devant $\pi$ représenté par $\Psi$.} (En d'autres termes (\footnote[24]{C'est abusif, cette stabilité est un énoncé nettement plus fort! Je le prouve page 99.}), $\Ind'_\pi(C) \subset C^\wedge$ est une sous-catégorie strictement pleine \textit{stable} par $\varinjlim$ grandes devant $\pi$.)

\msp{page 89}
Ce lemme est une tautologie : si $F$ désigne la limite inductive dans $C^\wedge$, il est clair par définition (en prenant un foncteur $I_0 \longrightarrow I$ cofinal, avec $I_0$ \textit{ordonnée} grande devant $\pi$, de sorte que la $\varinjlim$ peut être prise suivant $I_0$) que $F$ est dans $\Ind'_\pi(C)$, et alors c'est a fortiori une $\varinjlim$ dans $\Ind'_\pi C$.

Pour vérifier l'existence des $\varinjlim$ de cardinal $\leqslant \pi$ dans $\Ind'_\pi C$, dans le cas où $C$ lui-même est stable par lesdites limites, il suffit donc, grâce au corollaire 4.4.8, de prouver la cofinalité des foncteurs (4.4.6.3). Cela va résulter du

\textbf{Lemme 4.4.9.} \textit{Soient $C$ une petite catégorie, $\pi$ un cardinal infini, $J$ une petite catégorie telle que pour tous $j, j' \in \Ob J$, on ait $\card \Hom(j,j') \leqslant \pi$. Sous ces conditions, pour tout $j \in \Ob J$, le foncteur (4.4.6.3) \textup{admet un adjoint à gauche,} donc a fortiori il est cofinal (car les catégories colocalisées au dessus de $C/\varphi(j)$ ont un objet initial, a fortiori elles sont $0$-connexes).}

\msp{page 90}
\textsc{Démonstration de} 4.4.8. Revenons au diagramme (4.4.6.2), dont on déduit (4.4.6.3), nous allons le réécrire de façon simplifiée :
\[
\begin{tikzcd}
B \arrow[r, "\beta"] \arrow[d, shift left, "\varepsilon"] & B' \arrow[d, shift left, "\varepsilon'"] & \ni y'_0 \arrow[d, mapsto] \\
C \arrow[r, "\gamma"] \arrow[u, dotted, shift left, "\lambda"] & C' \arrow[u, dotted, shift left, "\lambda'"] & \ni z'_0\ .
\end{tikzcd}
\tag*{(4.4.9.1)}
\]
Je dis que $\varepsilon$, $\varepsilon'$ ont \textit{des adjoints à gauche} $\lambda$ resp. $\lambda'$, et que l'homomorphisme canonique de “changement de base”
\[
\lambda'\gamma \longrightarrow \beta\lambda \tag*{(4.4.9.2)}
\]
est un \textit{isomorphisme.} En effet, l'existence des adjoints de $\varepsilon$, $\varepsilon'$, à savoir des foncteurs $\varepsilon_{j\,!}$ et $\varepsilon'_{j\,!}$ relatifs à l'inclusion
\[
\varepsilon_j \ :\ e \longrightarrow J
\]
et aux catégories de foncteurs $\uHom(J,C)$ resp. $\uHom(J,\Ind'_\pi C)$, résultant de l'existence des sommes de type $\Hom(j,j')$ ($j' \in \Ob J$) dans $C$ (déjà acquise), \textit{et dans} $\Ind'_\pi C$\textit{ ;} et le fait que le foncteur $C \longrightarrow \Ind'_\pi C$ commute
\msp{page 91}
auxdites sommes (chose également déjà acquise, puisque $\card\, (\Hom(j,j')) \leqslant \pi$) implique alors que (4.4.4.2) est un isomorphisme. Moyennant la vérification en suspens (l'existence % [sic] (4.4.4.2) -> (4.4.9.2), the isomorphism in question is the change-of-base arrow (4.4.9.2)
\printedpage{22}
% head: A. Grothendieck : Dérivateurs, Ch. XVIII, Édition des Universités de Montpellier II et Paris VII
% foot: Version 17 décembre 2025
des sommes de cardinal $\leqslant \pi$ dans $\Ind'_\pi(C)$), le lemme 4.4.9 résultera du résultat plus général :

\textbf{Lemme 4.4.10.} \textit{Considérons un carré commutatif (4.4.9.1) (traits pleins) de catégories, où $\varepsilon$ admet un adjoint à gauche $\lambda$, $\varepsilon'$ un adjoint à gauche $\lambda'$, et la flèche de changement de base (4.4.9.2) est un isomorphisme. Alors pour tout $y'_0$ dans $B'$, posant $\varepsilon'(y'_0) = z'_0$, le foncteur induit par $(\varepsilon, \varepsilon')$,}
\[
\overline{B} \ =\ B/y'_0 \xrightarrow{\bar{\varepsilon}} \overline{C} = C/z'_0\ ,
\]
\textit{admet un adjoint à gauche $\overline{\lambda}$, donné par}
\[
\begin{tikzcd}[column sep=small]
\overline{\lambda}(\gamma(z) \xrightarrow{v} z'_0) \ \ = & \Bigl(\lambda'\gamma(z) \simeq \beta(\lambda(z)) \arrow[r, "{\lambda'(v)}"] \arrow[ddr, "{\substack{\overline{\lambda}(v) \,=\, \text{composé}\\ \text{adj} \circ \lambda'(v)}}"'] & \lambda'(z'_0) \arrow[d, equal] \\
 & & \lambda'\varepsilon'(y'_0) \arrow[d, "\text{adj.}"] \\
 & & y'_0\Bigr)\ \ .
\end{tikzcd}
\]

\msp{page 92}
La démonstration se fait par un AQT \texttt{[âne qui trotte]} fastidieux (j'avoue ne pas avoir poussé jusqu'au dernier détail \ldots).

\dgrblock{dgrb976}{The adjunction on categories of presentations}
\begin{dcomment}{exposition}{p92:La démonstration se fait par un AQT}{The adjunction on categories of presentations}
Grothendieck needs to present an entire diagram by objects of $C$, while keeping the prescribed value of that diagram at a chosen vertex. The passage to comma categories in Lemma 4.4.10 is what makes these two requirements compatible. Its proof is indeed formal, but the formal statement includes the invertibility of the specified mate, not merely the existence of four functors making a square commute.

Write the square as $\varepsilon'\beta=\gamma\varepsilon$, and write $\eta,\epsilon$ and $\eta',\epsilon'$ for the units and counits of $\lambda\dashv\varepsilon$ and $\lambda'\dashv\varepsilon'$. Let
\[
\chi:\lambda'\gamma\longrightarrow\beta\lambda
\]
be the invertible mate used in the source. Fix $y'\in B'$. An object of $(\gamma\downarrow\varepsilon'y')$ is a pair $(c,v)$ with $v:\gamma c\to\varepsilon'y'$. Define
\[
q_v=\epsilon'_{y'}\,\lambda'(v)\,\chi_c^{-1}:
\beta\lambda c\longrightarrow y'.
\]
Thus the proposed left adjoint sends $(c,v)$ to $(\lambda c,q_v)$. On a morphism $a:(c,v)\to(d,w)$, the equation $w\gamma a=v$ and naturality of $\chi$ show that $\lambda a$ is a morphism of the required comma category. This already supplies the functor whose existence the displayed formula asserts.

We check the adjunction by retaining the equation defining a morphism in the comma category. Given $(b,u)$ with $u:\beta b\to y'$, an arrow from $(\lambda c,q_v)$ to $(b,u)$ is an arrow $f:\lambda c\to b$ such that $u\beta f=q_v$. Its transpose under $\lambda\dashv\varepsilon$ is
\[
k=\varepsilon(f)\eta_c:c\longrightarrow\varepsilon b.
\]
The definition of the mate gives
\[
\varepsilon'(\chi_c)\eta'_{\gamma c}=\gamma(\eta_c).
\]
Consequently the transpose, under $\lambda'\dashv\varepsilon'$, of $u\beta f\chi_c$ is $\varepsilon'(u)\gamma(k)$. The transpose of $q_v\chi_c$ is $v$. Since transposition is a bijection, the required equation on $f$ is equivalent to
\[
\varepsilon'(u)\gamma(k)=v.
\]
This is exactly the condition that $k$ be a morphism from $(c,v)$ to the image of $(b,u)$ in $(\gamma\downarrow\varepsilon'y')$. Both transposition bijections are natural in the objects, so this proves the comma-category adjunction, including its action on morphisms.

In the application, $\varepsilon$ is evaluation at $j$. Its left adjoint sends $c$ to the diagram $j'\mapsto\coprod_{J(j,j')}c$. The comparison $\chi$ is invertible because the embedding of $C$ preserves these coproducts. Lemmas 4.4.11 and 4.4.12 establish their existence in the completion without presupposing general cocompleteness. A functor with a left adjoint is final: for each object in its target, the unit exhibits an initial object in the corresponding comma category. This explains why the adjunction provides precisely the cofinality needed for the simultaneous presentation.
\end{dcomment}


Pour terminer la démonstration de 4.4.9, il reste à prouver un dernier ( ?)

\textbf{Lemme 4.4.11.} \textit{Soit $C$ une petite catégorie, stable par sommes de cardinal $\leqslant \pi$. Alors $\Ind'_\pi C$ est stable par sommes de cardinal $\leqslant \pi$ (et, pour mémoire, le foncteur $\alpha : C \longrightarrow \Ind'_\pi C$ y commute).}

Appliquant le corollaire 4.4.8 au lemme 4.4.7, on est ramené à prouver la cofinalité de (4.4.6.3) quand $J$ est discrète, plus le fait que $\uHom(J,C)/\varphi$ est grande devant $\pi$, dans ce cas. Mais la donnée de $\varphi$ revient à celle d'une famille $(\varphi_j)_{j\in J}$ d'objets de $\Ind'_\pi(C)$, et le foncteur (4.4.6.3) s'explicite comme la projection
\[
\prod_{j\in J} (C/\varphi_j) \longrightarrow C/\varphi_j
\]
du produit total sur un des facteurs. Donc on est ramené à ceci :

\msp{page 93}
\textbf{Lemme 4.4.12.} \textit{Soit $(C_j)_{j\in J}$ une famille de petites catégories grandes devant $\pi$, indexée par un petit ensemble $J$. Alors $\prod C_j$ est grande devant $\pi$, et les projections (pour $j_0 \in J$) $\prod C_j \longrightarrow C_{j_0}$ sont cofinales.}

Prouvons d'abord la cofinalité, qui est le plus important pour nous. C'est évident si $J = \{j_0\}$, sinon on a $J = J' \cup \{j_0\}$, et $\prod_{j\in J} C_j \simeq C' \times C_{j_0}$, où $C' = \prod_{j\in J'} C_j$, donc il faut voir que $C' \times C_{j_0} \longrightarrow C_{j_0}$ est cofinal, i.e. que toute catégorie $\xi\backslash(C' \times C_{j_0})$ (pour $\xi \in \Ob C_{j_0}$) est 0-connexe, or elle est isomorphe à $C' \times (\xi\backslash C_{j_0})$, et $\xi\backslash C_{j_0}$ étant 0-connexe, il reste à voir que $C'$ est 0-connexe. On fera attention à cet égard qu'un produit \textit{infini} de catégories 0-connexes n'est pas forcément 0-connexe, donc la cofinalité n'est pas totalement tautologique! Mais si on prouve que la catégorie $\prod_{j\in J} C_j$ est
\printedpage{23}
% head: A. Grothendieck : Dérivateurs, Ch. XVIII, Édition des Universités de Montpellier II et Paris VII
% foot: Version 17 décembre 2025
\msp{page 94}
grande devant $\pi$, donc filtrante, cela implique qu'elle est 0-connexe, et on peut appliquer ce même résultat à $C' = \prod_{j\in J'} C_j$, et on gagne.

Donc on est ramené à prouver que $\prod C_j$ est grande devant $\pi$. Ceci se voit soit sur le critère $\mathrm{PF}_\pi\,1$, $\mathrm{PF}_\pi\,2$ de 4.4.4 b, soit en se ramenant au cas où les $C_j$ sont ordonnées, mais cela a l'air plus long. Notons que l'on prouve de même qu'un (petit) produit de catégories filtrantes est encore une catégorie filtrante.

\dgrblock{dgrb977}{Products of indexing categories and the first coproducts}
\begin{dcomment}{exposition}{p94:Notons que l'on prouve de même}{Products of indexing categories and the first coproducts}
The coproduct argument isolates the part of the construction that can be done before arbitrary colimits have been established. Put $\kappa=\pi^+$. A category called ``grande devant $\pi$'' here is $\kappa$-filtered. If $(D_j)_{j\in J}$ is any small family of such categories, a diagram with fewer than $\kappa$ arrows in $\prod_jD_j$ projects to such a diagram in every factor. Choose a cocone in each factor and take their product. This proves that the product is $\kappa$-filtered, including its nonemptiness; the usual axiom of choice is being used in making these simultaneous choices.

The projection to $D_{j_0}$ is final for a similarly concrete reason. For $d\in D_{j_0}$ its comma category is
\[
(d\downarrow\operatorname{pr}_{j_0})
\simeq(d\downarrow D_{j_0})\times\prod_{j\ne j_0}D_j.
\]
The first factor is filtered and nonempty, and the remaining product is filtered and nonempty by the argument just given. Their product is therefore connected. The finality criterion follows. This proof also explains the warning about products of connected categories. Connectedness alone supplies zigzags of possibly unbounded length; filteredness supplies compatible cocones of one fixed shape.

For a literal example of that distinction, let $P_n$ be the poset on $0,1,\ldots,n$, with each even vertex below its adjacent odd vertices and no further nonidentity comparisons. Every $P_n$ is connected. In $\prod_{n\geq1}P_n$, however, the tuples $(0)_n$ and $(n)_n$ cannot be joined by a finite zigzag. A zigzag of length $r$ projects to one of length $r$ in every factor, whereas the distance from $0$ to $n$ in the underlying undirected graph of $P_n$ is $n$. Choosing $n>r$ gives a contradiction. Thus the source's caution is a real one, even for a product of finite posets.

Now let $|J|\leq\pi$ and $F_j\in\Ind'_\pi C$, assuming only that $C$ has coproducts with at most $\pi$ summands. Set $I=\prod_j(C/F_j)$. For $i\in I$, write $c_{i,j}\to F_j$ for its $j$th component and put $s_i=\coprod_jc_{i,j}$ in $C$. The object
\[
H=\varinjlim_{i\in I}y(s_i)
\]
belongs to $\Ind'_\pi C$, because $I$ is $\kappa$-filtered. Here $y$ denotes the Yoneda embedding, and this colimit is taken in the presheaf category. To identify $H$, take $G\in\Ind'_\pi C$. Such a $G$ sends the given coproducts in $C$ to products of sets: it is a $\kappa$-filtered colimit of representables, and these colimits commute with products having fewer than $\kappa$ factors. Hence
\[
\begin{aligned}
\Hom(H,G)
&\simeq\varprojlim_iG(s_i)
\simeq\varprojlim_i\prod_jG(c_{i,j})\\
&\simeq\prod_j\varprojlim_iG(c_{i,j})
\simeq\prod_j\Hom(F_j,G).
\end{aligned}
\]
The last isomorphism uses the finality of the projections just proved. All the maps are induced by the evident coproduct cocone, so this proves the required universal property. For an empty family, $I$ is terminal and $s_i$ is the initial object of $C$; the same calculation includes the empty coproduct. There is therefore no circular use of the cocompleteness theorem in Lemma 4.4.11.
\end{dcomment}


\textbf{Remarques.}
\begin{enumerate}[label={}, leftmargin=*]
\item[1.] Cet argument ne montre nullement (même si $C$ est stable par $\varinjlim$ de cardinal $\leqslant \pi$) que pour tout $\varphi \in \uHom(J,\Ind(C))$, le foncteur similaire à (4.4.6.3)
\[
\uHom(J,C)/\varphi \longrightarrow C/\varphi(j)
\]
soit cofinal, et je doute qu'il le soit toujours. La difficulté provient du fait que $C \longrightarrow \Ind\, C$ ne commute pas aux sommes de cardinal $\leqslant \pi$, seulement aux sommes finies,
\msp{page 95}
donc le lemme 4.4.10 ne peut être invoqué.
\item[2.] Nous venons de prouver, via 4.4.8 et 4.4.9 (dont la démonstration est achevée avec 4.4.12), que si $C$ est une petite catégorie stable par $\varinjlim$ de cardinal $\leqslant \pi$, alors $\Ind_\pi C$ est également stable par ces limites. De plus, les systèmes inductifs de ce type $J$ dans $\Ind_\pi C$ admettent une représentation indicielle
\[
\psi \ :\ I \times J \longrightarrow C\ ,
\]
avec $I$ grande devant $\pi$. Et on peut toujours supposer, quitte à prendre $I' \longrightarrow I$ cofinal avec $I'$ ensemble ordonné grand devant $\pi$, que $I$ est de plus un ensemble ordonné.
\end{enumerate}

\dgrblock{dgrb978}{An ordinary ind-object whose action has no stage}
\begin{dcomment}{examples}{p94:et je doute qu'il le soit toujours}{An ordinary ind-object whose action has no stage}
The doubt on page 94 concerns the loss of a cardinal bound when $\Ind_\pi(C)$ is replaced by the ordinary ind-completion. It has a concrete negative answer. Fix a small skeleton $C$ of the category of at most countable sets. This category has all countable colimits, so it satisfies the hypothesis under discussion with $\pi=\aleph_0$. In $\widehat C$, define
\[
F(c)=\{f:c\longrightarrow\mathbb Z\mid f(c)\text{ is finite}\}.
\]
The presheaf $F$ is an ordinary ind-object. Indeed, the finite subsets $T\subset\mathbb Z$ form a filtered poset, and
\[
F=\varinjlim_{T\subset\mathbb Z,\ T\text{ finite}}y(T)
\]
pointwise: a map lies in the displayed colimit exactly when its image is contained in one finite $T$. Translation on $\mathbb Z$ defines an action on $F$, hence a diagram $\varphi:B\mathbb Z\to\Ind'(C)$.

An object of $C^{B\mathbb Z}/\varphi$ consists of a countable $\mathbb Z$-set $c$ and an equivariant map of presheaves $y(c)\to F$. By Yoneda, the map corresponds to a function $f:c\to\mathbb Z$ with finite image. Equivariance says
\[
f(n\cdot x)=n+f(x)\qquad(n\in\mathbb Z).
\]
If $c$ were nonempty, the image of $f$ would therefore be a nonempty finite subset of $\mathbb Z$ invariant under every translation. Taking any element in that image produces all its translates, a contradiction. Thus every such $c$ is empty.

Evaluation at the unique object of $B\mathbb Z$ consequently has image only the initial presentation $y(\varnothing)\to F$. In $C/F$, the element $0\in F(1)$ defines a presentation $y(1)\to F$, and there is no morphism from this presentation to $y(\varnothing)\to F$. The corresponding comma category for evaluation is empty. Evaluation is therefore not final, exactly as Grothendieck suspected.

The example also locates the obstruction. At $c=\mathbb N$, the identity-shaped enumeration of an infinite set cannot be represented at a finite stage. Accordingly, $F$ does not send a countable coproduct of singletons to the corresponding product: its value there consists of finite-image sequences. It is not an object of $\Ind'_{\aleph_0}C$. The stronger filteredness in the positive theorem is what prevents this failure of simultaneous realization. The distinction between a formal ind-object and the colimit of its values already appeared in the commentary on Chapter~VII, manuscript p.~123 (``A colimit need not be the same ind-object''); here it obstructs a specific comparison of diagram categories.
\end{dcomment}


\dgrblock{dgrb979}{The modern comparison for diagram categories}
\begin{dcomment}{later}{p94:et je doute qu'il le soit toujours}{The modern comparison for diagram categories}
Grothendieck's simultaneous presentation is the essential-surjectivity part of a familiar modern comparison. With $\kappa=\pi^+$ and $J$ having fewer than $\kappa$ arrows, it gives
\[
\operatorname{Ind}_{\kappa}(C^J)
\longrightarrow
\operatorname{Ind}_{\kappa}(C)^J
\]
when $C$ has $\kappa$-small colimits. Here $\operatorname{Ind}_{\kappa}$ uses the modern regular-cardinal convention, so it is Grothendieck's $\Ind_\pi$, not his ordinary $\Ind$.

\citet[Theorem 1.2 and Corollary 2.2]{Henry2025} proves that this comparison is an equivalence under precisely these hypotheses, and characterizes the diagram shapes for which the statement holds uniformly in $C$. The paper also shows why removing the colimit hypothesis on $C$ requires additional restrictions on $J$; its Theorem 1.3 treats the corresponding well-founded shapes. These are results about the chosen presentability rank and about realizing diagrams by stages. They should be distinguished from the mere existence of some accessibility rank for a functor category. We use the theorem numbering of the checked arXiv version 2, dated 9 October 2024; the published article appeared in 2025 under the title cited in the bibliography.
\end{dcomment}


C'est là le moyen donc, pour la démonstration du

\textbf{Théorème 4.4.13.} \textit{Soit $C$ une catégorie équivalente à une petite catégorie, et stable par $\varinjlim$ de cardinal $\leqslant \pi$, où $\pi$ est un cardinal infini donné. Alors :}
\msp{page 96}

\begin{enumerate}[label={}, leftmargin=*]
\item[\textit{a)}] \textit{(Pour mémoire.) Le foncteur canonique}
\[
C \longrightarrow \Ind_\pi(C) \xrightarrow{\ \sim\ } \Ind'_\pi(C)
\]
\textit{commute aux $\varinjlim$ de cardinal $\leqslant \pi$.}
\item[\textit{b)}] \textit{(Pour mémoire.) Les systèmes inductifs de cardinal $\leqslant \pi$ dans $\Ind_\pi(C)$ (i.e. les foncteurs $\varphi : J \longrightarrow \Ind_\pi(C)$, avec $\card \Fl(J) \leqslant \pi$) admettent une représentation indicielle}
\[
\psi \ :\ I \times J \longrightarrow C\ ,
\]
\textit{avec $I$ ensemble ordonné grand devant $\pi$, et admettent une limite inductive dans $\Ind_\pi C$ qui se calculent par}
\[
\underbrace{\varinjlim_{j\in J} \varphi(j)}_{\text{dans } \Ind_\pi(C)} \ \simeq\ \text{“}\,\varinjlim_{i\in I}\,\text{”}\Psi(i)
\]
\end{enumerate}

% ---- from der-ch18_p026-030.tex
\printedpage{24}
% head: A. Grothendieck : Dérivateurs, Ch. XVIII, Édition des Universités de Montpellier II et Paris VII
% foot: Version 17 décembre 2025

\textit{avec}
\[
\Psi(i) \;=\; \varinjlim_{j\in J}{}_{C}\,\psi(i,j)\;.
\]
\textit{c)} $\Ind_\pi C$ \textit{est stable par petites limites inductives} (\footnote{rajouter aussi, en d), le corollaire 4.4.16, page 101.}).

Il reste à établir ce dernier point. Comme on a déjà les $\varinjlim$ finies par b, il reste à prouver l'existence des petites $\varinjlim$ filtrantes, ou au choix des sommes infinies, ce qui sera plus joli en fait. Mais
\msp{page 97}
on pourra aussi bien traiter le cas général d'une $\varinjlim$ de type $J$ quelconque ($J$ petite). Pour ceci, soit
\[
I \;=\; I(J) \;=\; \begin{aligned}[t]
&\text{ensemble des sous-catégories } J' \text{ de } J \text{ telles que}\\
&\card\Fl(J')\leqslant\pi,
\end{aligned}
\]
ordonnons $I$ par inclusion. Il est évident que $I\neq\emptyset$ (car $\emptyset\in I$), et même que $I$ \textit{est grand devant} $\pi$, car si $(J'_\alpha)_{\alpha\in A}$ avec $\card(A)\leqslant\pi$ est une famille d'éléments de $I$, alors la sous-catégorie de $J$ engendrée par les $J'_\alpha$ est encore dans $I$ (vérification immédiate), et majore les $J'_\alpha$. En fait, la considération de $I(J)$ associée à $J$ va nous permettre de prouver ceci, qui généralise 4.4.13 c), compte tenu de b) :

\textbf{Corollaire 4.4.14} (\footnote{Comparer aussi 4.11.3, page 270.}). \textit{Soit $\mathcal{M}$ une catégorie, $\pi$ un cardinal infini tel que}
\begin{enumerate}
\item[a)] \textit{$\mathcal{M}$ est}
\msp{page 98}
\textit{stable par $\varinjlim$ de cardinal $\leqslant\pi$, et}
\item[b)] \textit{$\mathcal{M}$ est stable par $\varinjlim$ sur les ensembles ordonnés $I$ grands devant $\pi$.}
\end{enumerate}
\textit{Alors $\mathcal{M}$ est stable par petites $\varinjlim$ .}

(\footnote{\textbf{Corollaire 4.4.14.1.} $\mathcal{M}$ étant comme dans 4.4.14, et $f:\mathcal{M}\longrightarrow\mathcal{E}$ étant un foncteur qui commute aux $\varinjlim$ de cardinal $\leqslant\pi$, et aux petites $\varinjlim$ grandes devant $\pi$, alors $f$ commute aux petites $\varinjlim$ .}).

En effet, si $\varphi:J\longrightarrow\mathcal{M}$ est donné, on en construit
\[
\psi\;:\;I\;=\;I(J)\;\longrightarrow\;\mathcal{M}
\]
par
\[
\psi(J')\;=\;\varinjlim_{J'}\varphi(J')\;,
\]
qui existe par hypothèse sur $\mathcal{M}$. Il est alors immédiat que $\varinjlim\limits_{J}\varphi$ existe si et seulement si $\varinjlim\limits_{I}\psi$ existe, or cette existence est assurée du fait que $I$ est grand devant $\pi$, et l'hypothèse sur $\mathcal{M}$.

\dgrblock{dgrb980}{Why bounded colimits suffice for cocompleteness}
\begin{dcomment}{exposition}{p98:qui existe par hypothèse sur}{Why bounded colimits suffice for cocompleteness}
The passage from bounded colimits to arbitrary small colimits is the structural conclusion of this section. Its proof uses a directed family of portions of the indexing category, so it concerns the diagram and its relations together. A family of subsets of objects would not retain the equations imposed by its arrows.

For a small category $J$, let $\mathcal I_\pi(J)$ be the poset of its subcategories with at most $\pi$ arrows, with the empty subcategory included. A union of at most $\pi$ such arrow sets still has size at most $\pi$. Closing it under identities and finite compositions does not increase this bound, because $\pi$ is infinite. The resulting subcategory is an upper bound. Thus $\mathcal I_\pi(J)$ is $\pi^+$-directed.

Given $X:J\to\mathcal M$, suppose that the colimit $X_K=\varinjlim_K X$ exists for every $K\in\mathcal I_\pi(J)$. These objects form a diagram indexed by $\mathcal I_\pi(J)$. We verify directly that its colimit is the colimit of $X$. A compatible family of maps $X_K\to Z$ gives a map $X(j)\to Z$ by using the subcategory containing only the identity of $j$. For an arrow $a:j\to j'$, choose a subcategory containing $a$ and its two objects. Its cocone forces the required equality between the two maps to $Z$. Conversely, a cocone on $X$ restricts to every $K$, hence factors uniquely through every $X_K$, and the resulting maps are compatible with inclusions. These two constructions are inverse and natural in $Z$. Therefore
\[
\varinjlim_J X\simeq
\varinjlim_{K\in\mathcal I_\pi(J)}\left(\varinjlim_KX\right)
\]
whenever the right side exists. This proves Corollary 4.4.14 without an interchange of limits whose existence has not yet been established.

The same cocone argument proves the preservation statement in footnote 27. If a functor preserves the inner colimits and the outer directed colimit, it preserves the colimit of $X$. No hypothesis about exactness, generators or representability is used. This makes the result useful beyond the particular ind-completion being constructed.

There is a size distinction that becomes consequential on page 110. If $|\operatorname{Fl}J|\leq\rho$, then the set $\mathcal I_\pi(J)$ can have size $\rho^\pi$, rather than $\rho$. Thus the same decomposition proves a bound-preserving assertion only if this extra exponent has been controlled, or if a different argument supplies a smaller final indexing family. The unbounded cocompleteness theorem is sound; it does not by itself establish the later assertion that a pointwise cardinal bound survives every bounded colimit.
\end{dcomment}


Pour achever de prouver le théorème 4.4.13, on est donc ramené au cas particulier des limites inductives $\varinjlim\limits_{J}\varphi(j)$, quand $\varphi:J\longrightarrow\Ind'_\pi C$ est donné, avec $J$ ensemble
\printedpage{25}
% head: A. Grothendieck : Dérivateurs, Ch. XVIII, Édition des Universités de Montpellier II et Paris VII
% foot: Version 17 décembre 2025
\msp{page 99}
ordonné grand devant $\pi$. Mais ceci va résulter du résultat plus précis :

\textbf{Corollaire 4.4.15} (\footnote{comparer à 4.7.3.1, qu'il y aurait lieu d'expliciter ici.}). \textit{Soit $C$ une petite catégorie. Alors la sous-catégorie strictement pleine} $\Ind'_\pi C$ \textit{de $C^\wedge$ est stable dans $C^\wedge$ par $\varinjlim$ grandes devant $\pi$, en d'autres termes, ces $\varinjlim$} existent dans $\Ind'_\pi(C)$, \textit{et le foncteur d'inclusion y commute.} % doubt resolved: existent is upright in the PDF; see corrections.csv, PDF leaf 27

\textsc{Démonstration}. Soit $I$ ensemble ordonné, grand devant $\pi$, et $i\longmapsto F_i$ un foncteur $I\longrightarrow\Ind'_\pi(C)$, $F\;=\;\varinjlim\limits_{I}F_i$ dans $C^\wedge$, prouvons que $F\in\Ob\Ind'_\pi(C)$, i.e. (4.4.3) que $C/F$ est grande devant $\pi$. J'utilise le critère b) de 4.4.4, il faut vérifier $\PF_\pi\,1$, $\PF_\pi\,2$. Soient
\[
x_\alpha\xrightarrow{\;u_\alpha\;}F\;=\;\varinjlim F_i
\]
des objets de $C/F$ indexés par $A$ avec $\card A\leqslant\pi$. Chaque $u_\alpha$ provient de $x_\alpha\longrightarrow F_{i_\alpha}$, et en prenant un majorant commun des $i_\alpha$ ($I$ grand devant $\pi$), on trouve un $i\in I$ et des $v_\alpha:x_\alpha\longrightarrow F_i$. Comme $F_i$ est dans $\Ind'_\pi(C)$, i.e. $C/F_i$ grande devant $\pi$, considérant les $v_\alpha$ comme des objets de $C/F_i$, ils sont majorés par un objet $x\longrightarrow F_i$ de $C/F_i$. On a donc
\[
\begin{tikzcd}
x_\alpha \arrow[r] \arrow[d] & F_{i_\alpha} \arrow[d] \arrow[dr, dotted] & \\
x \arrow[r] & F_i \arrow[r] & F\;,
\end{tikzcd}
\]
qui montre que les $x_\alpha\longrightarrow F$ sont majorés par $x\longrightarrow F$,
\msp{page 100}
ce qui prouve $\PF_\pi\,1$. Prouvons $\PF_\pi\,2$, i.e. soient $x\xrightarrow{\;u\;}F$, $y\xrightarrow{\;v\;}F$ deux objets de $C/F$, et $(f_\alpha)_{\alpha\in A}$ une famille de cardinal $\leqslant\pi$ de flèches de $u$ dans $v$. On sait que $u$, $v$ proviennent de $u_i:x\longrightarrow F_i$, $v_i:y\longrightarrow F_i$ pour $i$ assez grand.
\[
\begin{tikzcd}[column sep=small]
x \arrow[rr, "f_\alpha"] \arrow[dr, "u"'] & & y \arrow[dl, "v"] \\
& F &
\end{tikzcd}
\]
\[
\begin{tikzcd}[column sep=small]
x \arrow[rr, "f_\alpha"] \arrow[dr, "u_i"'] & & y \arrow[dl, "v_i"] \\
& F_i &
\end{tikzcd}
\]

Pour tout $\alpha$, on peut choisir $i_\alpha\geqslant i$ tel que $v_if_\alpha$ et $u_i$ deviennent égaux en composant avec $F_i\longrightarrow F_{i_\alpha}$, et comme $\card A\leqslant\pi$ et $I$ grand devant $\pi$, il existe un majorant $i'$ des $i_\alpha$. Au total, on peut supposer (remplaçant $i$ par $i'$) que $v_if_\alpha=u_i$ pour tout $\alpha$. Alors $u_i$, $v_i$ étant considérés comme objets de $C/F_i$, les $f_\alpha$ sont des flèches de $C/F_i$ de $u_i$ dans $v_i$. Comme $C/F_i$ est grand devant $\pi$ (par l'hypothèse $F_i\in\Ob\Ind'_\pi(C)$), on trouve qu'il existe un égalisateur $(y,v_i)\xrightarrow{\;g\;}(z,w_i)$.
\printedpage{26}
% head: A. Grothendieck : Dérivateurs, Ch. XVIII, Édition des Universités de Montpellier II et Paris VII
% foot: Version 17 décembre 2025
\[
\begin{tikzcd}
x \arrow[r, shift left=2, "(f_\alpha)"] \arrow[r] \arrow[r, shift right=2] \arrow[dr, "u_i"'] & y \arrow[r, "g"] \arrow[d, "v_i"'] & z \arrow[dl, "w_i"] \\
& F_i &
\end{tikzcd}
\]

Alors l'image de $g$ dans $C/F$ (via $F_i\longrightarrow F$) égalise les flèches $f_\alpha$, en tant que flèches de $C/F$, ce qui prouve $\PF_\pi\,2$ et achève la démonstration de 4.4.14, et par là aussi celle du théorème 4.4.13, et du lemme 4.3. Par là est aussi achevée la démonstration de la première assertion du lemme
\msp{page 101}
4.3.7 (p. 77), ce qui achève la démonstration du théorème 4.3.6 (p. 70-71). Il reste à établir la propriété universelle énoncée dans 4.3.7, que je vais rappeler :

\dgrblock{dgrb981}{Filtered closure and its exact strength}
\begin{dcomment}{exposition}{p100}{Filtered closure and its exact strength}
Grothendieck now proves the closure property that footnote 24 had correctly distinguished from the definition of an ind-object. A single $\kappa$-filtered colimit of representables belongs to $\operatorname{Ind}_\kappa C$ by definition. The closure assertion says that a $\kappa$-filtered colimit of objects already in this category belongs to it again. It is this second assertion that allows arbitrary cocompleteness to be deduced from the preceding decomposition.

Here is the argument in a form that displays the two uses of the same cardinal. Let $\kappa=\pi^+$, let $I$ be $\kappa$-filtered, and let $F=\varinjlim_iF_i$ in $\widehat C$, with every category of elements $C/F_i$ $\kappa$-filtered. Replace $I$ by a final $\kappa$-directed poset if desired. Given fewer than $\kappa$ objects of $C/F$, choose representatives of their defining elements at stages of $I$. A common upper bound places all of them in one $C/F_i$. Filteredness of that category supplies a common target, whose image is the desired target in $C/F$.

For the second filteredness condition, start with fewer than $\kappa$ parallel arrows between two objects of $C/F$. Place the two defining elements in one $F_i$. Each equality saying that an arrow is over $F$ becomes true at some later stage, because equality in a filtered colimit of sets is eventual equality. A common upper bound makes all these equations true together. The corresponding arrows now lie in $C/F_{i'}$, where they can be simultaneously coequalized. Their images coequalize the original arrows. This proves both filteredness conditions and hence $F\in\operatorname{Ind}_\kappa C$.

The word ``\emph{égalisateur}'' on page 100 denotes a morphism that makes the given arrows equal after composition. No universal coequalizer is asserted or needed. Similarly, no colimits in $C$ are used in this closure proof. The hypothesis that $C$ has $\kappa$-small colimits belongs to the cocompleteness theorem, not to Corollary 4.4.15 itself.

The inclusion into $\widehat C$ creates these $\kappa$-filtered colimits, but generally does not preserve all ordinary filtered colimits. For instance, let $C$ again consist of countable sets and take $\pi=\aleph_0$. The completion $\Ind'_\pi C$ identifies with sets via $X\mapsto(c\mapsto X^c)$: every set is the $\aleph_1$-filtered union of its countable subsets, and a map out of a countable set factors through one such subset. The increasing union of finite subsets of $\mathbb N$ has colimit $\mathbb N$ in this completion. Its colimit in $\widehat C$, however, contains only finite-image maps $c\to\mathbb N$. At $c=\mathbb N$ it omits the identity map. This is the precise failure used in the preceding example and the reason the proof of Corollary 4.4.16 must retain the larger filteredness threshold.
\end{dcomment}


\textbf{Corollaire 4.4.16} (\footnote{Pourrait être incorporé en partie d) au théorème 4.4.13, car c'est un complément important, et pas tellement évident finalement.}). \textit{Soient $C$ une catégorie équivalente à une petite catégorie, stable par $\varinjlim$ de cardinal $\leqslant\pi$, où $\pi$ est un cardinal infini donné, et considérons le foncteur d'inclusion canonique}
\[
C\xrightarrow{\;\alpha\;}\Ind_\pi(C)\;. \tag*{(4.4.16.1)}
\]
\textit{Ce foncteur est} 2\textit{-universel pour les foncteurs $f:C\longrightarrow\mathcal{E}$, où $\mathcal{E}$ est une $\mathcal{U}$-catégorie stable par petites $\varinjlim$ , et $f$ un foncteur commutant à ces limites inductives. En d'autres termes, on a}
\begin{enumerate}
\item[\textit{1$^\circ$)}] \textit{$\alpha$ possède les propriétés précédentes, i.e.} $\Ind_\pi(C)$ \textit{est stable par $\varinjlim$ , et le foncteur $\alpha$ commute aux $\varinjlim$ de cardinal $\leqslant\pi$.}
\item[\textit{2$^\circ$)}] \textit{Pour toute $\mathfrak{U}$-catégorie $\mathcal{E}$ stable par petites limites inductives, le foncteur}
\[
\begin{aligned}
g\;&\longmapsto\;g\circ\alpha\\
\uHom_{!}(\Ind_\pi C,\mathcal{E})\;&\longrightarrow\;\uHom_{!\,\pi}(C,\mathcal{E})
\end{aligned}
\tag*{(4.4.16.2)}
\]
\textit{est une équivalence de catégories, où l'indice $!$ (resp. $!\,\pi$) indique la sous-catégorie strictement pleine de la catégorie} $\uHom$
\msp{page 102}
\textit{envisagée, formée des foncteurs qui commutent aux petites $\varinjlim$ (resp. à celles de cardinal $\leqslant\pi$).}
\end{enumerate}

On peut donc dire aussi que $\Ind_\pi(C)$ joue le rôle d'une “enveloppe” stable par petites $\varinjlim$ de la catégorie $C$, quand les $\varinjlim$ de cardinal $\leqslant\pi$, lesquelles existent déjà dans $C$, sont respectées par les plongement de $C$ dans son enveloppe. C'est un résultat nullement tautologique qu'il suffit pour cela d'ajouter à $C$ des $\varinjlim$ \textit{grandes devant} $\pi$ d'objets de $C$, et que plus précisément, la catégorie enveloppe est réalisée par $\Ind_\pi C$ (ou, au choix et moins pléthoriquement, par $\Ind'_\pi C\subset C^\wedge$). % [sic] les plongement -> le plongement, number agreement

\dgrblock{dgrb982}{Density determines the continuous extensions}
\begin{dcomment}{exposition}{p102:C'est un résultat nullement tautologique}{Density determines the continuous extensions}
The universal property is the main result of the construction: the passage from $C$ to $\Ind_\pi C$ freely supplies all small colimits while preserving the colimits of size at most $\pi$ that were already present. There is a local error in the proof that should be separated from this correct theorem. The unrestricted restriction functor displayed across pages 102--103 is not fully faithful. What is fully faithful is its restriction to functors preserving the canonical presentations by objects of $C$.

For the required version, write $\mathcal M=\Ind'_\pi C$ and use a small skeleton of $C$. Each $F\in\mathcal M$ has its canonical presentation
\[
F\simeq\varinjlim_{(c,u)\in C/F}y(c).
\]
If $g,h:\mathcal M\to\mathcal E$ preserve these colimits, a natural transformation $\theta:g y\to h y$ extends at $F$ by taking the colimit of its components along $C/F$. Naturality for $F\to F'$ follows by postcomposition $C/F\to C/F'$ and the defining cocones. The extension is unique because maps out of the displayed colimit are determined by their composites with its structure maps. On representables it restricts to $\theta$, since $C/y(c)$ has its terminal object $(c,\mathrm{id})$. This proves full faithfulness on the required functor categories. It is the precise density statement of SGA 4, Expos\'e I, Proposition 7.8 \citep{SGA4I1972}. Cocontinuous functors satisfy its hypothesis; so do functors preserving the relevant $\pi^+$-filtered presentations.

The failure without the preservation hypothesis is visible in sets. Let $C$ be a small skeleton of sets of cardinal at most an infinite $\pi$, so that $\Ind'_\pi C\simeq\mathrm{Set}$. Define a functor $Q:\mathrm{Set}\to\mathrm{Set}$ by taking $\mathcal P(X)$ and identifying all subsets of cardinal at most $\pi$ to a single distinguished element $*$. Direct image under a function is well-defined on this quotient: the image of a subset of cardinal at most $\pi$ still has cardinal at most $\pi$. For every $X\in C$, $Q(X)$ is a singleton. The identity transformation of $Q$ and the transformation that sends every value to $*$ therefore have identical restrictions to $C$. They are different on any $X$ with $|X|>\pi$, since the class of $X$ is then different from $*$. Restriction is not even faithful on arbitrary functors.

This example explains why strict generation cannot justify the unrestricted assertion. A generating family detects morphisms of objects; density expresses objects as particular colimits; a functor must respect those colimits for its values and transformations to be reconstructed from the generators. The source uses the correct restricted conclusion throughout the rest of the proof. Replacing the one unrestricted assertion by the density argument above repairs the proof without changing Corollary 4.4.16.
\end{dcomment}


\textsc{Démonstration}. Le 1°) est contenu dans 4.4.13 a) et c). Prouvons donc 2°). Il est évident que $\alpha$ est un foncteur pleinement fidèle dont l'image essentielle engendre strictement
\printedpage{27}
% head: A. Grothendieck : Dérivateurs, Ch. XVIII, Édition des Universités de Montpellier II et Paris VII
% foot: Version 17 décembre 2025
$\mathcal{M}=\Ind_\pi C$, d'où résulte que le foncteur (4.4.16.2) est pleinement fidèle, étant induit par le foncteur pleinement fidèle
\[
\alpha^*_{\mathcal{E}}\;:\;\uHom(\mathcal{M},\mathcal{E})\;\longrightarrow\;\uHom(C,\mathcal{E})
\]
\msp{page 103}
sur les sous-catégorie pleines envisagées. Il reste donc à montrer que (4.4.16.2) est essentiellement surjectif. Soit donc $f:C\longrightarrow\mathcal{E}$ dans $\uHom_{!\,\pi}(C,\mathcal{E})$, considérons son extension canonique % [sic] les sous-catégorie pleines -> les sous-catégories pleines, number agreement
\[
\left\{
\begin{aligned}
&Lf\;:\;\Ind(C)\;\longrightarrow\;\mathcal{E}\\
&Lf((X_i)_{i\in I})\;=\;\varinjlim_{i\in I}f(X_i)
\end{aligned}
\right.
\tag*{(4.4.16.3)}
\]
(cf. SGA 4 I 8.7.3), et soit
\[
\psi(f)\;=\;Lf|\Ind_\pi(C)\;. \tag*{(4.4.16.4)}
\]
Rappelons que $Lf$ commute aux petites $\varinjlim$ filtrantes, mais comme le foncteur d'inclusion
\[
\Ind_\pi C\;\longrightarrow\;\Ind(C) \tag*{(4.4.16.5)}
\]
ne commute pas à ces $\varinjlim$ , il n'est pas clair a priori que $\psi(f)$ commute auxdites limites, ni a fortiori qu'il commute aux petites $\varinjlim$ sans plus. Mais on a vu que (4.4.16.5) commute aux $\varinjlim$ \textit{grandes devant} $\pi$, ce qui est contenu dans le corollaire 4.4.15 (vu que les limites inductives \textit{filtrantes} dans $\Ind(C)$ sont les limites \textit{dans} $C^\wedge$) : dire que $\Ind'_\pi(C)\longrightarrow\Ind'(C)$ commute à un certain type de $\varinjlim$ filtrantes, revient à dire qu'il en est ainsi du composé avec $\Ind'_\pi(C)\hookrightarrow C^\wedge$, ce qui est établi par 4.4.15. De ceci résulte que $\psi(f)$ \textit{commute aux} $\varinjlim$ \textit{grandes devant} $\pi$. Pour établir que $g=\psi(f)$ commute
\msp{page 104}
aux petites $\varinjlim$ , il reste donc à voir qu'il commute aux $\varinjlim$ de cardinal $\leqslant\pi$, en vertu de 4.4.14.1 (marge page 98). Pour l'établir, on utilise le théorème 4.4.13 b, en représentant un foncteur $\varphi:J\longrightarrow\Ind_\pi(C)$ (où $\card\Fl J\leqslant\pi$) sous forme indicielle par
\[
\Phi\;:\;I\times J\;\longrightarrow\;C\;, \tag*{(4.4.16.6)}
\]
$I$ étant un ensemble ordonné grand devant $\pi$. On a donc
\[
\varphi(j)\;=\;\varinjlim_{i\in I}{}_{\Ind_\pi C}\,\Phi(i,j)\;,
\]
et comme $g=\psi(f)$ commute aux $\varinjlim$ de type $I$ par ce qu'on a vu, on aura
\[
g(\varphi(j))\;=\;\varinjlim_{i\in I}{}_{(\mathcal{E})}\;g(\Phi(i,j))\;,
\]
d'où
\[
\varinjlim_{j\in J}{}_{(\mathcal{E})}\,g(\varphi(j))\;=\;\varinjlim_{j\in J}\;\varinjlim_{i\in I}\;\underbrace{g(\Phi(i,j))}_{=\,f(\Phi(i,j))}\qquad\text{(limites dans $\mathcal{E}$)}\;. \tag*{$(*)$}
\]
\printedpage{28}
% head: A. Grothendieck : Dérivateurs, Ch. XVIII, Édition des Universités de Montpellier II et Paris VII
% foot: Version 17 décembre 2025

D'autre part, on a
\[
\begin{aligned}
\varinjlim_{j\in J}{}_{\Ind_\pi}\,\varphi(j)\;&\simeq\;\varinjlim_{j\in J}\;\varinjlim_{i\in I}\;\Phi(i,j)\\
&\simeq\;\varinjlim_{i\in I}\;\underbrace{\varinjlim_{j\in J}\;\Phi(i,j)}_{\substack{\text{peut être calculé}\\ \text{dans $C$, soit}\\ \Psi(i)\,\in\,\Ob C}}\;,
\end{aligned}
\]
donc
\[
\varinjlim_{j\in J}{}_{\Ind_\pi}\,\varphi(j)\simeq\varinjlim_{i\in I}\;\Psi(i)\;, \tag*{$(**)$}
\]
\msp{page 105}
d'où en appliquant $g$, qui commute aux $\varinjlim$ de type $I$, et se rappelant maintenant que $g|C\simeq f$,
\[
g(\,\varinjlim_{J}{}_{\Ind_\pi C}\,\varphi(j))\;\simeq\;\varinjlim_{i\in I}\;f(\Psi(i))\;. \tag*{$(*{*}*)$}
\]
On a d'autre part
\[
\Psi(i)\;\overset{\text{déf}}{=}\;\varinjlim_{j\in J}\;\Phi(i,j)\;,\qquad\qquad \varinjlim\;\text{ dans } C\;,
\]
et comme $f$ commute aux $\varinjlim$ de type $J$ (car $\card\Fl J\leqslant\pi$), on en conclut
\[
f(\Psi(i))\;\simeq\;\varinjlim_{j\in J}\;f(\Phi(i,j))\;,
\]
donc par $(*{*}*)$
\[
g(\,\varinjlim_{j\in J}\;\varphi(j))\;\simeq\;\varinjlim_{i\in I}\;\varinjlim_{j\in J}\;f(\Phi(i,j))\;. \tag*{$(*{*}{*}*)$}
\]
En composant avec $(*)$ et utilisant à nouveau la commutativité de $\varinjlim\limits_{I}$ à $\varinjlim\limits_{J}$ , on trouve donc
\[
g(\,\varinjlim\;\varphi(j))\;\overset{\sim}{\longleftarrow}\;\varinjlim\;g(\varphi(j))\;,
\]
ce qui prouve que $g$ commute aux limites inductives de type $J$, et achève de prouver que $g$ est \texttt{[dans]} $\uHom_{!}(\Ind_\pi(C),\mathcal{E})$. Comme $g|C\simeq f$, on en conclut que le foncteur (4.4.16.2) est essentiellement surjectif, ce qui achève de prouver 4.4.16. De plus, on a construit un foncteur quasi-inverse explicite, $\psi$,
\msp{page 106}
s'insérant dans le diagramme commutatif
\[
\begin{tikzcd}[column sep=large, row sep=large]
\uHom_{!\,\pi}(C,\mathcal{E}) \arrow[r, "\psi", "(\sim)"'] \arrow[d, hook, "\text{incl. pl. fid.}"] & \uHom_{!}(\Ind_\pi(C),\mathcal{E}) \arrow[r] \arrow[d, hook, dotted, "\text{pl. fid.}"] & \uHom_{!\,\pi}(\Ind_\pi(C),\mathcal{E}) \\
\uHom(C,\mathcal{E}) \arrow[r, "L", "\substack{\sim\\ \text{(cf. SGA 4 I 8.7.3)}}"'] & \uHom_{!\,\mathrm{filt}}(\Ind(C),\mathcal{E}) \arrow[ur, "R\ \text{(restriction)}"'] &
\end{tikzcd}
\qquad.
\tag*{(4.4.16.6)} % [sic] (4.4.16.6) -> (4.4.16.7)?, the label (4.4.16.6) is already used for Phi on the preceding page
\]

\dgrblock{dgrb983}{The extension and the codomain of restriction}
\begin{dcomment}{exposition}{p106:s'insérant dans le diagramme commutatif}{The extension and the codomain of restriction}
The explicit inverse constructed here makes the universal property usable. For a functor $f:C\to\mathcal E$ preserving colimits of cardinal at most $\pi$, put
\[
\overline f(F)=\varinjlim_{(c,u)\in C/F}f(c).
\]
This is the restriction of the usual extension $Lf$ on the ordinary ind-completion. The latter is characterized by preservation of filtered colimits, as in SGA 4, Expos\'e I, Proposition 8.7.3 \citep{SGA4I1972}. It follows from Corollary 4.4.15 that $\overline f$ preserves $\pi^+$-filtered colimits. The proof on pages 104--105 supplies the remaining preservation property rather than assuming it from the inclusion into $\Ind(C)$.

To see the mechanism, present a diagram $X:J\to\Ind_\pi C$, with $|\operatorname{Fl}J|\leq\pi$, by $c:I\times J\to C$, where $I$ is $\pi^+$-filtered. Put $s_i=\varinjlim_{j\in J}c(i,j)$ in $C$. Then all the indicated colimits exist, and the two comparisons are
\[
\begin{aligned}
\overline f\left(\varinjlim_jX_j\right)
&\simeq\varinjlim_i f(s_i)
\simeq\varinjlim_i\varinjlim_j f(c(i,j)),\\
\varinjlim_j\overline f(X_j)
&\simeq\varinjlim_j\varinjlim_i f(c(i,j)).
\end{aligned}
\]
Both are the colimit in $\mathcal E$ of the same diagram indexed by $I\times J$. The isomorphism between them is therefore the actual comparison induced by the colimit cocone, not merely an unspecified isomorphism between two objects. Corollary 4.4.14.1 then gives preservation of every small colimit. Natural transformations extend by the same construction, so this proves the asserted equivalence of categories, including its two-dimensional content.

There is a second, independent typing error in the final diagram. Its arrow $R$, on the whole bottom-middle category, need not land in $\uHom_{!\,\pi}(\Ind_\pi C,\mathcal E)$. A constant singleton functor to sets preserves all nonempty filtered colimits, but it does not preserve the initial object. It is therefore an object of the stated source of $R$ whose restriction is not in the stated target. One may repair the diagram by replacing that target with the full functor category, or by restricting the source of $R$ to the extensions of functors that preserve the specified colimits in $C$. The dotted downward arrow remains meaningful: it is the fully faithful functor $g\mapsto L(g|_C)$ supplied by the two universal properties. These qualifications preserve the explicit inverse $\psi$ while making every arrow of the diagram well-typed.
\end{dcomment}


% ---- from der-ch18_p031-035.tex
\setcounter{footnote}{29}

\printedpage{29}
% head: A. Grothendieck : Dérivateurs, Ch. XVIII, Édition des Universités de Montpellier II et Paris VII
% foot: Version 17 décembre 2025

\section*{4.5\quad Filtration cardinale canonique d'une catégorie $\mathfrak{U}$-paratopos}

\textbf{Définition 4.5.1.} On appelle \textit{catégorie $\mathfrak{U}$-paratopos} une catégorie $\mathcal{M}$ $\mathfrak{U}$-pseudo-topos et accessible, i.e. qui satisfait aux conditions équivalentes du théorème 4.3.6 (p. 70). On dit que c'est une \textit{catégorie $\pi$-$\mathfrak{U}$-paratopos}, si elle est $\pi$-accessible, i.e. si elle admet une petite sous-catégorie génératrice dont les objets soient $\pi$-accessible, ou ce qui revient au même, si % [sic] soient π-accessible -> soient π-accessibles, accord du pluriel
$\mathcal{M}$ est équivalente à une catégorie de la forme $\Ind'_{\pi}(\mathcal{M}_{\pi})$, où $\mathcal{M}_{\pi}$ est une petite catégorie (ou : une catégorie équivalente à une petite catégorie), et stable par $\varinjlim$ de cardinal $\leqslant \pi$. (Et on peut prendre pour $\mathcal{M}_{\pi}$ la sous-catégorie pleine de $\mathcal{M}$ formée des objets
\msp{page 107}
$\pi$-accessibles).

Dans la suite, nous allons, plus généralement, considérer une catégorie $C$ équivalente à une petite catégorie, stable par $\varinjlim$ de cardinal $\leqslant \pi_0$ (\footnote{Il n'y a pas lieu de le supposer tout de suite.}), et la catégorie
\[
\mathcal{M} \ \overset{\text{déf}}{=}\ \Ind'_{\pi_0}C \ \subset\ C^{\wedge}\ ,
\tag*{(4.5.1)}
\]
où $\pi_0$ est un cardinal infini donné. Je me propose d'introduire sur $\mathcal{M}$ une filtration cardinale canonique :

\textbf{Proposition 4.5.2} (\footnote{N'est prouvé (et n'est vrai) que si $\pi^{\pi_0} = \pi$, voir rectification 4.5.6, page 116.}). \textit{Soit $F \in \Ob\mathcal{M} \subset \Ob C^{\wedge}$, et soit $\pi$ un cardinal $\geqslant \pi_1 =$} $\sup(\pi_0, \card\Fl C_0)$\textit{, où $C_0$ est une petite sous-catégorie pleine de $C$ telle que $C_0 \overset{\sim}{\hookrightarrow} C$. Les conditions suivantes sont équivalentes :}
\begin{enumerate}[leftmargin=3em]
\item[\textit{a)}] $\card\Fl(C_0/F) \leqslant \pi$.
\item[\textit{a$'$)}] $\card\Ob\,(C_0/F) \leqslant \pi$.
\item[\textit{a$''$)}] \textit{Pour tout $x$ dans $C_0$,} $\card F(x) \leqslant \pi$.
\item[\textit{b)}] \textit{Il existe une petite catégorie filtrante $I$ (pas nécessairement ordonnée),} grande devant $\pi_0$\textit{, avec} $\card I \leqslant \pi$\textit{, et un système inductif $(x_i)_{i\in I}$ dans $C$, tel que l'on ait}
\[
F \ \simeq\ \varinjlim_{i\in I} x_i\ ,
\tag*{(4.5.2.1)}
\]
\textit{où la $\varinjlim$ est prise dans $\mathcal{M}$, ou ce qui revient au même (en vertu de 4.4.15, $I$ étant grande devant $\pi_0$) dans $C^{\wedge}$, i.e. $F$ est isomorphe au foncteur ind-représentable défini par l'ind-objet $\underline{x} = (x_i)_{i\in I}$ de $C$.}
\item[\textit{c)}] \textit{$F$ est $\pi$-accessible dans $\mathcal{M}$.}
\item[\textit{c$'$)}] \textit{$F$ est $\pi$-accessible dans $C^{\wedge}$.}
\end{enumerate}
\msp{page 108}

\textsc{Démonstration}. On a les implications $\pm$ tautologiques
\[
\begin{tikzcd}[row sep=small]
 & \mathrm{a}' \arrow[r, Rightarrow] & \mathrm{a}'' & \\
\mathrm{a} \arrow[ur, Rightarrow] \arrow[dr, Rightarrow] & & & \\
 & \mathrm{b} \arrow[r, Rightarrow] & \mathrm{c}' \arrow[r, Rightarrow] & \mathrm{c}\ ,
\end{tikzcd}
\tag*{$(*)$}
\]
\printedpage{30}
% head: A. Grothendieck : Dérivateurs, Ch. XVIII, Édition des Universités de Montpellier II et Paris VII
% foot: Version 17 décembre 2025
où a$'$ $\Longrightarrow$ a$''$ provient de l'inclusion canonique $F(x) = \Hom_{\mathcal{M}}(x,F) \hookrightarrow \Ob C/F$, a $\Longrightarrow$ b s'obtient en prenant $I = C_0/F$ et notant que $I$ est stable par limites inductives de cardinal $\leqslant \pi_0$ ($C$ l'étant), donc est grande devant $\pi_0$ (\footnote{On n'a pas besoin de cette stabilité, de toutes façons $C_0/F$ est grande devant $\pi_0$, en vertu de 4.4.3, p. 80.}). On a b $\Longrightarrow$ c$'$ en vertu du fait que la sous-catégorie strictement pleine $C^{\wedge}_{\pi}$ de $C^{\wedge}$ définie par
\[
\left\{
\begin{array}{lcl}
C^{\wedge}_{\pi} & = & \text{sous-catégorie strictement pleine de } C^{\wedge} \text{ formée des}\\
 & & F \in \mathcal{M} \text{ qui sont } \pi\text{-accessibles.}
\end{array}
\right.
\tag*{(4.5.2.2)}
\]
est stable par $\varinjlim$ de cardinal $\leqslant \pi$ (SGA 4 I 9.9). Enfin c$'$ $\Longrightarrow$ c provient du fait que le foncteur d'inclusion $\mathcal{M} \hookrightarrow C^{\wedge}$ est pleinement fidèle et $\pi_0$-accessible, i.e. commute aux limites inductives grandes devant $\pi_0$, et a fortiori à celles grandes devant $\pi$. Par les implications du diagramme ci-dessus, il suffit donc de prouver qu'on a
% marginal mark (handwritten squiggle) in the left margin of the edition beside the following display
\[
\mathrm{c} \overset{(\mbox{\footnotemark[33]})}{\Longrightarrow} a \qquad\quad \text{et} \qquad\quad \mathrm{a}'' \ \Longrightarrow\ \mathrm{a}\ .
\]
\footnotetext[33]{n'est prouvé (et n'est vrai) que si $\pi^{\pi_0} = \pi$.}
\setcounter{footnote}{33}

\dgrblock{dgrb984}{Two meanings of the size of a presheaf}
\begin{dcomment}{exposition}{p108:Par les implications du diagramme ci-dessus}{Two meanings of the size of a presheaf}
The new section asks whether the intrinsic presentability rank of an object can be read from the size of its canonical presentation. This would turn the generators used to construct a paratopos into an effective cardinal filtration. The potential discrepancy is already visible in Proposition 4.5.2: presentability in $\mathcal M$ and presentability in the ambient presheaf category test different colimits.

Replace $C$ by the small category $C_0$ in the statement, and put $\theta=\pi_0$, $\mathcal M=\operatorname{Ind}_{\theta^+}C$, and $\kappa=\pi^+$, with $\pi\geq\max(\theta,|\operatorname{Fl}C|)$. The following equivalences require no exponentiation hypothesis:
\[
\begin{gathered}
|\operatorname{Fl}(C/F)|\leq\pi
\quad\Longleftrightarrow\quad
|\operatorname{Ob}(C/F)|\leq\pi
\quad\Longleftrightarrow\quad
\forall c\in C, |F(c)|\leq\pi,\\
\Longleftrightarrow\quad
F\text{ has a }\theta^+\text{-filtered presentation with at most }\pi\text{ arrows},\\
\Longleftrightarrow\quad
F\text{ is }\kappa\text{-presentable in }\widehat C.
\end{gathered}
\]
We give the details, since they isolate the part of the proposition which survives unchanged.

The objects of $C/F$ are pairs $(c,a)$ with $a\in F(c)$, so their number is bounded by $\pi$ exactly when all the values are bounded by $\pi$. A morphism of $C/F$ is determined by a morphism $a:c\to d$ of $C$ and an element of $F(d)$; its source element is then forced. Consequently
\[
|\operatorname{Fl}(C/F)|
=\sum_{a:c\to d}|F(d)|\leq\pi\cdot\pi=\pi.
\]
The reverse bound follows from identity arrows. If these bounds hold, the category $C/F$ itself is the required $\theta^+$-filtered presentation. Conversely, a presentation with at most $\pi$ arrows has values bounded by $\pi$: every representable has that bound, and a colimit of at most $\pi$ sets of that size again has that size.

Such a presentation also makes $F$ $\kappa$-presentable in $\widehat C$. Indeed, a colimit of at most $\pi$ representables represents a limit of at most $\pi$ evaluation functors, and $\kappa$-filtered colimits of sets commute with those limits. For the converse, write any presheaf as the colimit of its partial presentations indexed by the subcategories of $C/F$ with at most $\pi$ arrows. The preceding small-subcategory argument makes this indexing poset $\kappa$-directed. Each partial presentation is pointwise bounded by $\pi$. If $F$ is $\kappa$-presentable, its identity factors through one of them, making $F$ a retract of a pointwise bounded presheaf. Every value of $F$ then has cardinal at most $\pi$. This proves the final equivalence.

The class needed in (4.5.2.2) is the class of all $\kappa$-presentable objects of $\widehat C$. The printed condition $F\in\mathcal M$ in that definition is too restrictive for the assertion of closure under colimits taken in $\widehat C$. For example, if $C$ is the terminal category, then $\mathcal M$ contains only the singleton presheaf, whereas its empty colimit in $\widehat C=\mathrm{Set}$ is empty. The proof is repaired by using the ambient class, as above.

Finally, ambient $\kappa$-presentability implies presentability in $\mathcal M$, since the inclusion creates $\kappa$-filtered colimits. The converse is the unsettled implication $c\Rightarrow a$ in the source. Colimits in $\mathcal M$ may enlarge the values of a presheaf; preservation of the colimits that happen to be computed pointwise does not control this enlargement. The condition $\pi^{\pi_0}=\pi$ introduced on page 116 supplies the missing uniform estimate. It is not needed for the five equivalences just proved.
\end{dcomment}


\msp{page 109}

\textit{Prouvons} a$''$ $\Longrightarrow$ a. Considérons l'application
\[
\Fl C/F \ \longrightarrow\ \Fl C
\tag*{$(*)$}
\]
induit par le foncteur étale $C' = C/F \longrightarrow C$, alors l'ensemble des flèches au-dessus d'une % [sic] induit -> induite, accord avec « l'application »
flèche $\alpha : x \longrightarrow y$ de $C$ est en correspondance biunivoque avec $\Ob C'_y = F(y)$, puisqu'une telle flèche $\alpha'$ est connue quand on connaît son but, qui est un objet quelconque de $F(y)$ ($C' \longrightarrow C$ étant étale, i.e. fibrant à fibres discrètes). Donc le fibre de $(*)$ en $\alpha$ est de % [sic] le fibre -> la fibre, genre
cardinal $= \card F(y) \leqslant \pi$, donc on a
\[
\card\Fl C/F \ \leqslant\ \pi \underbrace{\card\Fl C}_{\leqslant \pi_1} \ \leqslant\ \pi \underbrace{\ \pi_1\ }_{\leqslant \pi} \ \leqslant\ \pi^2 = \pi\ ,
\]
ce qui prouve a.

Prouvons enfin c $\Longrightarrow$ a. Pour ceci, nous allons utiliser la première inclusion $\mathcal{M}_{\pi} \subset \Filt^{\pi}\mathcal{M}$ de SGA 4 I 9.16, qui s'applique dès lors que nous prouvons le

% marginal mark (handwritten squiggle) in the left margin of the edition beside Corollaire 4.5.3
\textbf{Corollaire 4.5.3} (\footnote{Pas prouvé, la condition c) des filtrations cardinales peut-être pas satisfaite. Sans doute \textit{faux} en général.}). \textit{Sous les conditions de 4.5.2, soit, pour tout cardinal $\pi \geqslant \pi_1$,} $\Filt^{\pi}(\mathcal{M})$ \textit{la sous-catégorie strictement pleine de $\mathcal{M}$ formée des $F \in \Ob\mathcal{M}$ qui satisfont la condition a) de 4.5.2. Alors la famille des} $(\Filt^{\pi})_{\pi\geqslant\pi_1}$ \textit{est une}
\msp{page 110}
\textit{filtration cardinale de $\mathcal{M}$.}

La condition a) de loc. cit. 9.12 (les $\uFilt^{\pi}$ sont équivalentes à des petites catégories) est évidente. Prouvons b), i.e. la stabilité de $\uFilt^{\pi}$ par $\varinjlim$ filtrantes $F = \varinjlim F_i$ \textit{grandes devant} $\pi_0$, de cardinal $\leqslant \pi$. Comme $\mathcal{M} \hookrightarrow C^{\wedge}$ commute à ces limites, pour $x$ dans $C$
\[
F(x) \ =\ \varinjlim_{i} F_i(x)\ ,
\]
\printedpage{31}
% head: A. Grothendieck : Dérivateurs, Ch. XVIII, Édition des Universités de Montpellier II et Paris VII
% foot: Version 17 décembre 2025
d'où $\card F(x) \leqslant \card\Ob I \cdot \pi \leqslant \pi \cdot \pi = \pi$, ce qui (par l'équivalence a $\Longleftrightarrow$ a$''$) implique $F \in \Ob\Filt^{\pi}$. Prouvons la condition c) de loc. cit. Pour ceci, nous aurons besoin du fait, plus précis que le b), qu'on vient d'établir :

\textbf{Corollaire 4.5.4} (\footnote{Prouvé seulement si $\pi^{\pi_0} = \pi$\,! Sans doute \textit{faux} en général.}). \textit{Pour tout $\pi \geqslant \pi_1$,} $\Filt^{\pi}$ \textit{est stable par $\varinjlim$ (dans $\mathcal{M}$) de cardinal $\leqslant \pi$ (pas nécessairement filtrantes, a fortiori pas nécessairement grandes devant $\pi_0$).}

Il suffit, en vertu de ce qui est déjà prouvé de du lemme 4.4.14 (\footnote{Il faut une variante du lemme 4.4.14, pour les $\varinjlim$ de cardinal $\leqslant \pi'$, qui se prouve de la même façon.}), de montrer que $\Filt^{\pi}$ % [sic] de du lemme -> du lemme, mot répété
est stable par $\varinjlim$ \textit{de cardinal} $\leqslant \pi_0$. Si $\varphi : J \longrightarrow \Filt^{\pi}$ est un foncteur avec $\card J \leqslant \pi_0$, on le met sous forme indicielle
\[
\psi\ :\ I \times J \ \longrightarrow\ C\ ,
\]
grâce à 4.14.3.b, avec $I$ grande devant $\pi_0$.

\msp{page 111}

Je rappelle qu'on pouvait prendre
\[
I \ =\ \uHom(J,C)/\varphi\ ,
\]
où $\uHom(J,C) \hookrightarrow \uHom(J, \mathcal{M} = \Ind'_{\pi}C)$, quitte à ne pas exiger que $I$ soit ordonnée. On a envie de prouver que $\card\Fl I \leqslant \pi$. L'ennui, c'est qu'on arrive seulement à la majoration
\[
\card\Fl I \ \leqslant\ \pi^{\pi_0}\ .
\]
(Même si $J$ est discrète, on n'a pas mieux, puisqu'alors
\[
I \ \simeq\ \prod_{j\in J} C/\varphi(j)\ , \qquad \text{d'où} \quad \Fl I = \prod_{j\in J} \Fl C/\varphi(j)\ ,
\]
et si $\card\Fl C/\varphi(j) = \pi$ pour tout $j$, on trouve $\card\Fl I = \pi^{\pi_0}$ -- il n'y a pas moyen de faire mieux\,!)

Donc je n'arrive pas à prouver directement 4.5.4, ce qui m'aurait arrangé pour prouver la condition c) de loc. cit. Mais à vrai dire, on a fait ici le travail de la démonstration du théorème 4.3.6, partie b $\Longrightarrow$ c -- on y a déjà prouvé que la filtration qu'on vient de construire (par la condition a) de 4.5.2) est une filtration cardinale. La condition c) de SGA 4 I 9.12 est prouvée à la page 76. Mais je vais refaire
\msp{page 112}
la démonstration plus sérieusement.

\dgrblock{dgrb985}{Counting simultaneous presentations}
\begin{dcomment}{exposition}{p112:la démonstration plus sérieusement.}{Counting simultaneous presentations}
Grothendieck's obstruction on page 111 is not a failure of the existence theorem for colimits. It concerns the amount of data needed to describe a compatible family of presentations at one time. This distinction is what the initial argument for a cardinal filtration had left uncontrolled.

Suppose $|\operatorname{Fl}C|\leq\pi$, $|\operatorname{Fl}J|\leq\theta$, and every value $\varphi(j)(c)$ has cardinal at most $\pi$. An object of $C^J/\varphi$ consists of a choice of objects and arrows of $C$ for the diagram, together with at most $\theta$ elements defining a compatible transformation to $\varphi$. There are at most $\pi^\theta$ such choices. An arrow in this comma category adds at most $\theta$ morphisms of $C$, subject to equations; imposing equations cannot increase the bound. Hence
\[
|\operatorname{Fl}(C^J/\varphi)|\leq\pi^\theta.
\]
For a discrete $J$ of size exactly $\theta$, the category is $\prod_{j\in J}(C/\varphi(j))$. If each factor has $\pi$ arrows, its arrow set has size $\pi^\theta$. Thus the counting argument itself cannot improve this estimate to $\pi$. The exact equality in the source's parenthesis tacitly takes $|J|=\pi_0$; for a smaller discrete $J$, the exponent is $|J|$.

There is a useful positive conclusion that does not require closing the entire argument prematurely. The pointwise bounded objects form an essentially small class: choose standard sets of cardinal at most $\pi$ at the set of objects of $C$, and then choose all structure maps satisfying the functor equations. They are also closed under $\theta^+$-filtered colimits indexed by categories with at most $\pi$ arrows, because those colimits are computed pointwise. These prove the first two estimates used on page 110. They do not prove that every object can be assembled from bounded ones through a $\pi^+$-filtered diagram.

To obtain the third property by the method on page 112, one must take a prescribed set of at most $\pi$ arrows in $C/F$ and enlarge it to a $\theta^+$-filtered subcategory without exceeding $\pi$ arrows. That is the content of Lemma 4.5.5. It is a genuinely stronger demand than adjoining a cocone to a single small diagram. Closing under all the required cocones may require $\pi^\theta$ simultaneous choices. The original unqualified Corollaries 4.5.3 and 4.5.4 should consequently be read with their retraction notes, not used as completed lemmas while proving Proposition 4.5.2.
\end{dcomment}


Soit donc $F \in \Ob\mathcal{M}$, considérons
\[
J \ =\ C/F\ ,
\]
c'est une catégorie stable par $\varinjlim$ de cardinal $\leqslant \pi_0$, donc grande devant $\pi_0$, et on a
\[
F \ \simeq\ \varinjlim_{j\in J} x_j \qquad\quad (\,\varinjlim \ \text{dans } \mathcal{M})
\]
\printedpage{32}
% head: A. Grothendieck : Dérivateurs, Ch. XVIII, Édition des Universités de Montpellier II et Paris VII
% foot: Version 17 décembre 2025
en désignant par $j \longmapsto x_j$ le foncteur canonique $C/F \longrightarrow C$. Soit alors $I$ l'ensemble des sous-catégories $J'$ de $J$ qui sont \textit{de cardinal $\leqslant \pi$ et grandes devant $\pi_0$, $I$} ordonné par inclusion, et pour tout $J' \in I$, posons
\[
X_{J'} \ =\ \varinjlim_{j\in J'} x_j\ .
\]
Comme $J'$ est grande devant $\pi_0$, la limite inductive dans $\mathcal{M}$ est aussi une limite dans $C^{\wedge}$, d'où résulte aussitôt, comme $\card\Ob J' \leqslant \card\Fl J' \leqslant \pi$, que $X_{J'}$ est dans $\Filt^{\pi}$. Alors
\[
J' \ \longmapsto\ X_{J'} \ \in\ \Ob\Filt^{\pi}
\]
est un foncteur
\[
I \ \longrightarrow\ \Filt^{\pi}\ ,
\]
dont la limite inductive est celle des $x_j$, donc isomorphe à $F$. D'autre part, il est immédiat que $I$ est grand devant $\pi$, i.e. que si $(J'_{\alpha})_{\alpha\in A}$ est
\msp{page 113}
une famille d'éléments de $I$, celle-ci est majorée par un $J' \in I$. Ceci résulte du

% marginal mark (handwritten squiggle) in the left margin of the edition beside Lemme 4.5.5
\textbf{Lemme 4.5.5} (\footnote{N'est prouvé \textit{que} si $\pi^{\pi_0} = \pi$\,! Sans doute faux en général.}). \textit{Soient $\pi$, $\pi_0$, avec $\pi > \pi_0$, des cardinaux infinis, et soit $J$ une catégorie grande devant $\pi_0$. Alors pour toute partie $\Phi$ de} $\Fl J$\textit{, telle que} $\card\Phi \leqslant \pi$\textit{, il existe une sous-catégorie $J'$ de $J$} grande devant $\pi_0$\textit{, telle que} $\card J' \leqslant \pi$ (\footnote{C'est OK si on exige seulement $\card J' \leqslant \pi^{\pi_0}$.})\textit{, et} $\Fl J' \supset \Phi$.

Ce lemme étant admis, on a donc représenté un objet quelconque $F$ de $\mathcal{M}$ comme une $\varinjlim\limits_{i\in I} X_i$, avec $I$ ensemble ordonné grand devant $\pi$, les $X_i$ dans $\Filt^{\pi}$. Il reste à prouver que si $\pi' > \pi$ et $F \in \Filt^{\pi'}$, alors on peut prendre $I$ tel que $\card I \leqslant \pi'^{\pi}$. Mais $F \in \Filt^{\pi'}$ signifie que $J = C/F$ satisfait $\card\Fl J \leqslant \pi'$. Alors l'ensemble des parties de cardinal $\leqslant \pi$ de $\Fl J$ a un cardinal majoré par $\pi'^{\pi}$, a fortiori $\card I \leqslant \pi'^{\pi}$, ce qui achève de prouver que $(\Filt^{\pi})_{\pi\geqslant\pi_1}$ est une filtration cardinale sur $\mathcal{M}$.

\msp{page 114}

Pour prouver le lemme, on commence par saturer $\Phi$ par composition finie, on trouve une sous-catégorie $J'_0$ de $J$ telle que $\card\Fl J'_0 \leqslant \pi_0$ \ldots{} Mais je m'aperçois que pour prouver le lemme, j'ai encore besoin de $\pi^{\pi_0} = \pi$ -- une enveloppe de $J'$ grande devant $\pi_0$ va avoir un cardinal majoré (au mieux) par $\pi^{\pi_0}$.

Prenons p.ex.
\[
J \ =\ \mathfrak{P}(S) \qquad\quad \text{(ensemble ordonné)}
\]
avec $\card S = \pi > \pi_0$, donc $J$ est grand devant $\pi_0$. Considérons l'application
\[
S \ \longrightarrow\ J = \mathfrak{P}(S)\ , \qquad\quad s \longmapsto \{s\}\ ,
\]
et prenons pour $\Phi \subset \Fl J$ l'ensemble des $\mathrm{id}_{\{s\}}$. Donc on cherche une partie $J'$ de $J$, contenant $S$, et grande devant $\pi_0$ pour l'ordre induit par $J$. On voit tout de suite que la plus petite partie $J'$ faisant l'affaire est l'ensemble $\mathfrak{P}_{\pi_0}(S)$ des parties de $S$ de cardinal $\leqslant \pi_0$. Mais on trouve
\[
\card J' \cdot \pi_0! \ =\ \pi^{\pi_0}\ ,
\]
\printedpage{33}
% head: A. Grothendieck : Dérivateurs, Ch. XVIII, Édition des Universités de Montpellier II et Paris VII
% foot: Version 17 décembre 2025
\msp{page 115}
où
\[
\begin{aligned}
\pi_0! \ &=\ \card\Aut\,(E_0) \qquad\quad (E_0 \text{ ensemble de cardinal } \pi_0)\\
&=\ \pi_0^{\pi_0}\\
&=\ 2^{\pi_0}\ ,
\end{aligned}
\]
donc
\[
\card J' \cdot 2^{\pi_0} \ =\ \pi^{\pi_0}\ .
\tag*{$(*)$}
\]
D'autre part, il est immédiat que $\card J' \geqslant \card\mathfrak{P}(E_0) = 2^{\pi_0}$ (où $E_0$ est une partie de $S$ de cardinal $\pi_0$), donc le premier membre de $(*)$ est égal à $\card J'$, donc on trouve
\[
\card J' \ =\ \pi^{\pi_0}\ ,
\]
sorry\,!

Finalement, je me suis convaincu que la proposition 4.5.2, et les corollaires et lemmes précédents par lesquels j'ai essayé de l'établir, n'est vrai que pour ces cardinaux $\pi > \pi_0$ tels que
\[
\pi^{\pi_0} \ =\ \pi\ ,
\]
ce qui, moyennant l'hypothèse du continu généralisée, équivaut à
\[
\pi^{\pi_0} \ <\ 2^{\pi}\ (\ =\ \pi^{\pi}\ )
\]
(ou encore à $\pi^{\pi_0} \neq 2^{\pi}$ $(= \pi^{\pi})$). Je vais donner
\msp{page 116}
ici un énoncé rectifié.

\dgrblock{dgrb986}{A corrected obstruction to a small directed hull}
\begin{dcomment}{development}{p115}{A corrected obstruction to a small directed hull}
The purpose of the example on pages 114--115 is to show that the cardinal cost of making an indexing category sufficiently filtered cannot always be avoided. Its chosen ambient poset, however, makes the proposed counterexample fail. In $J=\mathcal P(S)$, the family consisting of all singletons together with $S$ has a greatest element. It is therefore directed at every cardinal threshold, has cardinal $|S|$, and contains every prescribed singleton. The family $[S]^{\leq\pi_0}$ is not the smallest directed family with that property. Its cardinality, correctly calculated in the source, is not a lower bound for every possible directed hull.

The required obstruction is a statement about cofinality. We give a direct repair of this passage. We checked the relevant arguments in SGA 4, Expos\'e I, sections 8--9, and the diagram-completion comparison in \citet{Henry2025}; neither supplies a correction to this particular example. The following elementary application of diagonalization is offered as a repair, without a claim of priority.

\textbf{Proposition (a singular-cardinal obstruction).}\ours{}
Let $\mu>\aleph_0$ be a singular cardinal of countable cofinality. The poset
\[
J=[\mu]^{\leq\aleph_0},
\]
ordered by inclusion, is $\aleph_1$-directed. There is no $\aleph_1$-directed subposet $J'\subseteq J$ of cardinal at most $\mu$ containing every singleton.

\textit{Proof.} A countable union of countable sets is countable, so $J$ is $\aleph_1$-directed. If $J'$ were as stated, then for every countable $A\subseteq\mu$, the family of singletons $\{\{a\}:a\in A\}$ would have an upper bound in $J'$. Thus $J'$ would be cofinal in $J$.

Choose an increasing sequence of infinite cardinals $\mu_n<\mu$ with supremum $\mu$. If $|J'|\leq\mu$, we can write $J'=\bigcup_{n<\omega}\mathcal H_n$ with $|\mathcal H_n|\leq\mu_n$. Each union $U_n=\bigcup\mathcal H_n$ has cardinal at most $\mu_n\cdot\aleph_0<\mu$. Choose $a_n\in\mu\setminus U_n$, and put $A=\{a_n:n<\omega\}$. For any $H\in J'$, some $n$ has $H\in\mathcal H_n$, and then $a_n\notin H$. Consequently no member of $J'$ contains $A$, contradicting cofinality. This proves the proposition.\hfill$\square$

Taking $\pi_0=\aleph_0$, $\pi=\mu$, and the prescribed arrows to be the identities of all singletons gives an actual counterexample to Lemma 4.5.5 as originally stated. One may take $\mu=\beth_\omega$, which also lies above the elementary size bounds that occur earlier in the section. The argument establishes failure in ZFC and does not assume the generalized continuum hypothesis.

For comparison, the corrected upper estimate can be proved constructively at the level of categories. Let $J$ be $\theta^+$-filtered, let $\mu>\theta$ satisfy $\mu^\theta=\mu$, and start with at most $\mu$ specified arrows. Close them under identities and finite compositions. At each successor stage, choose cocones in $J$ for all diagrams with at most $\theta$ arrows in the current subcategory, adjoin their vertices and legs, and again close under composition. There are at most $\mu^\theta=\mu$ such diagrams, since both their shapes and their values can be encoded by at most $\theta$ pieces of data from sets of size at most $\mu$. Each stage therefore has at most $\mu$ arrows.

Continue this construction for $\theta^+$ stages and take the union. Its cardinal remains at most $\mu$, since $\theta^+\leq\mu$. Every diagram with at most $\theta$ arrows in the union occurs at one stage, by regularity of $\theta^+$, and acquires a cocone at the next stage. The union is therefore $\theta^+$-filtered and gives the desired hull. Without a fixed-point hypothesis, the same construction with the bound $\nu=\mu^\theta$ works because $\nu^\theta=\nu$. This proves the enlarged bound recorded in footnote 38. The two arguments distinguish the valid upper estimate from the invalid minimality assertion and provide both the positive and negative parts needed at this point of the manuscript.
\end{dcomment}


\textbf{Proposition 4.5.6.}
\begin{enumerate}[leftmargin=3em]
\item[\circled{1\textsuperscript{o}}] \textit{Sous les hypothèses générales de 4.5.2 sur $C$, $\pi_0$, $\pi_1$, $\pi$, les conditions a, a$'$, a$''$, b, c$'$ de cette proposition sont équivalentes (moyennant $\pi \geqslant \pi_1$)} (\footnote{\textbf{N.B.} Pour cette équivalence, l'hypothèse que $C$ soit stable par $\varinjlim$ de cardinal $\leqslant \pi$ est inutile, et aussi pour que c$'$) implique c), et lui soit équivalente si $\pi = \pi^{\pi_0} \geqslant \pi_1$, cf. 4.7.1, p. 161.})\textit{, et impliquent la condition c ($F$ est $\pi$-accessible). Quand $\pi^{\pi_0} = \pi$ (notamment quand $\pi = 2^{c}$, avec $c \geqslant \pi_0$), alors toutes les conditions a à c$'$ de 4.5.2 sont équivalentes.}
\item[\circled{2\textsuperscript{o}}] \textit{Plus généralement et plus précisément, pour $\pi \geqslant \pi_1$ donné, considérons les conditions sur $\pi$ :}
\begin{enumerate}[leftmargin=2em]
\item[\textit{1)}] $\pi^{\pi_0} = \pi$.
\item[\textit{2)}] $\Filt^{\pi}(\mathcal{M})$ \textit{(définie dans 4.5.3) est stable par $\varinjlim$ de cardinal $\leqslant \pi$.}
\item[\textit{3)}] \textit{$\forall$ $F$ dans $\mathcal{M}$, on a $F \simeq \varinjlim\limits_{I} X_i$, avec $I$ ensemble ordonné grand devant $\pi$, les $X_i$ dans} $\Filt^{\pi}(\mathcal{M})$\textit{. De plus, si $\pi' > \pi$ et $F \in$} $\Filt^{\pi'}(\mathcal{M})$\textit{, on peut prendre $I$ de cardinal $\leqslant \pi'^{\pi}$.}
\item[\textit{4)}] \textit{Comme 3), sans la majoration de cardinaux, et en supposant seulement $I$ catégorie grande devant $\pi$ (pas nécessairement catégorie ordonné).} % [sic] catégorie ordonné -> catégorie ordonnée, accord
\item[\textit{5)}] \textit{La condition c) de 4.5.2 implique les autres conditions (elles sont équivalentes, par 1°), en d'autres termes}
\[
\mathcal{M}_{\pi} \ \subset\ \Filt^{\pi} \qquad\quad \textit{(donc } \mathcal{M}_{\pi} = \Filt^{\pi}\textit{, puisqu'on sait que } \Filt^{\pi} \subset \mathcal{M}_{\pi}\textit{)}\ ,
\]
\end{enumerate}
\end{enumerate}

% ---- from der-ch18_p036-040.tex
\setcounter{footnote}{39}

\printedpage{34}
% head: A. Grothendieck : Dérivateurs, Ch. XVIII, Édition des Universités de Montpellier II et Paris VII
% foot: Version 17 décembre 2025

\msp{page 117}

\dgrblock{dgrb987}{Two measurements of size}
\begin{dcomment}{exposition}{p116}{Two measurements of size}
Grothendieck is separating the size of an object from the size of the collection of all maps into it from his chosen test objects. This distinction is the reason for the correction of Proposition 4.5.2. We retain his infinite cardinals $\pi_0\leq\pi$, and write $\lambda=\pi_0^+$ and $\mu=\pi^+$. Thus ``$\pi$-accessible object'' means $\mu$-presentable object, and a category ``large relative to $\pi$'' is $\mu$-filtered. All assertions here concern ordinary categories.

Let $C_0$ be a small skeleton of $C$, with at most $\pi_1$ arrows, and put $\mathcal M=\Ind'_{\pi_0}(C)$. Its embedding in $C^\wedge$ identifies an object $X$ with the presheaf $c\mapsto\Hom(c,X)$. For $\pi\geq\pi_1$, the following three bounds are equivalent:
\[
 |\Fl(C_0/X)|\leq\pi,
 \qquad |\Ob(C_0/X)|\leq\pi,
 \qquad |\Hom(c,X)|\leq\pi\quad(c\in C_0).
\]
Indeed, an object of $C_0/X$ is a pair $(c,u)$, while an arrow over $a:c\to d$ is determined by its target $v:d\to X$, its source being $va$. Counting first objects and then arrows proves the equivalence. Such a presheaf is $\mu$-presentable in the presheaf category: its category-of-elements presentation uses at most $\pi$ representables. Conversely, evaluation at $c$ preserves presentability once $|\Hom_C(c,d)|\leq\pi$; the precise adjunction is explained in the next block. This proves the equivalence with condition c$'$ without any cardinal exponent assumption.

The distinction becomes concrete in sets. Take $C_0$ to represent sets of cardinality at most $\pi_0$. An object $X$ of $\Ens$ is $\mu$-presentable exactly when $|X|\leq\pi$. The forward implication follows by writing $X$ as the $\mu$-filtered union of its subsets of size at most $\pi$ and factoring $\mathrm{id}_X$ through one of them. The converse follows because a function from at most $\pi$ elements into a $\mu$-filtered colimit can be lifted at one stage, and equality of two such functions can be witnessed at one later stage. On the other hand,
\[
 X\in\Filt^\pi
 \quad\Longleftrightarrow\quad |X|^{\pi_0}\leq\pi.
\]
For $|X|\geq2$, the exponent on the right bounds all maps from the objects of $C_0$, and is itself attained by maps from a set of size $\pi_0$; the empty and singleton cases satisfy the same conclusion directly. Consequently equality $\Filt^\pi=(\Ens)_\pi$ is equivalent to $\pi^{\pi_0}=\pi$. This proves a genuine necessity result, independently of the unsuccessful power-set argument on manuscript page 114.

The corrected estimate
\[
 \Filt^\pi\subseteq\mathcal M_\pi
       \subseteq\Filt^{\pi^{\pi_0}}
\]
retains useful information at every cardinal. Its second inclusion follows by replacing $\pi$ by $\rho=\pi^{\pi_0}$: since $(\pi^{\pi_0})^{\pi_0}=\pi^{\pi_0}$, the corrected equivalence applies at $\rho$, and $\mathcal M_\pi\subseteq\mathcal M_\rho$. There is no assumption of the generalized continuum hypothesis in this argument.

There is one local repair in the proof on pages 117--118. A fixed $\Filt^\pi$ is not closed under arbitrarily large $\lambda$-filtered colimits: arbitrarily large unions of small sets already contradict this. What is needed there is only closure under retracts. It holds immediately, since $\Hom(c,X)$ is a retract of $\Hom(c,Y)$ whenever $X$ is a retract of $Y$. The implication 4)$\Rightarrow$5) is therefore correct, with this replacement of the printed justification.
\end{dcomment}


\begin{quote}
\textit{où $\mathcal{M}_\pi$ désigne la sous-catégorie strictement pleine de $\mathcal{M}$ formée des objets $\pi$-accessibles.}
\end{quote}

\begin{quote}
\textit{Ceci posé, on a les implications}
\[
\mathrm{1)} \ \Longrightarrow\ \mathrm{2)} \ \Longrightarrow\ \mathrm{3)} \ \Longrightarrow\ \mathrm{4)} \ \Longrightarrow\ \mathrm{5)} \ ,
\]
\textit{et dans le cas particulier où $\mathcal{M} = $} (Ens), \textit{$C$ = sous-catégorie pleine formée des ensembles de cardinal $\leqslant \pi_0$, on a aussi} 5) $\Longrightarrow$ 1), \textit{i.e. toutes ces conditions sont équivalentes.}
\end{quote}

\textsc{Démonstration.}

1) $\Longrightarrow$ 2). Cf. démonstration de 4.5.4.

2) $\Longrightarrow$ 3). La catégorie $J = \Filt^\pi/F$ est stable par $\varinjlim$ de cardinal $\leqslant \pi$, $\Filt^\pi$ l'étant, donc $J$ est grande devant $\pi$, d'ailleurs $F = \varinjlim\limits_{j\in J} X_j$, où $j \longmapsto X_j$ est le foncteur source $\Filt^\pi/F \longrightarrow \Filt^\pi$. D'où la représentation voulue de $F$, quitte à remplacer $J$ par un ensemble ordonné cofinal $I$, p.ex. en prenant pour $I$ l'ensemble des sous-catégories $J$ de $I$ telles que $\card \Fl J \leqslant \pi$ \texttt{[cette deuxième partie de la phrase “p.ex. ...” demande correction]}. La majoration $\card I \leqslant \pi'^{\pi}$, pour $F \in \Ob \Filt^{\pi'}$, se fait par AQT \texttt{[âne qui trotte]}.

L'implication 3) $\Longrightarrow$ 4) est tautologique.

4) $\Longrightarrow$ 5) (\footnote{comparer le lemme 4.5.14, page 133.}), car si $F \in \Ob \mathcal{M}_\pi$, on trouve par 4) que $F$ est facteur direct d'un objet de $\Filt^\pi$, et $\Filt^\pi$ étant stable par $\varinjlim$
\msp{page 118}
grandes devant $\pi_0$, cela implique que $\Filt^\pi$ est stable par facteurs directs, d'où $F \in \Ob \Filt^\pi$.

Enfin, si $\mathcal{M} = \Ens$, $C = \Ens(\pi_0)$. On note que
\[
\Filt^\pi \ = \ \{ \ X \in \Ob \Ens \mid \card\,(X)^{\pi_0} \ \leqslant \ \pi \ \} \ .
\]
D'autre part, on trouve facilement
\[
\begin{aligned}
(\Ens)_\pi \quad &\overset{\text{déf}}{=} \quad \text{sous-catégorie pleine de (Ens) formée} \\
&\phantom{\overset{\text{déf}}{=}} \quad \text{des objets $\pi$-accessibles} \\
&= \quad \{ X \in \Ob \Ens \mid \card X \leqslant \pi \} \ .
\end{aligned}
\]
Donc on a $(\Ens)_\pi \subset \Filt^\pi$ si et seulement si $\pi^{\pi_0} \leqslant \pi$, i.e. $\pi^{\pi_0} = \pi$, q.e.d. (\footnote{Cf. démonstration de la partie 1°) de 4.5.6, plus bas (bas de page et page suivante).}).

\textbf{Corollaire 4.5.7.} \textit{Pour tout cardinal $\pi \geqslant \pi_1$, on a}
\[
\Filt^\pi \ \subset \ \mathcal{M}_\pi \ \subset \ \Filt^{\pi^{\pi_0}} \ . \tag*{(4.5.7.1)}
\]

L'inclusion $\Filt^\pi \subset \mathcal{M}_\pi$ est l'implication a $\Longrightarrow$ c de la partie 1°) de 4.5.6. Pour prouver $\mathcal{M}_\pi \subset \Filt^{\pi^{\pi_0}}$, on pose $\pi' = \pi^{\pi_0}$ et on note que $\pi'^{\pi_0} = \pi^{\pi_0 \cdot \pi_0} = \pi^{\pi_0} = \pi'$. Donc on peut
\printedpage{35}
% head: A. Grothendieck : Dérivateurs, Ch. XVIII, Édition des Universités de Montpellier II et Paris VII
% foot: Version 17 décembre 2025
appliquer 1° à $\pi'$ au lieu de $\pi$, et on trouve que $\mathcal{M}_{\pi'} = \Filt^{\pi'}$. Comme $\mathcal{M}_\pi \subset \mathcal{M}_{\pi'}$, on en conclut $\mathcal{M}_\pi \subset \Filt^{\pi'}$, q.e.d.

J'ai oublié de démontrer le 1°) dans 4.5.6. Vu le diagramme d'implications $(\ast)$, p. 108, et l'implication $\mathrm{a}'' \Longrightarrow \mathrm{a}$ déjà prouvé (p. 109),
\msp{page 119}
il reste à prouver seulement
\[
\mathrm{c}' \ \Longrightarrow \ \mathrm{a}'' \ ,
\]
i.e. que $F$ $\pi$-accessible \textit{dans} $C^\wedge$ implique $\card F(x) \leqslant \pi \ \forall \ x \in C$. Ceci va résulter du résultat un peu plus général :

\textbf{Lemme 4.5.8.}
\begin{itemize}[leftmargin=3em]
\item[\textit{1°)}] \textit{Dans la catégorie $\mathcal{E} = $} (Ens), \textit{les objets $\pi$-accessibles ($\pi$ un cardinal infini) sont les ensembles $X$ tels que} $\card X \leqslant \pi$.
\item[\textit{2°)}] (\footnote{Cf. généralisation 4.14.18, p. 391 b) et c).}). \textit{Soit $C$ une catégorie équivalente à une petite catégorie, telle que $\forall \ x, y \in \Ob C$, on ait} $\card \Hom(x,y) \leqslant \pi$. \textit{Soit $\mathcal{E}$ une $\mathfrak{U}$-catégorie stable par $\varinjlim$ grandes devant $\pi$, et faiblement $\pi$-accessible (i.e. admettant une sous-catégorie génératrice formée d'objets $\pi$-accessibles), et stable par produits de cardinal $\leqslant \pi$. Alors pour qu'un objet $F$ de $\uHom(C,\mathcal{E})$ soit $\pi$-accessible, il} faut \textit{que $\forall \ x \in \Ob C$, $F(x)$ soit un objet $\pi$-accessible de $\mathcal{E}$, et cette condition est aussi suffisante si} $\card \operatorname{\text{réd}} \Fl(C) \leqslant \pi$ (\footnote{même \textit{sans} condition de stabilité de $\mathcal{E}$ par produits de cardinal $\leqslant \pi$.}).
\end{itemize}

\textsc{Démonstration.} 1°) $\card X \leqslant \pi \Longrightarrow X$ $\pi$-accessible dans $\mathcal{E} = (\Ens)$, car $X = \coprod_{x \in X} \{x\}$, et on applique SGA 4 I 9.9. Inversement, supposons $X$ $\pi$-accessible, et écrivons $X \simeq \varinjlim\limits_{i \in I} X_i$, où $I$ est l'ensemble des parties de $X$ qui sont de cardinal $\leqslant \pi$. Comme $I$ est grande devant $\pi$, l'application identique $X \xrightarrow{\ \mathrm{id}\ } X$ se factorise par un ensemble $X_i$, donc $X_i = X$, donc $\card X = \card X_i \leqslant \pi$, q.e.d.

\msp{page 120}

2°) Considérons, pour $x \in \Ob C$, le foncteur
\[
e \xrightarrow{\ \varepsilon_x\ } C \ ,
\]
d'où un foncteur
\[
\varepsilon_x^* \ : \ \uHom(C,\mathcal{E}) \longrightarrow \mathcal{E} \ , \qquad F \longmapsto F(x) \ ,
\]
admettant un adjoint à droite
\[
\varepsilon_{x\,*} \ : \ \mathcal{E} \ \longrightarrow \ \uHom(C,\mathcal{E}) \ ,
\]
donné par
\[
\varepsilon_{x\,*}(E)(y) \ = \ E^{\Hom(x,y)}
\]
(sous réserve d'existence des produits de cardinal $\leqslant \pi$). Ce foncteur $\varepsilon_{x\,*}$ est $\pi$-accessible en vertu de loc. cit. 9.8.

Ceci posé, soit $F$ un objet $\pi$-accessible de $\uHom(C,\mathcal{E})$, on voit que $\varepsilon_x^*(F)$ est $\pi$-accessible, $\varepsilon_x^*$ ayant un adjoint à droite $\pi$-accessible :
\[
E \ \longmapsto \ \Hom_{\mathcal{E}}(\underbrace{F(x)}_{\varepsilon_x^*(F)}, E) \ \simeq \ \Hom_{\uHom(C,\mathcal{E})}(F, \varepsilon_{x\,*}(E))
\]
\printedpage{36}
% head: A. Grothendieck : Dérivateurs, Ch. XVIII, Édition des Universités de Montpellier II et Paris VII
% foot: Version 17 décembre 2025
est composé du foncteur $\pi$-accessible $\varepsilon_{x\,*}$, et du foncteur covariant $\pi$-accessible représenté par $F$ dans $\uHom(C,\mathcal{E})$.

Inversement, si les $F(x)$ sont $\pi$-accessibles, prouvons que $F$ l'est, moyennant $\card \Fl(C) \leqslant \pi$. Se fait par AQT \texttt{[âne qui trotte]}, je ne vais pas l'expliciter ici.

\msp{page 121}

\dgrblock{dgrb988}{The presentability of a diagram}
\begin{dcomment}{exposition}{p120}{The presentability of a diagram}
Grothendieck wants to recover the sizes of the values of a presheaf from the size of the presheaf as an object. Evaluation supplies the necessary bridge, but its adjoint must have the correct variance. For a covariant functor $F:C\to\mathcal E$, the right adjoint to evaluation at $c$ is
\[
 R_c(E)(d)=\prod_{a:d\to c}E.
\]
The edition writes $E^{\Hom(c,d)}$. That expression has the opposite variance in $d$ and is not the asserted right adjoint. To verify the corrected formula directly, a map $t:F(c)\to E$ determines the component $F(d)\to R_c(E)(d)$ whose projection at $a:d\to c$ is $tF(a)$. Naturality recovers every component from the projection at $\mathrm{id}_c$, establishing the adjunction. This is the pointwise right Kan extension along $\{c\}\to C$; compare \citet[Proposition 6.1.5 and Theorem 6.2.1]{Riehl2016}.

Suppose $\mathcal E$ has the products and $\mu$-filtered colimits used here, and has a conservative family of representables preserving those colimits. This is the meaning of the generating-family hypothesis inherited from SGA~4, rather than merely faithfulness of those representables. Products indexed by at most $\pi$ elements commute with $\mu$-filtered colimits in $\mathcal E$: applying a conservative representable reduces the assertion to sets. In sets, lift the at most $\pi$ coordinates at a common stage to obtain surjectivity, and witness their at most $\pi$ equalities at a common later stage to obtain injectivity. Thus $R_c$ preserves $\mu$-filtered colimits when every $\Hom(d,c)$ has cardinality at most $\pi$. If $F$ is $\mu$-presentable, then
\[
 \Hom_{\mathcal E}(F(c),-)
   \cong\Hom_{[C,\mathcal E]}(F,R_c(-))
\]
preserves these colimits, so $F(c)$ is $\mu$-presentable.

For the converse, assume a skeleton of $C$ has at most $\pi$ arrows. The missing AQT argument is an equalizer computation. For any $G:C\to\mathcal E$,
\[
 \operatorname{Nat}(F,G)=
 \operatorname{Eq}\left(
 \prod_{c\in\Ob C}\Hom(F(c),G(c))
 \rightrightarrows
 \prod_{a:c\to d}\Hom(F(c),G(d))\right).
\]
The two maps send a family $(t_c)$ to $(G(a)t_c)$ and $(t_dF(a))$. If all $F(c)$ are $\mu$-presentable, every factor preserves $\mu$-filtered colimits in $G$, and so does the displayed equalizer: the total number of coordinates and compatibility conditions is at most $\pi$. Since diagram-category colimits are pointwise, this proves presentability of $F$. No products in $\mathcal E$ are required for this direction; only products of hom-sets are used. The source's parenthetical qualification is therefore exact.

This calculation also explains why the bound on the number of arrows cannot be silently replaced by a bound on the number of objects. One must make every naturality equation hold at the same stage. A one-object category may carry arbitrarily many such equations through its endomorphisms.
\end{dcomment}


\textbf{Proposition 4.5.9.} \textit{Soit $C$ une catégorie équivalente à une petite catégorie, et soit $\mathcal{M} = $} $\Ind'_{\pi_0}(C)$, \textit{$\pi_0$ un cardinal infini donné. On suppose $C$ stable par $\varinjlim$ de cardinal $\leqslant \pi_0$. Pour tout cardinal $\pi \geqslant \pi_0$, soit $\mathcal{M}_\pi \supset C$ la sous-catégorie pleine de $\mathcal{M}$ formée des objets $\pi$-accessibles de $\mathcal{M}$. Soit d'autre part} $\pi_1 = \sup(\pi_0, \card \Fl C)$. \textit{Alors $(\mathcal{M}_\pi)_{\pi \geqslant \pi_1}$ est une filtration cardinale de $\mathcal{M}$} (\footnote{\textbf{N.B.} Toutes les conditions des filtrations cardinales sont vérifiées à l'exclusion, quand $\pi \geqslant \pi_0$, de celle de la majoration de cardinaux $\card I \leqslant \pi'^{\pi}$ condition c).})\textit{, et de plus, pour $\pi \geqslant \pi_0$, $\mathcal{M}_\pi$ est stable dans $\mathcal{M}$ par $\varinjlim$ de cardinal $\leqslant \pi$.}

\textsc{Démonstration.} La stabilité de $\mathcal{M}_\pi$ est un rappel, cf. SGA 4 I 9.9. Comme on a
\[
\mathcal{M}_\pi \ \subset \ \Filt^{\pi^{\pi_0}}
\]
pour $\pi \geqslant \pi_1$ (cf. 4.5.7), on voit donc que les $\mathcal{M}_\pi$ sont équivalentes à des petites catégories (condition a) des filtrations cardinales). La condition b) des filtrations cardinales est contenue dans la propriété de stabilité des $\mathcal{M}_\pi$ qu'on vient de rappeler. Prouvons enfin la condition c). Soit donc $\pi \geqslant \pi_0$, et soit $x$ dans $\mathcal{M}$, considérons la catégorie
\[
I \ = \ \mathcal{M}_\pi/x \ ,
\]
elle est stable par $\varinjlim$ de cardinal $\leqslant \pi$, d'après la stabilité similaire de $\mathcal{M}_\pi$, donc $I$ est grande devant $\pi$ et équivalente à une petite catégorie, d'autre part, on a
\msp{page 122}
\[
x \ = \ \varinjlim_{i \in I} x_i \ ,
\]
les $x_i$ dans $\mathcal{M}_\pi$ ( $\varinjlim$ dans $\mathcal{M}$ des foncteurs $I = \mathcal{M}_\pi/x \longrightarrow \mathcal{M}_\pi$). Quitte à remplacer $I$ par une catégorie ordonnée cofinale, cela prouve la partie qualitative de c. Supposons enfin $x \in \Ob \mathcal{M}_{\pi'}$, où $\pi' > \pi$, montrons que l'on a alors
\[
\card I_0, \ \card I' \ \leqslant \ \pi'^{\pi} \qquad \text{(\footnotemark, \footnotemark)} \ , \tag*{(4.5.9.1)}
\]
\addtocounter{footnote}{-1}%
\footnotetext{où $I_0$, $I'_0$ sont des sous-catégories réduites de $I$, $I'$ telles que $I_0 \hookrightarrow I$, $I'_0 \hookrightarrow I'$ soient des équivalences \texttt{[Grothendieck a biffé l'indice 0 de $I'$ dans 4.5.9.1]}.}%
\stepcounter{footnote}%
\footnotetext{On ne le prouve que si $\pi \geqslant \pi_1$.}%
ce qui achèvera de prouver que l'on a une filtration cardinale. Il nous faut une majoration des cardinaux.

\textbf{Lemme 4.5.10} (\footnote{Donner énonce commun avec 4.5.11, et prouver d'abord 4.5.11.}). \textit{Soit $\pi$ un cardinal $\geqslant \pi_1$. Alors on a} % [sic] énonce -> énoncé, missing accent (in footnote 47)
\[
\card \Fl((\Filt^\pi)_0) \ \leqslant \ 2^\pi \ . \tag*{(4.5.10.1)}
\]

On a en effet, en désignant par $C^\wedge(\pi)$ la sous-catégorie strictement pleine de $C^\wedge$ formée des $F : C^\circ \longrightarrow \Ens$ avec $\card F(x) \leqslant \pi$ pour tout $x$ dans $C$ (i.e. $F$ se factorise par $\Ens(\pi)$), et par $C^\wedge(\pi)_0$ une
\printedpage{37}
% head: A. Grothendieck : Dérivateurs, Ch. XVIII, Édition des Universités de Montpellier II et Paris VII
% foot: Version 17 décembre 2025
\msp{page 123}
sous-catégorie réduite exhaustive,
\[
\card C^\wedge(\pi)_0 \ \leqslant \ 2^\pi \ , \tag*{(4.5.10.2)}
\]
d'où résulte (4.5.10.1), puisque pour $\pi \geqslant \pi'$, $\Filt^\pi$ est une sous-catégorie pleine de $C^\wedge(\pi')$. Pour établir (4.5.10.2), on note que l'ensemble des classes à isomorphie près de $\Ens(\pi')$ s'identifie à l'ensemble des cardinaux $\leqslant \pi'$, lequel ensemble est de cardinal $\leqslant \pi'$. Remplaçant $\Ens(\pi)$ par $\Ens(\pi')_0$, $C$ par $C_0$, on doit majorer $\card \Fl(\uHom(C_0, \Ens(\pi')_0))$. Or $\Ob \uHom(C_0, \Ens(\pi')_0)$ est contenu dans l'ensemble des applications de $\Fl C_0$ dans $\Fl(\Ens(\pi')_0)$, le premier est majoré par $\pi_1$, le deuxième par $\pi^\pi \times \pi \times \pi = \pi^{2\pi} = \pi^\pi = 2^\pi$ (\textbf{N.B.} $\pi < 2^\pi$, d'où $\pi^\pi \leqslant (2^\pi)^\pi = 2^{\pi^2} = 2^\pi$), donc $\Ob \uHom$ majoré par $(2^\pi)^{\pi_1} = 2^{\pi\pi_1} = 2^\pi$. D'autre part, si $F$, $G$ sont dans $\uHom(C_0, \Ens(\pi_0))$, $\Hom(F,G)$ est majoré par $(\pi^\pi)^{\pi_1} = \pi^\pi = 2^\pi$, donc $\Fl \uHom$ est majoré par $2^\pi \times 2^\pi \times 2^\pi$, q.e.d.

\textbf{Corollaire 4.5.11.} \textit{Soit $F \in \Ob \Filt^\pi$, $F' \in \Ob \Filt^{\pi'}$, avec $\pi' \geqslant \pi \geqslant \pi_1$. Alors}
\[
\card \Hom(F, F') \ \leqslant \ \pi'^{\pi} \ . \tag*{(4.5.11.1)}
\]

\msp{page 124}

C'est essentiellement la même démonstration, légèrement plus simple. (On aurait dû commencer par là.)

\textbf{Corollaire 4.5.12.} \textit{Si $\pi' \geqslant \pi \geqslant \pi_1$, on a pour $x \in \Ob \mathcal{M}_{\pi'}$}
\[
\card\,(\mathcal{M}_\pi/x)_0 \ \leqslant \ \pi'^{(\pi^{\pi_0})} \ . \tag*{(4.5.12)}
\]

En effet,
\[
\card \Ob\,(\mathcal{M}_\pi/x)_0 \ \overset{\text{par 4.5.11}}{\leqslant} \ \underbrace{\underbrace{\card \Ob\,(\mathcal{M}_\pi)_0}_{\leqslant 2^{(\pi^{\pi_0})} \text{par 4.5.10}} \times (\pi'^{\pi_0})^{\pi^{\pi_0}}}_{= \pi'^{(\pi^{\pi_0})}} \ ,
\]
i.e.
\[
\card \Ob\,(\mathcal{M}_\pi/x)_0 \ \leqslant \ \pi'^{(\pi^{\pi_0})} \ ,
\]
d'où
\[
\card \Fl(\mathcal{M}_\pi/x)_0 \ \leqslant \ \pi'^{(\pi^{\pi_0})} \times \pi'^{(\pi^{\pi_0})} \times (\pi'^{\pi_0})^{(\pi^{\pi_0})} \ \leqslant \ \pi'^{(\pi^{\pi_0})} \ .
\]
Cela donne donc (revenant à 4.5.9.1)
\[
\card I_0 \ \leqslant \ \pi'^{(\pi^{\pi_0})} \ .
\]

D'autre part, $I'$ peut se définir en termes de $I_0$ comme l'ensemble des sous-catégories de $I_0$ qui sont de cardinal $\leqslant \pi$, donc leur nombre est majoré par $(\card I_0)^\pi \ \leqslant \ (\pi'^{\pi^{\pi_0}})^\pi = \pi'^{\pi\pi^{\pi_0}} = \pi'^{(\pi^{\pi_0})}$.

Mais ce n'est pas tout à fait la majoration exigée, puisque
\[
\pi'^{\pi} \leqslant \pi'^{(\pi^{\pi_0})} \qquad \text{(\footnotemark)},
\]
\footnotetext{\textbf{N.B.} On a égalité si on a $\pi' = 2^{c'}$ avec $c' \geqslant \pi^{\pi_0}$, soit $\pi = 2^c$ avec $c \geqslant \pi_0$.}%
et qu'il est possible, a priori, que l'inégalité
\printedpage{38}
% head: A. Grothendieck : Dérivateurs, Ch. XVIII, Édition des Universités de Montpellier II et Paris VII
% foot: Version 17 décembre 2025
\msp{page 125}
soit stricte ! Apparemment, les ``majorations''
\[
\mathcal{M}_\pi \ \subset \ \Filt^{\pi^{\pi_0}} \ , \qquad \mathcal{M}_{\pi'} \ \subset \ \Filt^{\pi'^{\pi_0}}
\]
sont trop grossières. Aussi il s'agit de mieux appréhender les catégories $\mathcal{M}_\pi$, pour $\pi \geqslant \pi_0$ (\footnote{Mais on peut utiliser 4.5.10, 4.5.11 pour conclure que pour $x \in \Ob \mathcal{M}_{\pi'} \subset \Ob \Filt^{\pi'^{\pi_0}}$, on a $\card \Fl(C_0/x) \leqslant (\pi'^{\pi_0})^{\pi_0} = \pi'^{\pi_0}$, on aura donc $x = \varinjlim\limits_{J} x_j$ ($x_j \in \Ob C$, $\varinjlim$ filtrante, $\card \Fl J \leqslant \pi'^{\pi_0}$), et en introduisant l'ensemble $I$ des sous-catégories $J' \subset J$ qui sont de cardinal $\leqslant \pi$, et posant $x_{J'} = \varinjlim\limits_{j \in J'} x_j$, on trouve $x = \varinjlim\limits_{J' \in I} x_{J'}$, où cette fois $I$ est un \textit{ensemble ordonné} grand devant $\pi$, les $x_{J'} \in \mathcal{M}_\pi$, et $\card I \leqslant (\card J)^\pi \leqslant (\pi'^{\pi_0})^\pi = \pi'^{\pi\pi_0} = \pi'^{\pi}$, c'est la bonne majoration !}). Je vais hasarder une

\dgrblock{dgrb989}{The canonical filtration and its cardinal bound}
\begin{dcomment}{exposition}{p124}{The canonical filtration and its cardinal bound}
Grothendieck's preferred filtration is now the intrinsic one: $\mathcal M_\pi$ consists of the $\pi^+$-presentable objects. The strength of Proposition 4.5.9 is that it supplies the required presentations at every sufficiently large $\pi$, although the more concrete test $|\Hom(c,X)|\leq\pi$ may fail. The footnote on manuscript page 125 supplies the right construction; counting the whole slice too early obscures it.

First, $\mathcal M_\pi$ is closed under colimits with at most $\pi$ arrows. If $A=\operatorname{colim}_{j\in J}A_j$ and $|\Fl J|\leq\pi$, then $\Hom(A,-)=\lim_{J^{\mathrm{op}}}\Hom(A_j,-)$, and a limit of at most $\pi$ set-valued functors preserving $\pi^+$-filtered colimits again preserves them. This is the elementary closure result used throughout these pages; the historical reference is \citet[Expos\'e I, 9.8--9.9]{SGA4I1972}. Essential smallness follows either from the presheaf bound $\mathcal M_\pi\subseteq\Filt^{\pi^{\pi_0}}$ or from the presentation theorem proved below.

For a direct construction, write $X=\operatorname{colim}_{j\in J}c_j$ using $J=C_0/X$. Let $P$ be the poset of subcategories $K\subseteq J$ with at most $\pi$ arrows, and set
\[
 X_K=\operatorname{colim}_{k\in K}c_k\quad\text{in }\mathcal M.
\]
The category $P$ is $\pi^+$-directed. A family of at most $\pi$ such subcategories has a union of at most $\pi$ arrows; closing the union under identities and finite composition still adds at most $\pi$ arrows. Every $X_K$ belongs to $\mathcal M_\pi$. Moreover $X=\operatorname{colim}_{K\in P}X_K$: a compatible family of maps from the $X_K$ is exactly a cocone on the original diagram, since every object and arrow of $J$ occurs in some $K$. This proof does not assert that the subcategories $K$ are themselves $\pi_0^+$-filtered.

If $X\in\mathcal M_{\pi'}$, with $\pi'\geq\pi\geq\pi_1$, then $|\Fl(C_0/X)|\leq(\pi')^{\pi_0}$ by Corollary 4.5.7. The number of choices of $K$ is bounded by
\[
 |P|\leq\bigl((\pi')^{\pi_0}\bigr)^\pi
      =(\pi')^\pi.
\]
The number of order relations of this infinite poset has the same bound. This proves the quantitative clause of the cardinal filtration without assuming $\pi^{\pi_0}=\pi$.

The numerical lemmas on pages 122--124 are correct after normalizing the stray primes in the proof of 4.5.10. For presheaves with at most $\pi$ values at each object of a category with at most $\pi$ arrows, one chooses representative sets among the cardinals at most $\pi$, and then structure maps between them. There are at most $(2^\pi)^\pi=2^\pi$ resulting functors and at most $2^\pi$ natural transformations between each pair. Similarly, if the source values are bounded by $\pi$ and target values by $\pi'$, then
\[
 |\operatorname{Nat}(F,G)|\leq
 \prod_{c\in\Ob C_0}|G(c)|^{|F(c)|}
 \leq(\pi')^\pi.
\]
The weaker exponent in Corollary 4.5.12 is a valid intermediate estimate. It is not a counterexample to the sharper presentation bound: that bound counts the subdiagrams used above rather than every map from every presentable object at once.
\end{dcomment}


\textbf{Proposition 4.5.13.} \textit{Les hypothèses et notations sont celles de 4.5.9. Soit $\pi$ un cardinal $\geqslant \pi_0$, et soit $x \in \Ob \mathcal{M}$. Considérons les conditions :}
\begin{itemize}[leftmargin=3em]
\item[\textit{a})] \textit{Il existe une catégorie $I$, avec} $\card \Fl I \leqslant \pi$, \textit{et un foncteur $i \longmapsto x_i : I \longrightarrow C$, tel que $x$ soit isomorphe à $\varinjlim x_i$.}
\item[\textit{a}$'$)] \textit{Comme a), avec $I$ filtrante.}
\item[\textit{a}$''$)] \textit{Comme a), avec $I$ ordonnée filtrante.}
\item[\textit{b})] \textit{Il existe un foncteur} cofinal \textit{$I \longrightarrow C/x$, avec} $\card \Fl I \leqslant \pi$.
\item[\textit{b}$'$)] \textit{La catégorie filtrante $C/x$ est ``essentiellement de cardinal $\leqslant \pi$'', i.e. il existe une catégorie ordonnée} \texttt{[non, cf. b$''$)]} filtrante $I$, $\card \Fl I \leqslant \pi$, \textit{et un foncteur cofinal $I \longrightarrow C/x$.}
\item[\textit{b}$''$)] \textit{Comme b$'$), avec de plus $I$ ordonnée.}
\item[\textit{c})] \textit{Il existe une famille de flèches $(u_\alpha)_{\alpha \in A}$, $x_\alpha \longrightarrow x$, épimorphique stricte, indexée par un ensemble $A$ tel que} $\card A \leqslant \pi$.
\item[\textit{d})] \textit{On a $x \in \Ob \mathcal{M}_\pi$, i.e. $x$ est $\pi$-accessible dans $\mathcal{M}$.}
\end{itemize}

\msp{page 126}

\textit{Ceci posé :}
\begin{itemize}[leftmargin=3em]
\item[\circled{1\textsuperscript{\textit{o}}}] \textit{Les six premières conditions $a$, $a'$, $a''$, $b$, $b'$, $b''$ sont équivalentes, soit $P(\pi)$, et impliquent les deux dernières,}
\[
\begin{tikzcd}[column sep=small]
 & P(\pi) \arrow[dl, Rightarrow] \arrow[dr, Rightarrow] & \\
c & & d \ .
\end{tikzcd}
\]
\textit{Si $\pi \geqslant \pi_1$, alors $P(\pi)$ est équivalente à $c$.}
\item[\circled{2\textsuperscript{\textit{o}}}] \textit{Soit $\Filt'^{\pi} \subset \mathcal{M}$ la sous-catégorie strictement pleine de $\mathcal{M}$ formée des objets qui satisfont la condition $P(\pi)$ ci-dessus, et $\Filt^\pi$ la sous-catégorie strictement pleine définie dans 4.5.3 (p. 109). On a alors}
\[
\Filt^\pi \ \subset \ \Filt'^{\pi} \ \subset \ \mathcal{M}_\pi \ , \tag*{(4.5.13.1)}
\]
\textit{et pour qu'on ait égalité des trois membres, il suffit qu'on ait $\pi^{\pi_0} = \pi$. (Cf. prop. 4.5.6, page 116). De plus, $\Filt'^{\pi}$ est stable dans $\mathcal{M}$ par} sommes \textit{de cardinal $\leqslant \pi$.}
\item[\circled{3\textsuperscript{\textit{o}}}] \textit{Les conditions suivantes sur $C$, $\pi_0$, $\pi$ sont équivalentes.}
\end{itemize}

% ---- from der-ch18_p041-045.tex
\setcounter{footnote}{49}

\printedpage{39}
% head: A. Grothendieck : Dérivateurs, Ch. XVIII, Édition des Universités de Montpellier II et Paris VII

\begin{itemize}
\item[\textit{1)}] $\Filt'^{\pi}$ \textit{est stable dans $\mathcal{M}$ par} $\varinjlim$ \textit{de cardinal $\leqslant \pi$.}
\item[\textit{1$'$)}] $\Filt'^{\pi}$ \textit{est stable par conoyaux de double-flèches.}
\item[\textit{2)}] \textit{Tout objet $x$ de $\mathcal{M}$ est de la forme $\varinjlim\limits_{I} x_i$, où $x_i$ dans} $\Filt'^{\pi}$\textit{, $I$ grande devant $\pi$.}
\item[\textit{3)}] \textit{La condition d) de 1°) implique $P(\pi)$ (donc lui est équivalente), i.e. on a}
\[ \mathcal{M}_{\pi} \ = \ \Filt'^{\pi}. \]
\end{itemize}

\msp{page 127}
\textit{Ces conditions sont satisfaites dans les trois cas suivants :}
\begin{itemize}
\item[\textit{A)}] \textit{$\pi^{\pi_0} = \pi$ (pour mémoire, c'est contenu dans 2°).}
\item[\textit{B)}] \textit{Les familles épimorphiques strictes dans $\mathcal{M}$ sont épimorphiques strictes universelles.}
\item[\textit{C)}] \textit{Les limites inductives filtrantes dans $\mathcal{M}$ commutent au changement de base (ce qui est le cas notamment si les} $\varinjlim$ \textit{filtrantes sont exactes à gauche)} (\footnote{Je suppose que cette condition revient à la même condition sur $C$, pour les $\varinjlim$ filtrantes de cardinal $\leqslant \pi$.}).
\end{itemize}

\dgrblock{dgrb990}{A small colimit need not be cofinal in the canonical presentation}
\begin{dcomment}{development}{p126}{A small colimit need not be cofinal in the canonical presentation}
Grothendieck is looking for a presentation whose size can be read directly from the object. There are two different requirements in Proposition 4.5.13: being the colimit of a small diagram of objects of $C$, and possessing a small cofinal diagram in $C/X$. The proof on page 129 establishes the first requirement with a filtered indexing poset; it does not establish cofinality in $C/X$. The distinction persists even in sets, where all the exactness hypotheses later imposed are satisfied.

\textbf{A counterexample to a)$\Rightarrow$b).}\ours{}
Take $\pi_0=\aleph_0$, let $C$ be a skeleton of the countable sets, and let $X$ have cardinality $\pi=\beth_\omega$. Here $\pi\geq\pi_1$, $C$ has countable colimits, and $\mathcal M=\Ens$. The set $X$ is the coproduct of its $\pi$ singleton subsets. It is also the filtered colimit of its finite subsets, indexed by a poset with $\pi$ arrows. Thus a), a$'$), and a$''$) all hold.

Suppose a category $I$ with at most $\pi$ arrows admitted a cofinal functor $I\to C/X$. Write its objects as maps $u_i:c_i\to X$. For every countable subset $A\subseteq X$, nonemptiness of the comma category $(A\hookrightarrow X)\downarrow I$ would give some $i$ with $A\subseteq u_i(c_i)$. The at most $\pi$ countable subsets $u_i(c_i)$ would therefore be cofinal among the countable subsets of $X$.

No such family exists. Indeed, write any family $\mathcal H$ of at most $\pi$ countable subsets as an increasing union $\bigcup_{n<\omega}\mathcal H_n$ with $|\mathcal H_n|<\pi$. Each union $\bigcup\mathcal H_n$ has cardinality less than $\pi$. Choose $s_n\in X\setminus\bigcup\mathcal H_n$. The countable set $\{s_n:n<\omega\}$ is contained in no member of $\mathcal H$: a member belonging to $\mathcal H_n$ omits $s_n$. This contradiction proves failure of b), hence also of b$'$) and b$''$). The choice $\beth_\omega$ avoids any continuum hypothesis; its countable cofinality and its being larger than $2^{\aleph_0}$ are sufficient.

For the corrected positive assertion, let $X=\operatorname{colim}_I c_i$ with $|\Fl I|\leq\pi$. Every finite set of arrows of $I$ gives a finite incidence diagram: introduce one vertex for each object used and one extra vertex for the source of each selected arrow, with its identity arrow to the source vertex and the given arrow to the target vertex. Its colimit belongs to $C$, because $C$ has finite colimits. Ordering finite sets of arrows by inclusion gives a filtered poset with at most $\pi$ arrows, and the colimit of these finite presentations is $X$. A cocone on all of them is exactly a cocone on $I$. This establishes a)$\Rightarrow$a$''$), including categories with loops: finite subsets need not generate finite subcategories, which is why the incidence-diagram description is preferable.

Thus a), a$'$), and a$''$) are equivalent. Each of b), b$'$), and b$''$) implies the corresponding colimit condition, but the two triples cannot be identified. For the rest of our commentary, $\Filt'^\pi$ denotes the one-step colimit class defined by a), the class used in the subsequent presentation arguments. This resolves an ambiguity created by the false equivalence in the printed definition; it is not a change to the source text.
\end{dcomment}


\textbf{Remarque.} La filtration de $\mathcal{M}$ pour les $\Filt'^{\pi}$ n'est autre que celle envisagée dans SGA 4 I 9.13, où on fait de plus l'hypothèse B) sur $\mathcal{M}$, qui après-coup n'était plus tellement de mon goût. Toujours est-il que d'après (4.5.13.1), cette filtration se rapproche plus de la ``bonne'' filtration par les $\mathcal{M}_{\pi}$, que celle que j'avais eu encore naguère tendance à lui préférer, celle par les $\Filt^{\pi}$. Elle a d'ailleurs sur cette dernière l'avantage, dans le cas envisagé B), d'être une filtration cardinale (ce qui n'est pas nécessairement le cas pour $\Filt^{\pi}$, du moins s'il existe des cardinaux infinis $\pi \geqslant \pi_0$ avec $\pi^{\pi_0} \neq \pi$). Je soupçonne maintenant qu'en général (et s'il existe des cardinaux infinis $\pi > \pi_0$ tels que $\pi^{\pi_0} \neq \pi$), en dehors des conditions A) B) C), la filtration par les $(\Filt'^{\pi})$ n'est \textit{pas} une
\msp{page 128}
filtration cardinale. En tous cas, 4.5.6 et la présente proposition, 3°, montrent que si soit $(\Filt^{\pi})_{\pi\geqslant\pi_2}$, soit $(\Filt'^{\pi})_{\pi\geqslant\pi_2}$ est une filtration cardinale ($\pi_2$ étant un cardinal donné $\geqslant \pi_1$), alors elle coïncide avec la filtration par les $\mathcal{M}_{\pi}$. Donc apparemment, dans tous les cas où je sais construire une filtration cardinale, sur un paratopos du moins (i.e. sur une catégorie $\mathcal{M}$ stable par petites $\varinjlim$ , et dont tous les objets sont accessibles -- l'existence d'une filtration cardinale impliquant de plus celle d'une petite sous-catégorie strictement génératrice \dots), celle-ci n'est autre que celle par les $\mathcal{M}_{\pi}$ -- chose dont je ne m'étais par aperçu en rédigeant SGA 4 I 9. % [sic] m'étais par aperçu -> m'étais pas aperçu, "par" for the negation "pas"

\dgrblock{dgrb991}{A precise reduction of universality to the generators}
\begin{dcomment}{development}{p127:n'est autre que celle envisagée dans SGA 4 I 9.13}{A precise reduction of universality to the generators}
Footnote 50 suggests testing condition C) inside the small generating category. As printed, that suggestion needs a qualification: $C$ is only assumed to have colimits with at most $\pi_0$ arrows, and it need not have pullbacks. With the latter hypothesis added, the reduction is valid, and the bound $\pi_0$ itself suffices. The proof below makes explicit which data must be approximated together.

\textbf{A bounded test for universal filtered colimits.}\ours{}
Let $\kappa$ be infinite, let $C$ be essentially small with pullbacks and all colimits indexed by categories with at most $\kappa$ arrows, and put
\[
 \lambda=\kappa^+,
 \qquad \mathcal M=\Ind'_{\kappa}(C).
\]
Then ordinary filtered colimits in $\mathcal M$ are universal if and only if filtered colimits in $C$ indexed by categories with at most $\kappa$ arrows are universal in $C$. Here universality in $C$ means stability under every base change that is a morphism of $C$. Thus the additional hypothesis is existence of pullbacks in $C$; no terminal object or colimits of the larger bound $\pi$ are needed. The marker identifies the present formulation and proof, without a priority claim.

We first recall the presheaf facts that control the proof. The category $\mathcal M$ consists of the presheaves on $C$ that carry colimits with at most $\kappa$ arrows in $C$ to limits of sets. This description follows from the canonical category-of-elements presentation: the preservation condition makes that category $\lambda$-filtered, and conversely a $\lambda$-filtered colimit of representables has the preservation condition. Limits of such presheaves, and $\lambda$-filtered colimits of them, are computed pointwise. Indeed, limits commute with limits, and $\lambda$-filtered set colimits commute with limits having at most $\kappa$ arrows. Consequently $\lambda$-filtered colimits in $\mathcal M$ commute with finite limits. The embedding $j:C\to\mathcal M$ preserves pullbacks, by Yoneda and the pointwise description of limits. It also preserves the prescribed colimits: for any such colimit $c=\operatorname{colim}_i c_i$ in $C$ and any $F\in\mathcal M$,
\[
 \Hom_{\mathcal M}(j(c),F)=F(c)
   =\lim_i F(c_i)
   =\lim_i\Hom_{\mathcal M}(j(c_i),F).
\]

We need a simultaneous diagram approximation. If $J$ has at most $\kappa$ arrows, every $J$-diagram in $\mathcal M$ is a $\lambda$-filtered colimit of $C$-valued $J$-diagrams. Here is a direct verification of this instance of the diagram construction used earlier. Let $E$ be the full subcategory of $\mathcal M^J$ consisting of the $C$-valued diagrams, up to isomorphism. It is essentially small and closed under colimits with at most $\kappa$ arrows, since those colimits are computed pointwise and preserved by $j$. For a diagram $D$, the category $E/D$ is therefore $\lambda$-filtered: take the colimit of any diagram of at most $\kappa$ arrows in that slice. The empty diagram is handled by its initial object.

To see that the canonical diagram $E/D\to\mathcal M^J$ has colimit $D$, use the free diagrams
\[
 L_a(c)(b)=\coprod_{J(a,b)}c
 \qquad(a,b\in J,\ c\in C).
\]
They belong to $E$, since every indexing set $J(a,b)$ has at most $\kappa$ elements, and $L_a$ is left adjoint to evaluation at $a$. A cocone from $E/D$ to a diagram $G$ therefore assigns compatible maps $c\to G(a)$ to all maps $c\to D(a)$ from $C$. The canonical presentation of $D(a)$ gives a unique map $t_a:D(a)\to G(a)$. For an arrow $a\to b$ in $J$, the induced maps $L_b(c)\to L_a(c)$ show that the $t_a$ satisfy naturality. Finally, for any $A\to D$ in $E/D$, testing with all maps $L_a(c)\to A$ shows that its given cocone map is $A\to D\xrightarrow{t}G$. Those tests suffice because the maps from $C$ to every $A(a)$ are jointly epimorphic. This proves the universal property and hence the simultaneous approximation. A small skeleton of $E/D$ supplies a small filtered index.

Assume now the bounded universality hypothesis in $C$. First consider a filtered category $I$ with at most $\kappa$ arrows, a diagram $D:I\to\mathcal M$ with a cocone to $Z$, and a map $T\to Z$. We will prove the more flexible comparison
\[
 \operatorname{colim}_{i\in I}(T\times_Z D_i)
   \longrightarrow T\times_Z\operatorname{colim}_{i\in I}D_i
 \tag{*}
\]
is an isomorphism. Adjoin to $I$ the cocone vertex $z$, all arrows $i\to z$, and one further vertex $t$ with an arrow $t\to z$, imposing the cocone relations. The resulting shape still has at most $\kappa$ arrows. Apply the simultaneous approximation to this entire shape. We obtain a $\lambda$-filtered diagram, indexed by $p$, of data
\[
 D_p:I\to C,\qquad D_p\longrightarrow Z_p,
 \qquad T_p\longrightarrow Z_p,
\]
whose colimit is exactly the original data. In particular, one does not assume that $Z_p$ is the colimit of $D_p$.

Put $Q_p=\operatorname{colim}_I D_p$ in $C$. Applying the assumed universality to the base change
\[
 T_p\times_{Z_p}Q_p\longrightarrow Q_p
\]
gives, in $C$ and hence in $\mathcal M$,
\[
 \operatorname{colim}_{i\in I}(T_p\times_{Z_p}D_{p,i})
   \simeq T_p\times_{Z_p}Q_p.
\]
All objects and pullbacks in this equation lie in $C$, and the preceding facts ensure that its colimit and pullbacks are preserved by $j$. Take the $\lambda$-filtered colimit over $p$. On the left, interchange the two colimits and commute the $p$-colimit with each pullback. On the right, commute it with the displayed pullback and use
\[
 \operatorname{colim}_p Q_p
   =\operatorname{colim}_p\operatorname{colim}_i D_{p,i}
   \simeq\operatorname{colim}_i D_i.
\]
The resulting isomorphism is exactly (*). Setting $Z=\operatorname{colim}_I D_i$ proves universality for every bounded filtered diagram in $\mathcal M$.

It remains to remove the bound on $I$. Let $I$ be any small filtered category, and let $P$ be the poset of its filtered subcategories with at most $\kappa$ arrows. Every set of at most $\kappa$ arrows of $I$ lies in such a subcategory. Start by closing the selected arrows under identities and finite composition. Choose in $I$ a cocone for every finite diagram in the resulting subcategory, add its objects and arrows, and close under composition again. Repeat this operation countably many times. At each step there are at most $\kappa$ finite diagrams and at most $\kappa$ new arrows; the final union still has at most $\kappa$ arrows. Every finite diagram in the union occurs at some stage and receives a cone at the next, so the union is filtered. If the initial selection is empty, first choose one object of $I$.

The same construction applied to the union of at most $\kappa$ members of $P$ proves that $P$ is $\lambda$-directed. Every object and arrow of $I$ occurs in some member of $P$, so for any $D:I\to\mathcal M$ the cocone universal property gives
\[
 \operatorname{colim}_{K\in P}\operatorname{colim}_{i\in K}D_i
   \simeq\operatorname{colim}_{i\in I}D_i.
\]
Apply the already proved comparison (*) to each restricted diagram $D|_K$, with the fixed cocone to $Z$ and the fixed arrow $T\to Z$. Taking the $\lambda$-filtered colimit over $P$, and again commuting it with pullback, proves (*) for arbitrary filtered $I$. This is condition C).

Conversely, if condition C) holds in $\mathcal M$, apply it to a bounded filtered diagram in $C$ and a base-change arrow in $C$. The embedding $j$ preserves its colimit and pullbacks and reflects isomorphisms, so the corresponding comparison in $C$ is an isomorphism. This proves the converse and the theorem.

With $\kappa=\pi_0$, this gives a precise affirmative form of footnote 50. Without the pullback hypothesis the phrase ``the same condition on $C$'' needs a relative formulation using pullbacks calculated in $\mathcal M$; it is not the intrinsic statement just proved. The argument addresses the proposed reduction of C), and makes no assertion about the separate implication B)$\Rightarrow$C).
\end{dcomment}


\textsc{Démonstration de 4.5.13.}

\textit{Prouvons 1°.} On a les implications tautologiques
\[
\begin{tikzcd}[row sep=large]
\mathrm{b}'' \arrow[r, Rightarrow] \arrow[d, Rightarrow] & \mathrm{b}' \arrow[r, Rightarrow] \arrow[d, Rightarrow] & \mathrm{b} \arrow[d, Rightarrow] \\
\mathrm{a}'' \arrow[r, Rightarrow] & \mathrm{a}' \arrow[r, Rightarrow] & \mathrm{a} \text{ ,}
\end{tikzcd}
\qquad\qquad \text{(\footnotemark) ,}
\]
\footnotetext{compte tenu que $C$ est strictement génératrice dans $\mathcal{M}$, donc $x \simeq \varinjlim\limits_{i\in C/x} x_i$.}
et pour prouver l'équivalence de ces six conditions, il suffit de prouver
\[ \mathrm{a} \ \Longrightarrow \ \mathrm{b}'' \ . \]
% foot: Version 17 décembre 2025

\printedpage{40}
% head: A. Grothendieck : Dérivateurs, Ch. XVIII, Édition des Universités de Montpellier II et Paris VII

\msp{page 129}
Si on a
\[ x \ \simeq \ \varinjlim_{I} x_i \ , \qquad\quad \card \Fl I \ \leqslant \ \pi \ , \]
soit $I'$ l'ensemble des \textit{sous-diagrammes} $J \subset I$ finis de $I$, ordonné par inclusion, donc $I'$ est un ensemble ordonné filtrant, et du fait que $C$ est stable par $\varinjlim$ finies, on trouve un foncteur
\[ J \longmapsto x_J \ = \ \varinjlim_{i\in J} x_i \ : \ I' \longrightarrow C \ , \]
et on aura
\[ x \ \simeq \ \varinjlim_{J\in I'} x_J \ . \]
D'autre part, on a
\[
\begin{aligned}
\card I' \ &= \ \card \mathfrak{P}_f( \ \underbrace{\Fl(I)}_{\text{de card $\leqslant \pi$}} \ ) \\
&\leqslant \ \pi + \pi^2 + \ \cdots \ + \pi^n + \ \dots \\
&\leqslant \ \pi\aleph_0 \ = \ \pi \ ,
\end{aligned}
\]
q.e.d. Pour établir 1°, il reste à établir
\[
\begin{tikzcd}[column sep=small]
& P(\pi) \arrow[dl, Rightarrow] \arrow[dr, Rightarrow] & \\
\mathrm{c} & & \mathrm{d} \text{ ,}
\end{tikzcd}
\]
et l'implication inverse $\mathrm{c} \Longrightarrow P(\pi)$ quand $\pi \geqslant \pi_1$. Mais $P(\pi) \Longrightarrow \mathrm{c}$ est une tautologie, et $P(\pi) \Longrightarrow \mathrm{d}$ puisque $\mathcal{M}_{\pi}$ est stable par $\varinjlim$ de cardinal $\leqslant \pi$. Reste à prouver que $\mathrm{c} \Longrightarrow P(\pi)$ si $\pi \geqslant \pi_1$.

\dgrblock{dgrb992}{One step of colimits already gives all presentable objects}
\begin{dcomment}{development}{p128}{One step of colimits already gives all presentable objects}
The useful question remaining after the preceding counterexample is whether the one-step colimit class $\Filt'^\pi$ is smaller than $\mathcal M_\pi$. In fact the two classes coincide for every $\pi\geq\pi_0$. The general presentation theorem is known; it is discussed as Remark 1.30 of Ad\'amek--Rosick\'y in the primary proof of \citet{Barton2019Presentable}. We give a direct proof in Grothendieck's setting. The marker below identifies our formulation and proof, rather than a claim that the theorem is new.

\textbf{A bounded presentation theorem, with an explicit flattening proof.}\ours{}
Suppose $C$ is essentially small and closed under colimits with at most $\pi_0$ arrows, and $\mathcal M=\Ind'_{\pi_0}(C)$. Let $D_\pi$ be the full subcategory of objects expressible as a colimit in $\mathcal M$ of a $C$-valued diagram with at most $\pi$ arrows. Then
\[
 D_\pi=\mathcal M_\pi\qquad(\pi\geq\pi_0).
\]
In particular $D_\pi$ is closed under all colimits with at most $\pi$ arrows, without the extra hypotheses A), B), or C) of Proposition 4.5.13.

The point requiring proof is that a map into a presentation can itself be recorded using only boundedly many of its data. Let $Y=\operatorname{colim}_{j\in J}d_j$, with $d_j\in C$, and let $P_0(J)$ be the poset of subcategories of $J$ with at most $\pi_0$ arrows. This poset is $\pi_0^+$-directed: take unions of at most $\pi_0$ arrow sets and close under finite composition. For $K\in P_0(J)$, choose $d_K=\operatorname{colim}_K d_j$ in $C$. The embedding $C\to\mathcal M$ preserves these colimits. Their universal properties give compatible transition maps, and
\[
 Y=\operatorname{colim}_{K\in P_0(J)}d_K.
\]
This equality follows by identifying cocones, since every object and arrow of $J$ appears in some $K$. Every $c\in C$ is $\pi_0^+$-presentable in $\mathcal M$. Hence a map $c\to Y$ factors through some $d_K$. Two maps can be factored through the same $d_K$ by passing to a common upper bound. Notice that $K$ itself need not be filtered.

We now prove closure under coequalizers. Let $f,g:X\rightrightarrows Y$ be maps between objects of $D_\pi$, and choose presentations
\[
 X=\operatorname{colim}_{i\in I}c_i,
 \qquad Y=\operatorname{colim}_{j\in J}d_j,
 \qquad |\Fl I|,|\Fl J|\leq\pi.
\]
Writing $s_i:c_i\to X$ for the first cocone, choose, for every $i$, one $K_i\in P_0(J)$ and two arrows
\[
 u_i,v_i:c_i\rightrightarrows d_{K_i}
\]
whose composites into $Y$ are $fs_i$ and $gs_i$. Construct a graph of objects of $C$ as follows. Retain the objects and arrows of the $J$-diagram. For each $i$, add a fresh vertex with value $d_{K_i}$ and the complete colimit cone $d_j\to d_{K_i}$ for $j\in K_i$. Add a further fresh vertex with value $c_i$, and the two arrows $u_i,v_i$ into that $d_{K_i}$. The free category on this graph, equipped with the assigned composites in $C$, has at most $\pi$ arrows. Indeed, there are at most $\pi$ choices of $i$, each adds at most $\pi_0$ edges, and taking finite paths preserves the infinite bound $\pi$. One may also impose the original diagram relations; the bound is unchanged.

A cocone on the original $J$-diagram with target $Z$ is a map $h:Y\to Z$. Its extension to every new vertex $d_{K_i}$ is forced uniquely, because the displayed cone on $K_i$ is a genuine colimit cone in $C$ and in $\mathcal M$. Its extension to the new copy of $c_i$ exists precisely when $hfs_i=hgs_i$, and is then unique. These equalities for all $i$ are equivalent to $hf=hg$, since the presentation maps $s_i$ are jointly epimorphic. Thus cocones on the augmented diagram are exactly maps out of $\operatorname{coeq}(f,g)$. This proves that this coequalizer belongs to $D_\pi$. No compatibility between the choices of the $K_i$ is needed: only equations on the jointly epimorphic presentation maps are being imposed.

The class $D_\pi$ is also closed under coproducts of at most $\pi$ objects, by taking disjoint unions of their indexing categories. The standard expression of a colimit as a coequalizer of two coproducts now proves closure under every diagram with at most $\pi$ arrows. The inclusion $D_\pi\subseteq\mathcal M_\pi$ follows from closure of presentable objects under these colimits.

For the converse, express an arbitrary $X$ as a colimit of a small $C$-valued diagram. Its subcategories of at most $\pi$ arrows give a $\pi^+$-filtered presentation of $X$ by objects of $D_\pi$, exactly as in the preceding block. If $X\in\mathcal M_\pi$, its identity factors through one term, so $X$ is a retract of an object $Y\in D_\pi$. A retract $X\xrightarrow{i}Y\xrightarrow{r}X$, $ri=1$, is the coequalizer of $1_Y$ and $ir$. Closure under coequalizers gives $X\in D_\pi$, completing the proof.

This supplies the positive answer sought on pages 127--128 and 138. Iterated closure and one-step closure agree here; the bounded cofinal class and the pointwise presheaf bound can still be smaller. The three notions differ because only the first permits the additional relations introduced by the flattening construction.
\end{dcomment}


\msp{page 130}
Supposons donc qu'on ait une famille épimorphique stricte
\[ x_{\alpha} \longrightarrow x \qquad\quad (\alpha \in A,\ \card A \leqslant \pi) \ , \]
cela signifie donc que $x$ est limite inductive d'un système inductif indexé par la catégorie $B$ ayant comme ensemble d'objets $A \amalg (A\times A)$, les flèches non identiques étant représentées par les flèches canoniques
\[
\begin{tikzcd}[column sep=small]
& x_{\alpha} \times_x x_{\beta} \arrow[dl, "p_{\alpha,\beta}"'] \arrow[dr, "q_{\alpha,\beta}"] & \\
x_{\alpha} & & x_{\beta} \text{ .}
\end{tikzcd}
\]

(\textbf{N.B.} $B$ n'est par une catégorie ordonnée, car pour l'objet $(\alpha,\alpha)$, représenté par $x_{\alpha}\times_x x_{\alpha}$, % [sic] n'est par une -> n'est pas une, "par" for the negation "pas"
il y a \textit{deux} flèches dans l'objet $\alpha$, représenté par $x_{\alpha}$. Mais on a évidemment $\card \Fl(B) \leqslant \pi$.) Si on savait que $C$ est stable par produits binaires et par passage à des sous-objets (ici les $x_{\alpha} \times_x x_{\beta} \ \subset \ x_{\alpha} \times x_{\beta}$), on aurait prouvé la condition a). À défaut, utilisons encore le fait que $C$ est génératrice par épimorphismes (même stricts) et choisissons pour tout couple $(\alpha,\beta)$ un ensemble épimorphique de flèches :
\[ \underbrace{x_{\alpha,\beta,\gamma}}_{\in \Ob C_0} \xrightarrow{u_{\alpha,\beta,\gamma}} x_{\alpha} \times_x x_{\beta} \ , \]
% foot: Version 17 décembre 2025

\printedpage{41}
% head: A. Grothendieck : Dérivateurs, Ch. XVIII, Édition des Universités de Montpellier II et Paris VII
où $C_0$ est une sous-catégorie réduite exhaustive de $C$. La flèche $u_{\alpha,\beta,\gamma}$ est connue quand on connaît les composés
\msp{page 131}
\[
\begin{tikzcd}
& x_{\alpha} \\
x_{\alpha,\beta,\gamma} \arrow[ur, "v_{\alpha,\beta,\gamma}"] \arrow[dr, "w_{\alpha,\beta,\gamma}"'] & \\
& x_{\beta} \text{ ,}
\end{tikzcd}
\]
donc pour $x_{\alpha,\beta,\gamma} \in \Ob C_0$ fixé, les choix possibles pour $u_{\alpha,\beta,\gamma}$ sont de cardinal majoré par $\pi_1\pi_1 = \pi_1$, de plus on a $\card \Ob C_0 \leqslant \pi_1$, donc le nombre total des flèches $v_{\alpha,\beta,\gamma}$, $w_{\alpha,\beta,\gamma}$ (pour $\alpha$, $\beta$ fixés) est majoré par $\pi_1 \cdot \pi_1 = \pi_1$. Donc le nombre total des flèches $u_{\alpha,\beta,\gamma}$, pour $\alpha$, $\beta$, $\gamma$ variables, est majoré par $\pi \cdot \pi_1 = \pi$ (si $\pi \geqslant \pi_1$). Donc $x$ est représenté par une catégorie $I$ ayant comme objets l'ensemble somme de $A$ et de l'ensemble des indices triples $\alpha,\beta,\gamma$, et pour chaque tel indice, exactement deux flèches $v_{\alpha,\beta,\gamma}$, $w_{\alpha,\beta,\gamma}$ non identiques de source cet objet, et allant respectivement dans $x_{\alpha}$, $x_{\beta}$. On a donc $\card \Fl I \leqslant \pi$, et on a un foncteur $I \longrightarrow C$, $\alpha \longmapsto x_{\alpha}$, $(\alpha,\beta,\gamma) \longmapsto x_{\alpha,\beta,\gamma}$, tel que $x \ = \ \varinjlim\limits_{i\in I} x_i$. Cela achève la démonstration de 1°).

\msp{page 132}
\textit{Prouvons 2°.} La première inclusion (4.5.13.1) est tautologique, la deuxième ne fait qu'exprimer $P(\pi) \Longrightarrow \mathrm{d}$. Le fait que $\pi^{\pi_0} = \pi$ implique l'égalité des membres extrêmes est un rappel de 4.5.6. Reste à prouver que $\Filt'^{\pi}$ est stable dans $\mathcal{M}$ par sommes de cardinal $\leqslant \pi$. Mais si $(x_{\alpha})_{\alpha\in A}$ est une famille d'objets de $\Filt'^{\pi}$, avec $\card A \leqslant \pi$, et si $x = \coprod_{\alpha\in A} x_{\alpha}$, prenant pour tout $x_{\alpha}$ un foncteur $\varphi_{\alpha} : I_{\alpha} \longrightarrow C$, avec $\card \Fl I_{\alpha} \leqslant \pi$, tel que $x_{\alpha}$ soit la $\varinjlim$ dans $\mathcal{M}$ de ce foncteur, on trouve que $x$ est $\varinjlim$ du foncteur $\varphi : I \longrightarrow C$ défini par les $\varphi_{\alpha}$, où $I = \coprod_{\alpha} I_{\alpha}$ (somme dans Cat). Il est évident que $\card \Fl(I) \leqslant \pi\cdot\pi = \pi$, ce qui achève de prouver 2°.

\dgrblock{dgrb993}{Strict epimorphic presentations and the number of relations}
\begin{dcomment}{exposition}{p131}{Strict epimorphic presentations and the number of relations}
Grothendieck's counting argument on pages 130--131 gives a useful concrete form of the one-step presentation theorem. It should be read with the hypothesis $x_\alpha\in C$ in condition c); this is used in every hom-set estimate, although it is absent from the printed statement. Without it the identity of an arbitrary $X$ would satisfy c), making the proposed characterization vacuous.

Suppose $(c_\alpha\to X)_{\alpha\in A}$ is a strict epimorphic family, $|A|\leq\pi$, with all $c_\alpha\in C$. Because $\mathcal M$ has pullbacks, strictness means that $X$ is the colimit of the family together with all intersections $c_\alpha\times_Xc_\beta$ and their two projections. An intersection need not lie in $C$. To replace it, use all maps $c\to c_\alpha\times_Xc_\beta$ with $c\in C_0$. Such a map is determined by its two composites $c\to c_\alpha,c\to c_\beta$. For $\pi\geq\pi_1$, there are at most $\pi_1$ choices for each fixed pair $(\alpha,\beta)$, and at most $\pi^2\pi_1=\pi$ choices altogether.

Form a category with one vertex for each $c_\alpha$ and one vertex for each selected $c$, with two arrows from the latter vertex to its two targets. Its colimit is $X$. To verify this, a cocone gives maps $c_\alpha\to Z$ agreeing after every selected map into an intersection. The family of selected maps is jointly epimorphic on that intersection, so the two maps from the intersection to $Z$ agree. Strictness of the original family then gives one and only one map $X\to Z$. This proves condition a) directly and justifies the arrow count. The two projections for $\alpha=\beta$ must both be retained; the indexing category is not generally a poset.

Conversely, the legs of any colimit diagram form a strict epimorphic family: a family compatible with all equalities in the common target is in particular compatible with every arrow of the diagram, and therefore descends uniquely by the colimit property. Thus, for $\pi\geq\pi_1$, condition c), with its missing domain requirement restored, characterizes $D_\pi=\mathcal M_\pi$.

The presentation theorem also sharpens the counts left unfinished around (4.5.9.1). Throughout this count assume $\pi'\geq\pi\geq\pi_1$. When $\pi\geq\pi_1$, there are at most $2^\pi$ isomorphism classes of diagrams of at most $\pi$ arrows with values in $C_0$, and hence at most $2^\pi$ objects in a skeleton of $\mathcal M_\pi$. If $A\in\mathcal M_\pi$ and $X\in\mathcal M_{\pi'}$, then a one-step presentation of $A$ gives
\[
 |\Hom(A,X)|\leq
 \prod_{i\in\Ob I}|\Hom(c_i,X)|
 \leq\bigl((\pi')^{\pi_0}\bigr)^\pi
 =(\pi')^\pi.
\]
Taking $\pi'=\pi$ in the hom estimate also bounds the arrows in a skeleton of $\mathcal M_\pi$ by $2^\pi$. An arrow of the slice is determined by such an arrow $A\to B$ and a map $B\to X$. Consequently a skeleton of $\mathcal M_\pi/X$ has at most $(\pi')^\pi$ arrows. The sharper bound is therefore available even for the canonical slice, once the bounded presentation theorem has been established.
\end{dcomment}


\textit{Prouvons 3°.} On a les implications tautologiques $1 \Longrightarrow 1'$, et $1' \Longrightarrow 1$, puisque $\Filt'^{\pi}$ est stable par sommes de cardinal $\leqslant \pi$. Donc $1 \Longleftrightarrow 1'$, et implique que pour tout $x$ dans $\mathcal{M}$,
\msp{page 133}
$I \ = \ \Filt'^{\pi}/x$ est stable par $\varinjlim$ de cardinal $\leqslant \pi$, donc est grand devant $\pi$. Comme $x \simeq \varinjlim\limits_{i\in I} x_i$, cela prouve que $1 \Longrightarrow 2$. L'implication $2 \Longrightarrow 3$ est standard, et résulte de $\Filt'^{\pi} \subset \mathcal{M}_{\pi}$, je vais expliciter dans un

\textbf{Lemme 4.5.14.} \textit{Soient $\mathcal{M}$ une catégorie stable par petites} $\varinjlim$ \textit{, $\pi$ un cardinal infini, $\mathcal{M}_0$ une sous-catégorie strictement pleine de $\mathcal{M}$. Considérons les conditions}
\begin{itemize}
\item[\textit{a)}] \textit{$\mathcal{M}_0$ est une sous-catégorie strictement génératrice, et stable dans $\mathcal{M}$ par} $\varinjlim$ \textit{de cardinal $\leqslant \pi$.}
\item[\textit{b)}] \textit{Pour tout objet de $\mathcal{M}$, on a un isomorphisme}
\[ x \ \simeq \ \varinjlim_{i\in I} x_i \ , \]
\textit{les $x_i$ dans $\mathcal{M}_0$, $I$ grande devant $\pi$.}
% foot: Version 17 décembre 2025

\printedpage{42}
% head: A. Grothendieck : Dérivateurs, Ch. XVIII, Édition des Universités de Montpellier II et Paris VII
\item[\textit{b$'$)}] \textit{Comme b), avec $I$ catégorie} ordonnée.
\item[\textit{c)}] \textit{$\mathcal{M}_{\pi} \subset \mathcal{M}_0$ (où $\mathcal{M}_{\pi}$ est la sous-catégorie strictement pleine de $\mathcal{M}$ formée des objets $\pi$-accessibles).}
\end{itemize}
\textit{On a donc les implications}
\[ a \ \Longrightarrow \ b \ \Longleftrightarrow \ b' \ \Longrightarrow \ c \ , \]
\textit{et si on suppose $\mathcal{M}_0 \subset \mathcal{M}_{\pi}$, et $\mathcal{M}_{\pi}$ sous-catégorie strictement génératrice de $\mathcal{M}$, alors les quatre conditions $a$, $b$, $b'$, $c$ sont équivalentes, et équivalent à $\mathcal{M}_0 = \mathcal{M}_{\pi}$} (\footnote{cf. démonstration de proposition 4.5.6, $4 \Longrightarrow 5$, page 117.}).

\msp{page 134}
Enfin (revenant à la démonstration de 4.5.13, 3°), l'implication $3) \Longrightarrow 1)$ provient du fait que $\mathcal{M}_{\pi}$ est stable par $\varinjlim$ de cardinal $\leqslant \pi$.

Reste à prouver que ces conditions, i.e. la stabilité de $\Filt'^{\pi}$ par conoyaux de double-flèches, est satisfaite dans les cas A) B) C), le cas A) étant d'ailleurs déjà connu. Le cas B) résulte immédiatement de la transitivité des familles épimorphiques strictes universelles, et c'est d'ailleurs aussi contenu explicitement (stabilité de $\Filt'^{\pi}$ par $\varinjlim$ de cardinal $\leqslant \pi$) dans SGA 4 I 9.13 (sinon dans l'énoncé, du moins c'est signalé dans la démonstration). En rédigeant SGA 4 I 9, je savais donc qu'on a $\mathcal{M}_{\pi} \subset \Filt'^{\pi}$ pour tout $\pi$, mais l'inclusion, plus évidente pourtant dans le cas d'espèce, $\Filt'^{\pi} \subset \mathcal{M}_{\pi}$ (quand on suppose les objets de $C$ $\pi_1$-accessibles) m'avait apparemment échappé\,?

\dgrblock{dgrb994}{The retract missing from the abstract lemma}
\begin{dcomment}{development}{p134:Reste à prouver que ces conditions}{The retract missing from the abstract lemma}
Grothendieck isolates Lemma 4.5.14 to make the passage from presentations to presentability reusable. Its essential mechanism is correct: presentability factors the identity through one stage. What this produces is a retract, and a strictly full subcategory need not contain all retracts of its objects.

\textbf{A counterexample to b)$\Rightarrow$c) in Lemma 4.5.14.}\ours{}
Fix an infinite cardinal $\pi$, take $\mathcal M=\Ens$, and let $\mathcal M_0$ be the full subcategory of sets of cardinality at most $\pi$ except the singleton sets. Every object of $\mathcal M_0$ is $\pi$-accessible, and $\mathcal M_0$ is essentially small and strictly generating: its two-element sets already map jointly surjectively to every nonempty set. Also $\mathcal M_\pi$ is strictly generating, so the additional hypotheses in the last sentence of the printed lemma are satisfied.

For a set $X$ with at least two elements, use the poset of its subsets $A$ with $2\leq|A|\leq\pi$. This is $\pi^+$-directed, because a union of at most $\pi$ such subsets still has at most $\pi$ elements, and its colimit is $X$. The empty set has its constant presentation. For a singleton, index a diagram by the ordinal $\pi^+$, put the two-element set $\{0,1\}$ at every stage, and use the constant map with value $0$ for every strict transition. Identities remain identities, and composition is respected because the constant map is idempotent. The ordinal is $\pi^+$-directed, and the colimit of this diagram is a singleton. Thus b$'$), and hence b), hold for every $X$. Condition c) fails, since a singleton is $\pi$-accessible but is not in $\mathcal M_0$.

The corrected implication is
\[
 \mathrm{b}\quad\Longrightarrow\quad
 \mathcal M_\pi\subseteq
 \{\text{retracts of objects of }\mathcal M_0\}.
\]
Its proof is exactly the identity-factorization argument. Adding closure of $\mathcal M_0$ under retracts restores b)$\Rightarrow$c), and restores the final equivalence when the other hypotheses of the lemma hold. In the applications, essential smallness of the relevant slices has already been established; an abstract formulation should retain an appropriate small-presentation hypothesis when using $\mathcal M_0/X$ as an indexing category.

Under the small-presentation hypothesis just specified, the implication a)$\Rightarrow$c), which is the one used in Proposition 4.5.17, survives. Closure under colimits of size at most $\pi$ includes coequalizers, hence splitting of idempotents. Starting from the canonical dense presentation, the same closure makes $\mathcal M_0/X$ $\pi^+$-filtered; the resulting retract therefore belongs to $\mathcal M_0$. Likewise our proof that $D_\pi=\mathcal M_\pi$ established coequalizer closure before invoking identity factorization. The repair does not weaken the final canonical filtration theorem.
\end{dcomment}


Reste à traiter le cas C). $\Big($J'ai l'impression qu'il est un peu plus général que le cas B), i.e. que la condition B) sur $\mathcal{M}$ implique C), mais non inversement. Ainsi, la condition C) est vérifiée dès que $\mathcal{M}$ correspond à une espèce de structure algébrique définie en
\msp{page 135}
termes de $\varprojlim$ \textit{finies}, i.e. quand $\mathcal{M}$ est de la forme $\Ind(C)$, $C$ stable par $\varinjlim$ \textit{finies}, alors que je doute que la condition B) soit toujours satisfaite pour ces catégories. Par exemple dans la catégorie Cat (qui est un paratopos de cette espèce), les $\varinjlim$ filtrantes sont exactes à gauche, mais je ne crois pas que les familles épimorphiques stricts soient universelles. % [sic] épimorphiques stricts -> épimorphiques strictes, agreement with "familles"
Prenons p.ex. deux homomorphismes de groupes $G_0 \xrightarrow{u_0} G$, $G_1 \xrightarrow{u_1} G$ tels que $G_0 * G_1 \longrightarrow G$ soit surjectif, alors cette famille à deux éléments est épimorphique stricte dans Gr et dans Cat (\footnote{!, c'est \textit{faux} en général, mais vrai dans le cas particulier plus bas.}), mais n'est pas forcément épimorphique stricte universelle, même si p.ex. $G_0 = G_1 = H$, $G = H \times H$, et $u_0, u_1$ sont des injections $H \rightrightarrows H \times H$, données par $u_0(h) = (h,1)$, et $u_1(h) = (1,h)$. Si je prends l'application diagonale $H \longrightarrow H \times H$, les images inverses de $u_0$, $u_1$ sont le même objet de $\mathcal{M}/H$, savoir $e \longrightarrow H$, et cette famille n'est épimorphique stricte que si $H$ est le groupe unité.$\Big)$

Le cas C) va résulter du résultat plus précis :

\textbf{Corollaire 4.5.15.} \textit{Sous les conditions générales de 4.5.13, supposons que dans $\mathcal{M}$ les} $\varinjlim$
\msp{page 136}
filtrantes \textit{commutent au changement de base (p.ex. les} $\varinjlim$ \textit{filtrantes sont exactes). Alors pour tout $\pi \geqslant \pi_1$,} $\Filt'^{\pi}$ \textit{est formée des $x \in \Ob \mathcal{M}$ satisfaisant la condition suivante.}
% foot: Version 17 décembre 2025

\dgrblock{dgrb995}{A free product repairs the group example}
\begin{dcomment}{examples}{p135}{A free product repairs the group example}
Grothendieck wants to show that universality of filtered colimits is compatible with failure of universality for strict epimorphic families. The intended distinction is sound, but the two axes of $H\times H$ do not give the promised example. For nontrivial $H$, their common intersection is the trivial group. Maps from the two copies of $H$ to $H*H$ agree on that intersection, but do not extend to $H\times H$: the images of the two factors would have to commute. The normal form in a free product shows that nonidentity elements in the two factors do not commute. Thus the axes are not a strict epimorphic family, contrary to the assertion retained in footnote 53.

A nearby example proves exactly what is needed. Put
\[
 G=\langle a,b\mid a^2=b^2=1\rangle=C_2*C_2,
 \qquad A=\langle a\rangle,\quad B=\langle b\rangle.
\]
The inclusions $A\to G$ and $B\to G$ form a strict epimorphic family in groups. Any compatible pair of maps out of $A$ and $B$ extends uniquely to the free product; their intersection is trivial, so compatibility imposes only agreement on the unit. The element $ab$ has infinite order, since every nonzero power has a nonempty reduced word. Hence $K=\langle ab\rangle$ is infinite cyclic, and reduced-word normal form gives
\[
 K\cap A=1=K\cap B.
\]
Pulling the family back along $K\hookrightarrow G$ therefore gives two copies of $1\to K$. This family is not even jointly epimorphic: the trivial homomorphism $K\to C_2$ and the homomorphism sending $ab$ to the nonidentity element agree on both copies of $1$.

The same construction works in $\mathrm{Cat}$ by viewing groups as one-object categories. A compatible pair of functors $BA,B B\to\mathcal D$ must select the same object of $\mathcal D$, because their restrictions to the common terminal category agree. Their homomorphisms into that object's automorphism group extend uniquely to $A*B$. This proves strict epimorphy in $\mathrm{Cat}$ itself, with arbitrary target categories, rather than merely in its subcategory of groups. Pullback along $BK\to BG$ gives the same failure.

Finally, ordinary filtered colimits in $\mathrm{Cat}$ commute with finite limits. One direct verification uses the nerve. For every finite ordinal $[n]$, a functor $[n]\to\operatorname{colim}_i\mathcal A_i$ is specified by finitely many objects, arrows, and composition equalities, so it occurs at one stage of a filtered diagram; equality between two such functors is likewise witnessed at a later stage. Thus the nerve preserves filtered colimits. It also preserves limits and reflects isomorphisms. In simplicial sets the required commutation is pointwise the familiar set-theoretic fact, so it holds in categories. This settles the direction relevant to Grothendieck's example: condition C) does not imply condition B). The converse suggestion B)$\Rightarrow$C) is a separate assertion, not established by this example.
\end{dcomment}


\printedpage{43}
% head: A. Grothendieck : Dérivateurs, Ch. XVIII, Édition des Universités de Montpellier II et Paris VII
\begin{itemize}
\item[\textit{c$'$)}] \textit{Il existe une famille de flèches $(u_{\alpha})_{\alpha\in A}$, $u_{\alpha} : x_{\alpha} \longrightarrow x$, dans $\mathcal{M}$, qui est} épimorphique stricte universelle \textit{(i.e. couvrante pour la topologie canonique de $\mathcal{M}$), avec $x_{\alpha} \in \Ob C$ $\forall\ \alpha \in A$, et} $\card A \leqslant \pi$.
\end{itemize}
\textsc{Démonstration} : il faut prouver que cette condition équivaut à la condition c) de 4.5.13 1°). Comme on a $\mathrm{c}' \Longrightarrow \mathrm{c}$ trivialement, il reste à montrer que $\mathrm{c} \Longrightarrow \mathrm{c}'$. Mais on a vu dans 4.5.13 1°) que c implique a$''$, i.e. qu'on a un isomorphisme
\[ x \ = \ \varinjlim_{i\in I} x_i \ , \]
les $x_i$ dans $C$, et $I$ ensemble \textit{ordonné filtrant}, avec $\card I \leqslant \pi$. Il reste à voir que si $x$ dans $\mathcal{M}$ est une limite inductive \textit{filtrante} d'objets $x_i$, alors l'ensemble des flèches $x_i \longrightarrow x$ est épimorphique stricte \textit{universelle.} Or cela est immédiat, car pour $x' \longrightarrow x$, on a $x' = \varinjlim\limits_{i} x'_i$ (par hypothèse C), où $x'_i = x' \times_x x_i$, donc la famille $x'_i \longrightarrow x'$ est épimorphique stricte, q.e.d.

Prouvons maintenant que dans le cas C, la condition de 3°, 1 est satisfaite, i.e. que $\Filt'^{\pi}$ est stable par $\varinjlim$ de cardinal $\leqslant \pi$. Cela
\msp{page 137}
résulte aussitôt du

\textbf{Corollaire 4.5.16.} (de 4.5.15). \textit{Sous les conditions de 4.5.15 (i.e. la condition C) de 4.5.13, 3°), soit $x$ un objet de $\mathcal{M}$, et supposons qu'il existe une famille} épimorphique \underline{universelle} $(x_{\alpha} \longrightarrow x)_{\alpha\in A}$ \textit{de flèches de $\mathcal{M}$, avec $x_{\alpha} \in$} $\Ob \Filt'^{\pi}$ \textit{pour tout $\alpha \in A$, et} $\card A \leqslant \pi$. \textit{Alors $x \in$} $\Ob \Filt'^{\pi}$.

En effet, l'hypothèse sur $x_{\alpha}$ équivaut, en vertu de 4.5.15, à l'existence d'une famille épimorphique stricte universelle de flèches $(x_{\alpha\beta} \longrightarrow x_{\alpha})_{\beta\in B_{\alpha}}$, avec $\card B_{\alpha} \leqslant \pi$ et les $x_{\alpha\beta}$ dans $C$. Soit $B = \coprod_{\alpha\in A} B_{\alpha}$, on a alors $\card B \leqslant \pi^2 = \pi$, et d'autre part, la famille des flèches composées $x_{\alpha\beta} \longrightarrow x_{\alpha} \longrightarrow x$ indexée par $B$ est épimorphique stricte universelle (par transitivité de ce type de familles), donc $x \in \Ob \Filt'^{\pi}$, q.e.d.

Avec tout ça, on n'a toujours pas donné de caractérisation de $\mathcal{M}_{\pi}$ lui-même, qui est meilleure que $\Filt^{\pi}$ et $\Filt'^{\pi}$ du point de vue catégorique, car seuls les $\mathcal{M}_{\pi}$ définissent une filtration cardinale, et de plus ils ont les propriétés de stabilité désirables dans ce
\msp{page 138}
contexte, pour les $\varinjlim$ de cardinal $\leqslant \pi$. Le seul critère non tautologique que j'arrive à dégager dans le cas général (pour tout $\pi \geqslant \pi_0$, et sans hypothèse d'exactitude supplémentaire sur $\mathcal{M}$) est le suivant, beaucoup moins ``explicite'' ou ``calculatoire'' que les critères très concrets pour $\Filt^{\pi}$ et pour $\Filt'^{\pi}$ :

\dgrblock{dgrb996}{Universal epimorphisms can conceal too many relations}
\begin{dcomment}{development}{p137}{Universal epimorphisms can conceal too many relations}
Grothendieck's condition C) gives a useful covering interpretation of a small presentation. If $X=\operatorname{colim}_{i\in I}X_i$ is filtered and such colimits commute with pullback, then every pullback of its cocone is again a colimit cocone, hence is strictly epimorphic. Combined with the finite-diagram replacement above, this proves Corollary 4.5.15: for $\pi\geq\pi_1$, the one-step colimit class consists exactly of objects with a universal strict epimorphic family of at most $\pi$ maps from $C$.

Corollary 4.5.16 replaces ``strict epimorphic'' by ``epimorphic'' for the outer family, although its proof still uses strict descent. The weaker hypothesis does not suffice. The following example isolates the missing information: a surjection can cover every element while supplying none of the relations.

\textbf{A counterexample to Corollary 4.5.16 as printed.}\ours{}
Let $\mathcal R$ be the category whose objects are pairs $(S,R)$ with $R\subseteq S^{\mathbb N}$, and whose arrows are functions preserving this one countably-ary relation. We first verify the ambient hypotheses, so that the example does not rely on a category outside the scope of the proposition.

Limits have the usual underlying set-theoretic limits, with a tuple declared related precisely when every projection is related. A colimit has the underlying set-theoretic colimit, and its relation is the union of the images of the relations in the diagram. These descriptions follow directly by testing cones and cocones. All these colimits are universal. Indeed, for a colimit $(S,R)=\operatorname{colim}_i(S_i,R_i)$ and an arrow $(T,Q)\to(S,R)$, the underlying set of the pulled-back colimit is $T$, by universality of set colimits. If $q\in Q$, its image in $R$ comes from some $r_i\in R_i$, by the definition of the colimit relation. Pairing the coordinates of $q$ and $r_i$ gives a related tuple in $(T,Q)\times_{(S,R)}(S_i,R_i)$ mapping to $q$. Thus its colimit relation is all of $Q$, proving the assertion, including condition C).

Let $C$ be a skeleton of objects with countable underlying set and countable relation. It is closed under countable colimits by the preceding formula. Its objects are $\aleph_1$-presentable. To lift a map from such an object into an $\aleph_1$-filtered colimit, lift its countably many elements, choose witnesses for its countably many related tuples, and then witness the resulting countably many coordinate equalities at one common stage. This produces a relation-preserving lift; the same common-stage argument proves eventual equality of two lifts. Every object is an $\aleph_1$-filtered union of its countable partial substructures: take countable sets of elements and countable sets of related tuples, including the countably many coordinates of the chosen tuples. These choices form a countably directed poset. Therefore $\mathcal R\simeq\Ind'_{\aleph_0}(C)$, with $\pi_0=\aleph_0$. There are at most $2^{\aleph_0}$ arrows in a skeleton of $C$: code countable carriers by subsets of $\mathbb N$, countable relations by countable lists of sequences, and maps by functions between countable sets.

Epimorphisms in $\mathcal R$ are exactly the maps surjective on underlying sets. Surjectivity suffices because the forgetful functor is faithful. If a map misses an element, two functions to the two-element set equipped with its full relation can distinguish that element while agreeing on the image; both are morphisms. Such epimorphisms are universal, since pullbacks are set-theoretic on carriers.

Choose $\pi=\beth_\omega$ and a set $S$ of cardinality $\pi$. The arrow
\[
 q:(S,\varnothing)\longrightarrow(S,S^{\mathbb N}),
 \qquad q=\mathrm{id}_S,
\]
is universally epimorphic. Its source is $\pi^+$-presentable: its hom-functor is simply $\Hom_{\Ens}(S,U(-))$, where $U$ preserves colimits. Its target is not $\pi^+$-presentable. Indeed, $|S^{\mathbb N}|>\pi$. To see this directly, partition any purported list of at most $\pi$ sequences into an increasing union of sublists $L_n$ of cardinality less than $\pi$, and choose the $n$th coordinate outside the coordinates appearing there. The resulting sequence differs from every listed sequence.

Now write $(S,S^{\mathbb N})$ as the colimit of $(S,R_0)$ over all subsets $R_0\subseteq S^{\mathbb N}$ with $|R_0|\leq\pi$. This index is $\pi^+$-directed, and the transition maps are identities on carriers. If the target were presentable, its identity would factor through some $(S,R_0)$. The factor map would necessarily be the identity on $S$, which preserves the full relation only if $S^{\mathbb N}\subseteq R_0$, an impossibility. The source belongs to $D_\pi$ already as a coproduct of $\pi$ singleton structures with empty relation, whereas the target does not, since $D_\pi\subseteq\mathcal M_\pi$. The single arrow $q$ therefore violates precisely the conclusion of Corollary 4.5.16, with $\pi\geq\pi_1$ and all its other hypotheses satisfied.

The same category even satisfies B). A family is strictly epimorphic precisely when it is jointly surjective on carriers and every relation tuple of its target is the image of a relation tuple in one member. For necessity, strictness implies joint epimorphy, and the preceding two-point test proves joint surjectivity. Factor the family through the structure on the same carrier with only the image relation. The factor maps are compatible with the original family, so strictness supplies a map back to that structure. Joint surjectivity forces its underlying function to be the identity, hence forces every target tuple to lie in the image relation. Conversely, compatible maps out of the family agree on each underlying pairwise intersection, since the intersection structure has the set-theoretic pullback as carrier. They therefore descend to a unique function on the target carrier; the tuple condition makes this function a morphism. Both conditions are preserved by pullback, by the coordinate-pairing argument above. Thus the counterexample is entirely about the omission of strictness, not a failure of either exactness condition.

The positive statement is restored by requiring the outer family to be universally strict epimorphic. Choose the universal strict covers of its members furnished by Corollary 4.5.15. Their composites form a universal strict cover of the target by transitivity, with at most $\pi\cdot\pi=\pi$ terms, so 4.5.15 applies again. This is the covering-composition principle for the canonical topology \citep[Expos\'e II, 2.5--2.6]{SGA4I1972}. We found no published treatment of this particular counterexample to the manuscript assertion; the proof is supplied here without a claim of priority.
\end{dcomment}


\textbf{Proposition 4.5.17.} \textit{Soit $\pi$ un cardinal $\geqslant \pi_0$. Alors $\mathcal{M}_{\pi}$ est égal à la plus petite sous-catégorie strictement pleine de $\mathcal{M}$, contenant $C$, et stable par} $\varinjlim$ \textit{de cardinal $\leqslant \pi$.} (Il est important de noter qu'il \textit{ne suffit pas} de dire : stable par $\varinjlim\limits_{I}$ , avec \textit{$I$ grande devant} $\pi_0$, et $\card I \leqslant \pi$. On trouve alors $\Filt^{\pi}$, qui risque d'être une catégorie strictement plus petite que $\mathcal{M}_{\pi}$. Ce n'est pas non plus la catégorie des $\varinjlim\limits_{I} x_i$, avec $x_i \in C$ $\forall\ i \in \Ob I$, et $\card \Fl I \leqslant \pi$ -- puisque cette catégorie-là est $\Filt'^{\pi}$.)

\textsc{Démonstration}. Soit $\mathcal{M}'$ la catégorie de l'énoncé, on a donc $\mathcal{M}' \subset \mathcal{M}_{\pi}$, puisque $\mathcal{M}_{\pi}$ contient $C$ et est stable par les $\varinjlim$ de cardinal $\leqslant \pi$. Inversement, le lemme 4.5.14, $\mathrm{a} \Longrightarrow \mathrm{c}$.
% foot: Version 17 décembre 2025


\dgrblock{dgrb997}{The closure description and the later presentation theorem}
\begin{dcomment}{status}{p138}{The closure description and the later presentation theorem}
Grothendieck ends this portion with a characterization that is intrinsic to the category: start with its chosen small collection of $\pi_0^+$-presentable objects and close under colimits having at most $\pi$ arrows. Proposition 4.5.17 is correct. The generating collection lies in $\mathcal M_\pi$, and that class is closed under the prescribed colimits. Conversely, the canonical filtered presentation and the splitting of its resulting retract give the reverse inclusion. Only a)$\Rightarrow$c) of the corrected Lemma 4.5.14 is needed.

The additional warning that a single colimit of this size cannot suffice should not be retained as a present-day conclusion. Our explicit flattening proof shows that such one-step colimits already form $\mathcal M_\pi$. The general theorem is the assertion attributed to Ad\'amek--Rosick\'y, Remark 1.30, in \citet{Barton2019Presentable}; Barton's authored proof removes the retract using good diagrams and the fat small object argument. That proof also records why replacing arbitrary small colimits by colimits that are $\pi_0^+$-filtered is a stronger demand. The arithmetic obstruction exposed by sets remains real for the latter demand.

There are therefore two distinct outcomes to the exploratory discussion. The intrinsic filtration exists at every rank $\pi\geq\pi_1$, with the stated quantitative bound, and its one-step colimit description needs no extra exactness condition. The full canonical presheaf presentation can still have too many maps, and a small presentation need not be cofinal in it. The unrepaired implication in Lemma 4.5.14 and the group example do not invalidate these positive results; their corrections explain which arguments actually prove them.
\end{dcomment}
\begin{otherlanguage}{english}
\dgroverview{XVIII-a}{Chapter XVIII \textperiodcentered\ \textup{XVIII-a} \textperiodcentered\ Overview}
\begin{otherlanguage*}{english}
\nocite{Grothendieck1990XVIII}
The opening of Chapter XVIII asks what a category of coefficients must retain if the formal operations of sheaf theory are to survive outside a topos. Grothendieck begins with colimits and a small dense collection of test objects. The proposed tensor then forces him to ask whether continuity conditions define a reflective category, and that question leads to accessibility. The long investigation of cardinal filtrations is therefore part of the construction of coefficients. His letter to Thomason of 2 April 1991 explicitly describes this chapter as a return to the accessibility ideas of SGA~4, Expos\'e I, Section 9 \citep{Grothendieck1991Letter,SGA4I1972}.

These pages run from manuscript page 50 to page 138, corresponding to printed pages 1--43 of the edition. They contain Sections 4.1--4.5: pseudo-topoi and the tensor problem, its small-presheaf formulation, the proposed accessibility criterion, the restricted ind-completion, and the intrinsic cardinal filtration. The text often records an unsuccessful proof and then retains the objection which caused its abandonment. A faithful reading has to preserve this exploratory structure. An assertion subsequently withdrawn is not evidence that the intended theorem is false; conversely, the truth of the final theorem does not validate an intermediate equivalence.

We work in a fixed Grothendieck universe and use smallness relative to that universe. One translation governs the entire discussion. If $\pi$ is infinite, a category called \emph{grande devant $\pi$} is filtered with respect to diagrams having at most $\pi$ arrows. In the modern regular-cardinal convention it is $\pi^+$-filtered. Grothendieck's $\pi$-accessible objects are consequently $\pi^+$-presentable objects. His cocomplete categories with accessible small generators are locally presentable categories; the broader modern notion of an accessible category does not include cocompleteness. The cardinal attached to an object and one uniform cardinal for a chosen small generating family must also be distinguished.

A pseudo-topos has enough density for the dual special adjoint functor theorem. Colimit-preserving functors out of it have right adjoints, and a contravariant set-valued functor taking colimits to limits is representable. These facts make the opening construction plausible. They also expose a distinction in the motivating terminology. If morphisms of pseudo-topoi are described by their inverse-image functors, then the actual bicategorical product has underlying category $\mathcal M\times\mathcal N$: a cocontinuous functor out of that category is determined by its restrictions to $(m,0)$ and $(0,n)$. The proposed $\mathcal M\boxtimes\mathcal N$ represents functors cocontinuous in each variable separately. It is a tensor, with $\mathbf{Set}$ as unit, rather than this Cartesian product. For Grothendieck topoi the tensor nevertheless agrees with the product in the more restrictive category of topoi and geometric morphisms; Pitts's Theorem 2.3 proves the ordinary comparison \citep{Pitts1985}.

The tensor is first described as the category of bifunctors
\[
 \mathcal M^{\mathrm{op}}\times\mathcal N^{\mathrm{op}}
 \longrightarrow\mathbf{Set}
\]
which take colimits to limits separately. Density embeds it in a presheaf category on a small product of test categories. This proves local smallness and gives a useful formula for its continuous extensions, but it does not by itself provide the reflector. Nor does the naive external product of two representable presheaves usually satisfy the continuity condition: the empty colimit already forces both representing objects to be terminal. The example $\mathbf{Ab}\boxtimes\mathbf{Ab}\simeq\mathbf{Ab}$ makes the failure concrete. Its external tensor is the ordinary tensor product over $\mathbb Z$, and it cannot be recovered from the Cartesian product formula suggested in the manuscript.

Accessibility supplies the missing reflector. Choose a common regular cardinal and small full subcategories of presentable objects closed under colimits below that cardinal. Separate continuity can now be imposed by orthogonality to a \emph{set} of maps between small presheaves. The commentary constructs the reflection by alternately adjoining extensions and identifying extensions which agree on their domains. Presentability controls the length of the construction. The resulting category is again locally presentable, and reflection of representables supplies its elementary tensors. This gives a complete existence proof in the setting that the chapter is trying to isolate. Kelly's treatment of categories with prescribed colimits supplies the ordinary categorical background \citep[Section 6.5]{Kelly1982}. Lurie's tensor on presentable $\infty$-categories is a later comparison with the same universal pattern, not a substitute for the ordinary cardinal estimates here \citep[Section 4.8.1]{Lurie2017HA}.

The structural assertion behind Theorem 4.3.6 survives the failure of its first numerical proof. Starting from a small dense family of presentable objects, close that family under suitably small colimits. Every slice of the resulting category over an object is filtered at the required rank, and its canonical colimit is that object. The ind-object Hom formula proves full faithfulness. Thus the category is a restricted ind-completion. This argument does not require the provisional filtration by sizes of categories of elements to have the properties claimed for it on pages 72--76. The latter proof uses an unestablished directed-hull bound and imports a covering estimate from SGA~4 without its universality hypothesis. The manuscript's own corrective footnotes are justified.

Section 4.4 gives the complementary construction. When $C$ has colimits of diagrams with at most $\pi$ arrows, its $\pi^+$-ind-completion is cocomplete and is the free cocomplete category which preserves those existing colimits. The proof passes through cofinality, simultaneous presentations of diagrams, bounded coproducts, and the expression of an arbitrary indexing category as a directed union of small subcategories. We supply the omitted comma-adjunction argument and separate the existence of individual ind-presentations from closure under filtered colimits. The unrestricted claim that restriction on all functor categories is fully faithful is false; the needed restriction to functors preserving the density presentations is fully faithful and repairs the universal-property proof. Henry's recent theorem gives an exact modern comparison for these diagram presentations, with the hypotheses on the diagram shape and on $C$ stated explicitly \citep{Henry2025}.

The difference between ordinary and more highly filtered ind-objects is essential. Let $C$ consist of countable sets. Its $\aleph_1$-ind-completion is sets, but the ordinary ind-object obtained from finite subsets of $\mathbb Z$ has, at $c$, only the finite-image maps $c\to\mathbb Z$. Translation gives it a $\mathbb Z$-action which cannot occur at any nonempty equivariant representable stage. This answers negatively the specific cofinality question on manuscript page 94. It also explains why the restricted completion need not be closed under ordinary filtered colimits calculated in the presheaf category.

Section 4.5 measures an object in two ways. Write $\mathcal M_\pi$ for its intrinsic presentability layer, and $\Filt^\pi$ for the layer whose restricted nerve has at most $\pi$ maps from each test object, once the chosen small category itself is bounded by $\pi$. In sets, with tests of size at most $\pi_0$, these conditions read respectively
\[
 |X|\leq\pi,
 \qquad |X|^{\pi_0}\leq\pi.
\]
They need not coincide. The corrected comparison is
\[
 \Filt^\pi\subseteq\mathcal M_\pi
        \subseteq\Filt^{\pi^{\pi_0}}.
\]
The fixed-point condition $\pi^{\pi_0}=\pi$ restores equality. No continuum hypothesis is needed for the comparison itself. The counterexample sketched in the manuscript by a supposedly smallest directed family of subsets is not valid; the commentary replaces it with a diagonal argument at a singular cardinal of countable cofinality.

There is a second distinction within the proposed primed filtration. An object may be the colimit of a diagram of at most $\pi$ objects and arrows from $C$ without admitting a cofinal diagram of that size in $C/X$. For $C$ the countable sets and $|X|=\beth_\omega$, the coproduct of the singleton subsets gives the first presentation. A cofinal presentation would give at most $\beth_\omega$ countable subsets cofinal among all countable subsets of $X$, which the same diagonal argument excludes. Thus the two triples of conditions in Proposition 4.5.13 cannot be identified.

The positive outcome is stronger than the closing warning on page 138 suggests. The class of \emph{one-step} colimits of $C$-diagrams with at most $\pi$ arrows is already $\mathcal M_\pi$, for every $\pi\geq\pi_0$. The commentary proves closure under coequalizers by recording, for each generator in a presentation of the source, only the bounded portion of a target presentation through which its two images factor. This keeps the enlarged diagram within the original cardinal bound. Coproducts and splitting of idempotents finish the proof. The general presentation theorem is known; Barton's authored proof discusses its earlier attribution and the difficulty with the originally cited justification \citep{Barton2019Presentable}. Our argument is included as a direct repair in Grothendieck's setting, without a claim of priority. It also yields the sharper bounds on the number of objects and arrows which the chapter seeks.

Two later corrections prevent the descent arguments from obscuring that conclusion. Presentability factors an identity through a stage as a \emph{retract}; Lemma 4.5.14 must retain closure under retracts before replacing this by membership in the chosen subcategory. And a universal epimorphism need not transmit a bound on the number of relations. A category of sets with one countably-ary relation gives a counterexample to Corollary 4.5.16 even when all colimits are universal. Requiring the outer family to be universally \emph{strict} epimorphic restores the corollary. Both corrections leave the intrinsic filtration theorem intact.

Footnote 50 suggests testing universality of filtered colimits in the generating category. We prove a precise version: if $C$ also has pullbacks, then ordinary filtered colimits in $\Ind'_{\pi_0}(C)$ are universal exactly when the filtered colimits in $C$ with at most $\pi_0$ arrows are universal. The proof approximates the entire cocone and its base-change arrow together, and then expresses any filtered indexing category as a $\pi_0^+$-directed union of filtered subcategories with at most $\pi_0$ arrows. The additional pullback hypothesis makes the test intrinsic to $C$.

The remaining verification limits concern the disputed implication in Proposition 4.3.1 for arbitrary pseudo-topoi and the suggestion that universality of strict epimorphic families implies universality of ordinary filtered colimits. We do not describe these as literature-established open problems. The relevant older discussions of density, ind-objects, and the distinction between formal and realized colimits are cross-referenced where they are used.
\end{otherlanguage*}
\end{otherlanguage}
\dgrchunkend
